\documentclass[10pt]{article}
% \documentclass[10pt, draft]{article} % draft mode (no figures)

\usepackage[margin=1in]{geometry}
\usepackage{amsmath,amssymb,amsthm,bm,mathtools}
\usepackage{algorithm}
\usepackage{algpseudocode}
\usepackage{enumitem}
\usepackage{graphicx}
\usepackage{subcaption}
\usepackage{booktabs}
\usepackage{tabularx}
\usepackage{makecell}
\usepackage{float}
\usepackage{xcolor}
\usepackage{hyperref}
\usepackage[style=numeric,sorting=nyt]{biblatex}
\usepackage{dsfont}

\usepackage{color,soul}
\usepackage[textsize=tiny]{todonotes}
\usepackage[normalem]{ulem} % provides \sout
\usepackage[most]{tcolorbox} % compile-safe crossed blocks for structural table moves

\definecolor{impGreen}{RGB}{0,128,0}
\definecolor{latestPurple}{RGB}{128,0,128}




\addbibresource{references.bib}
\addbibresource{references_additions.bib}

\hypersetup{hidelinks}

\newtheorem{theorem}{Theorem}
\newtheorem{proposition}{Proposition}
\newtheorem{lemma}{Lemma}
\newtheorem{corollary}{Corollary}
\theoremstyle{definition}
\newtheorem{definition}{Definition}
\newtheorem{assumption}{Assumption}
\newtheorem{remark}{Remark}

\newcommand{\R}{\mathbb{R}}
\newcommand{\E}{\mathbb{E}}
\newcommand{\Cov}{\operatorname{Cov}}
\newcommand{\Corr}{\operatorname{Corr}}
\newcommand{\Var}{\operatorname{Var}}
\newcommand{\diag}{\operatorname{diag}}
\newcommand{\1}{\mathbf{1}}
\newcommand{\D}{\mathcal{D}}
\newcommand{\Hcal}{\mathcal{H}}
\newcommand{\Vcal}{\mathcal{V}}
\newcommand{\argminop}{\operatorname*{arg\,min}}
\newcommand{\argmaxop}{\operatorname*{arg\,max}}
\newcommand{\advisorchange}[1]{{\color{blue}#1}}
\newcommand{\blueheading}[1]{\texorpdfstring{\textcolor{blue}{#1}}{#1}}
\newenvironment{advisorchanges}{\begingroup\color{blue}}{\endgroup}
\newcommand{\oldtext}[1]{\textcolor{blue}{\sout{#1}}}
\newcommand{\newtext}[1]{\textcolor{impGreen}{#1}}
\newcommand{\latesttext}[1]{\textcolor{latestPurple}{#1}}

% Structured deletions in this revision are shown in blue with an X-shaped strike.
% Inline deletions continue to use \oldtext. The block form is used only where
% \sout cannot safely span a table, displayed mathematics, or multi-paragraph material.
\newenvironment{oldrevisionblock}{%
  \begingroup
  \begin{tcolorbox}[
    enhanced,breakable,
    colback=white,colframe=white,boxrule=0pt,
    left=0pt,right=0pt,top=1pt,bottom=1pt,
    coltext=blue,
    overlay unbroken={\draw[blue,line width=.45pt] (frame.south west)--(frame.north east);\draw[blue,line width=.45pt] (frame.north west)--(frame.south east);},
    overlay first={\draw[blue,line width=.45pt] (frame.south west)--(frame.north east);},
    overlay middle={\draw[blue,line width=.45pt] (frame.south west)--(frame.north east);},
    overlay last={\draw[blue,line width=.45pt] (frame.north west)--(frame.south east);}
  ]
}{\end{tcolorbox}\endgroup}


% PURPLE PIPELINE NOTE (not manuscript text):
% The existing scripts/populate_paper_v6_994.py in the testing branch is hard-coded
% to paper/paper_v6.tex and paper/img/generated_v6_preselection. Reuse its mature
% 994-tree extraction/formatting logic, but do not run it unmodified against an
% older manuscript. For the local current paper, either (i) make PAPER_TEX/PAPER_IMG
% CLI-configurable and point them to paper_v12.tex, or (ii) create a thin
% populate_paper_v12_994.py wrapper that imports the same helpers. Guard the source
% manuscript hash and fill only the explicit placeholders below.

% Allows the source to compile outside the project directory.
\newcommand{\safeincludegraphics}[2][]{%
  \IfFileExists{#2}{\includegraphics[#1]{#2}}{%
  \fbox{\parbox[c][0.19\textheight][c]{0.92\linewidth}{\centering
  Figure file not found in this folder.\\[2mm]{\scriptsize\ttfamily\detokenize{#2}}}}}}

% \title{Reducing Regressivity in Automated Valuation Models for Fair Property Taxation:\\
% A Covariance-Guided Correction for Assessor Workflows}

\title{Fair Property Taxation:
A Covariance-Guided Correction for Assessor Workflows}

\author{
Nicol\'as Acevedo\thanks{Operations Research Center, MIT, Cambridge, MA 02139. \texttt{nacevedo@mit.edu}}
\qquad
Haihao Lu\thanks{Sloan School of Management, MIT, Cambridge, MA 02139. \texttt{haihao@mit.edu}}
\\[1ex]
Michael Wagner\thanks{Cook County Assessor's Office, Chicago, IL 60602. \texttt{Michael.Wagner2@cookcountyil.gov}}
\qquad
Timothy Sparer\thanks{Cook County Assessor's Office, Chicago, IL 60602. \texttt{Timothy.Sparer@cookcountyil.gov}}
}
\date{}

\begin{document}
\maketitle


\begin{abstract}
% \textcolor{blue}{\sout{Property value assessments determine how local property-tax burdens are distributed. Assessor offices increasingly use machine learning-based automated valuation models to estimate values for large and varied property inventories, but strong average prediction does not prevent \emph{regressivity}: lower-value properties can receive systematically higher assessment-to-sale ratios than higher-value properties. Working with data scientists and assessors at the Cook County Assessor's Office (CCAO), we develop a correction that can be added to its existing LightGBM residential workflow without replacing its data, features, or review process. The method adds a training penalty for a systematic relation between assessment errors and sale prices. We consider a direct squared-covariance penalty and a sample-additive, covariance-guided objective that is compatible with standard boosting software. Applied to a Python translation of the CCAO workflow, the validation-selected sample-additive model substantially improves three standard measures of vertical equity: their departures from neutral values fall by about 54\% to 88\% relative to ordinary LightGBM. Predictive fit changes little, while mean absolute error rises by about 3.5\%. The current comparison measures the full implemented custom-objective path and does not yet isolate the penalty from every implementation difference. The correction targets a global first-order pattern. It does not set the \oldtext{overall assessment level}\newtext{overall valuation level}, remove local inequity, or guarantee that every ratio-study measure will fall within its reference range. It should therefore be used with temporal validation, uncertainty analysis, subgroup ratio studies, and comparison with simpler post-hoc recalibration.}}

% \textcolor{impGreen}{
Property value assessments determine how local property-tax burdens are distributed among owners.
Assessor offices increasingly use machine learning to
estimate values for large and heterogeneous property inventories, given their strong predictive power. However, predictive accuracy alone does not secure fairness in the estimated values. A common problem in   


\emph{regressivity}: valuation-to-sale ratios can be systematically higher for lower-valued properties and lower for higher-valued properties.
Working with data scientists and assessors at the Cook County Assessor's Office (CCAO), we study whether this application-specific pattern can be reduced through the model-training objective without replacing the surrounding data, features, validation, or review workflow. We use the first-order association between log valuation errors and log sale prices as a training regularization mechanism. We consider a direct squared-covariance penalty and develop a sample-additive upper-bound surrogate that is compatible with standard boosting implementations. 
% \oldtext{The correction is intentionally targeted: it controls a global first-order sale-price-related pattern but does not determine overall valuation level, remove nonlinear or local inequity, or guarantee that assessor-facing vertical-equity measures fall within their reference ranges.}
\newtext{The correction is intentionally targeted at a global first-order sale-price-related pattern: it does not by itself determine overall valuation level, directly target nonlinear or local inequity, or guarantee that assessor-facing vertical-equity measures fall within their reference ranges.} 
% \oldtext{We evaluate the correction using chronological validation, held-out and forward evaluation, complementary ratio-study and dependence diagnostics, and value-group ratio profiles.}
\newtext{We evaluate the correction using chronological validation, held-out and forward evaluation, assessor-facing ratio-study measures, mechanism and residual-structure diagnostics, and value-group ratio profiles.} 
% \oldtext{Across the complete Direct and Surrogate paths, stronger regularization moves the targeted first-order and several assessor-facing vertical-equity diagnostics toward neutrality, while the ratio-shape and distance-correlation results show that global neutrality need not eliminate nonlinear or local value-related structure.} 
% \oldtext{Across the complete Direct and Surrogate paths, stronger regularization moves the targeted first-order and several assessor-facing vertical-equity diagnostics toward neutrality. The ratio-shape results show that near-neutral first-order or grouped diagnostics need not imply a flat value-related profile, while distance correlation shows that broader residual--price dependence can remain.}
\newtext{Across substantial portions of the Direct and Surrogate paths, increasing regularization attenuates the targeted first-order pattern and moves several assessor-facing vertical-equity diagnostics toward neutrality, while stronger regularization exposes tradeoffs in predictive performance and ratio shape. Near-neutral first-order or grouped diagnostics do not imply a flat value-related profile, and broader residual--price dependence can remain.}
% \oldtext{Exploratory multi-jurisdiction analyses separately examine transferability and comparison with simple post-hoc recalibration. No penalty strength is chosen.}
% \oldtext{Centered post-hoc recalibration provides a low-cost benchmark and can reproduce a large first-order correction, underscoring that the incremental value of in-processing lies in the attainable accuracy--equity path rather than in slope correction alone. Exploratory multi-jurisdiction analyses separately examine transferability and heterogeneity. We report complete regularization paths and do not select a deployment penalty strength.}
\newtext{A fixed-prediction-space analysis clarifies that the two objectives are related but not equivalent: the Direct penalty follows a rank-one covariance-correction path, whereas the Surrogate is a log-price-distance-weighted projection that can combine multiple learnable correction directions. Exploratory multi-jurisdiction analyses separately examine transferability and heterogeneity. We report complete regularization paths and do not select a deployment penalty strength.}

\end{abstract}



\section{Introduction}
\label{sec:introduction}

Property taxes are one of the largest sources of local government revenue in the United States. They help finance public schools, roads and streets, and other essential local services \cite{IAAO2020PropertyTaxPolicy}. Using 2016 U.S. Census Bureau data, \textcite{Berry2021} reports that property taxes accounted for roughly 72\% of local tax revenue and 47\% of local own-source general revenue, and totaled about \$500 billion annually.

% \textcolor{blue}{\st{The U.S. property tax system is generally value-based}}\textcolor{impGreen}
{Property taxation in the United States is generally based on property values}: local assessors estimate property values, which are then used, together with jurisdiction-specific assessment and taxation rules, to determine each taxpayer's share of the property-tax burden \cite{IAAO2020PropertyTaxPolicy,IAAO2025Mass}. Consequently, systematic overassessment of some properties and underassessment of others can shift tax liability from one group of property owners to another, creating or exacerbating inequities \cite{Berry2021,IAAO2020PropertyTaxPolicy}.

The {valuation} of groups of properties for tax purposes is known as \emph{mass appraisal}. The International Association of Assessing Officers (IAAO) defines mass appraisal as ``the process of valuing a group of properties as of a given date using common data, standardized methods, and statistical testing'' \cite{IAAO2025Mass}. Within this process, assessors use statistical models to estimate property values across large property groups. {When implemented in computer-based valuation systems, such models are commonly referred to as automated valuation models (AVMs)} \cite{IAAOAITaskForce2022,BidansetRakow2022}. Recent mass-appraisal research increasingly examines nonlinear machine-learning (ML) methods, including gradient-boosted trees, alongside traditional regression approaches \cite{Rootetal2023,Sevgen2024100111,Zilli2024}. %\textcolor{blue}{\st{Given their high performance in predictive power.}}

However, when evaluating the performance of a mass-appraisal valuation model, predictive accuracy alone is insufficient. The IAAO recommends the use of \emph{ratio studies}, which compare {appraised values} with observed sale prices. 
% \textcolor{blue}
% {\st{A central quantity in these studies is the assessment-to-sale ratio, defined as the assessed value divided by the sale price.}} 
% \textcolor{blue}{\sout{A central quantity in these studies is the ratio of appraised value to sale price. Because the CCAO model studied here predicts fair market value before the statutory conversion to assessed value, we evaluate its predictions on a common market-value scale, using the model-estimated market value divided by the observed sale price; we refer to this quantity as the assessment ratio throughout the paper. Ratio studies use such ratios to evaluate assessment level and uniformity, including whether similar properties are treated consistently (horizontal equity) and whether assessment ratios vary systematically across property-value levels (vertical equity).}}
% \oldtext{A central quantity in these studies is the ratio of a valuation estimate to a market-value proxy such as sale price. Because the CCAO model studied here predicts \emph{fair market value} before the statutory conversion to assessed value, we evaluate its model-estimated market value relative to the observed sale price and refer to this quantity as the \emph{valuation-to-sale ratio}, or \emph{valuation ratio} for brevity \cite{CCAOModelResAVM2026}.}
\newtext{A central quantity in these studies is the ratio of a valuation estimate to a market-value proxy such as sale price. Because the CCAO model studied here predicts \emph{fair market value} before the statutory conversion to assessed value and the subsequent steps that determine taxable value and tax liability, we evaluate its model-estimated market value relative to the observed sale price and refer to this quantity as the \emph{valuation-to-sale ratio}, or \emph{valuation ratio} for brevity \cite{CCAOModelResAVM2026,CCAOAssessmentNotice2026,CCAOPropertyTaxSystem2026}.} Ratio studies use such ratios to evaluate valuation level and uniformity, including whether similar properties are valued consistently (\emph{horizontal uniformity}) and whether valuation ratios vary systematically across property-value levels (\emph{vertical equity}) \cite{IAAO2013Ratio,IAAO2025Mass,IAAO2026ExposureRatio}.

% \textcolor{blue}{\sout{A particularly relevant form of vertical inequity is assessment regressivity: lower-value properties tend to receive higher assessment ratios than higher-value properties. This pattern can place a disproportionate share of the assessment, and potentially tax, burden on owners of lower-valued properties. Moreover, regressivity may remain hidden in aggregate measures of predictive accuracy because these measures do not capture whether assessment errors vary systematically with property value. A valuation model can therefore achieve strong predictive accuracy while still exhibiting systematic price-related bias and, consequently, inequitable assessment outcomes.}}

{A particularly relevant form of vertical inequity is assessment \emph{regressivity}: valuation ratios tend to be higher for properties with lower underlying market values and lower for properties with higher market values. Because market value is latent, empirical vertical-equity diagnostics necessarily rely on observed sales or other market-value proxies \cite{IAAO2026ExposureRatio}. This pattern can place a disproportionate share of the assessment, and potentially tax, burden on owners of lower-valued properties \cite{McMillen2020,Berry2021,AvenancioLeonHoward2022}. Moreover, regressivity may remain hidden in aggregate measures of predictive accuracy because these measures do not identify whether valuation errors vary systematically across the value distribution. A valuation model can therefore achieve strong predictive accuracy while still exhibiting a systematic value-related ratio pattern.}

% IS THIS PARAGRAPH WORTH KEEPING???
% \textcolor{blue}{\st{For property valuation modeling, this means that regressivity cannot be treated only as an ex-post diagnostic. If the goal is to preserve predictive performance while reducing systematic price-related bias, vertical equity must also enter model training or calibration as an explicit criterion. Implementing this objective is especially challenging in the mass-appraisal context: current sale prices are observed only when properties transact, so values for all properties in the jurisdiction must be inferred from a smaller and selected set of verified sales. Properties differ in size, age, condition, construction, use, and amenities; location effects operate at multiple geographic levels and vary over time; and administrative records may be missing, outdated, or inconsistent. At the same time, assessment systems must produce values within fixed assessment schedules while remaining auditable and defensible in taxpayer appeals. These constraints limit the types of equity corrections that are useful in practice: a method should reduce regressivity while preserving predictive performance, require limited changes to existing valuation workflows, and remain feasible to implement, audit, and defend within assessor timelines. Moreover, because valuation practices, available data, and administrative requirements differ across jurisdictions, the specific implementation of such a correction should also be adaptable to the local assessment context.}}
% \textcolor{blue}{\cite{IAAO2020SalesVerification,IAAO2018AVM,IAAO2025Mass,Almy2014}}

% \textcolor{blue}{\sout{This motivates treating regressivity as an explicit criterion in model training or calibration rather than only as an ex-post diagnostic. In mass appraisal, a useful correction must operate within existing sales-based modeling, validation, and review processes: current prices are observed only for properties that transact, while valuation systems must produce values within fixed assessment schedules and remain compatible with professional and administrative review. Accordingly, the correction should reduce regressivity while preserving predictive performance and requiring minimal changes to existing valuation workflows.}}

{This motivates asking whether regressivity that remains after fitting a strong predictive model can be reduced through an explicit training or calibration criterion rather than treated only as an ex-post diagnostic. In mass appraisal, a useful correction must operate within existing sales-based modeling, validation, and review processes: current prices are observed only for properties that transact, while valuation systems must produce values within fixed assessment schedules and remain compatible with professional and administrative review \cite{IAAO2020SalesVerification,IAAO2018AVM,IAAO2025Mass}. Accordingly, the correction should reduce regressivity while preserving predictive performance and requiring minimal changes to existing valuation workflows.}

This paper arises from a collaboration between the Cook County Assessor's Office (CCAO) and academic researchers, motivated by an operational goal: reducing regressivity in the CCAO's residential valuation workflow. The CCAO currently uses a county-wide LightGBM (Light Gradient-Boosting Machine) model \cite{lightgbm} as its primary residential valuation model, {following earlier township-level linear-regression and gradient-boosting workflows} \cite{CCAOModelResAVM2026}. The current model is trained on screened residential sales and estimates a \emph{fair market value} using property characteristics, location, environmental variables, and market trends. It operates within a broader assessment pipeline that separates data ingestion, model training, assessment, evaluation, interpretation, and finalization \cite{CCAOModelResAVM2026}.

% \textcolor{blue}{\st{Further, the value produced by the valuation model is only one input into the final property-tax bill. In Cook County, the Assessor's estimate of fair market value is first converted to an assessed value according to property classification, for example residential or commercial. A state equalization factor is then applied, and eligible exemptions can further reduce the resulting equalized assessed value. Local tax rates are ultimately applied to this taxable equalized assessed value to determine the tax bill. These rates are calculated by the Cook County Clerk based on the tax levies submitted by the relevant taxing districts and the taxable property values within those districts. In this paper, we focus solely on the valuation stage of the property-tax system, rather than the subsequent steps that convert estimated market values into final tax bills.}}
% \textcolor{blue}{\cite{CCAOAssessmentNotice2026,CCAOPropertyTaxSystem2026}}

% \textcolor{impGreen}
{The model's fair-market-value estimate is only one input into the final property-tax bill. In Cook County, that value is converted to assessed value according to property classification, equalized using the state equalization factor, reduced by eligible exemptions, and ultimately combined with local tax rates determined from taxing-district levies and taxable property values \cite{CCAOAssessmentNotice2026,CCAOPropertyTaxSystem2026}. We therefore focus on the valuation stage rather than the downstream steps that convert estimated market values into final tax liabilities.}

% \textcolor{blue}{\sout{Within the valuation stage, the central problem addressed in this paper is the interaction between predictive performance and regressivity. In the CCAO workflow considered, the gradient boosting model substantially improves predictive accuracy and ratio dispersion relative to the linear-regression baseline, but its vertical-equity diagnostics move in a more regressive direction. This creates a model-design problem: rather than choose between the linear and gradient-boosting baselines, we seek to retain the predictive advantage of gradient boosting while explicitly reducing regressivity. From an operational perspective, the correction should also fit within CCAO's existing modeling and validation workflow without requiring a fundamental change to CCAO's existing modeling and validation workflow.}}

{Within the valuation stage, the central problem addressed in this paper is the interaction between predictive performance and regressivity. In our controlled comparison within the translated CCAO workflow, unpenalized LightGBM substantially improves predictive accuracy and lowers aggregate ratio dispersion relative to the linear-regression baseline, but its vertical-equity diagnostics move in a more regressive direction. This creates a model-design problem: rather than choose between the two baselines, we seek to retain the predictive advantages of LightGBM while reducing the regressive value-related pattern. From an operational perspective, the correction should fit within CCAO's existing modeling and validation workflow without requiring a fundamental redesign of the surrounding valuation pipeline \cite{CCAOModelResAVM2026}.}

% \textcolor{blue}{\st{Within the valuation stage, the central problem addressed in this paper is the interaction between predictive performance and vertical equity. In the CCAO workflow considered, the gradient boosting model substantially improves predictive accuracy and ratio dispersion relative to linear regression-based models, but its vertical equity diagnostics move in a more regressive direction. This tradeoff between a simpler linear model and a more complex nonlinear model moves the discussion from a model-selection problem to a model-design problem: the goal is to preserve the predictive advantage of the nonlinear model while explicitly controlling its regressive behavior. From an operational perspective, we aim to achieve this goal without requiring a fundamental change to CCAO's existing modeling and validation workflow.}}
% \textcolor{blue}{\cite{CCAOModelResAVM2026}}

% \textcolor{blue}{\sout{We address this challenge by modifying the training objective of the valuation model so that regressivity becomes an explicit model-design consideration. Specifically, we formulate a direct squared-covariance penalty that controls global price-related dependence between valuation errors and property values, and then develop a sample-additive surrogate motivated by the same dependence-control objective that is compatible with LightGBM's custom-objective interface. For a fixed penalty formulation, the correction changes only the training objective and introduces one additional tuning parameter: the penalty strength. The correction requires no changes to the underlying data, feature set, or surrounding validation and evaluation workflow. The penalty strength can be selected through chronological validation using the predictive and ratio-study measures already used to evaluate mass-appraisal models. The resulting approach is therefore designed to reduce regressivity while remaining compatible with the existing machine-learning valuation pipeline. The framework is intentionally targeted: it does not, by itself, determine \oldtext{overall assessment level}\newtext{overall valuation level} or guarantee local or subgroup equity.}}

{We address this challenge by modifying the valuation model's training objective. Specifically, we use the global first-order association between log valuation errors and log sale prices as a trainable mechanism: we formulate a direct squared-covariance penalty and develop a sample-additive upper-bound surrogate that is compatible with LightGBM's custom-objective interface. For a fixed penalty formulation, the correction changes only the training objective and introduces one additional tuning parameter, the penalty strength. It requires no changes to the underlying data, feature set, or surrounding validation and evaluation workflow. \oldtext{For a later deployment decision, penalty strength can be calibrated through chronological validation only after an explicit operating criterion is specified; the present paper does not infer a universal or temporally invariant penalty range from the regularization path itself.}
\latesttext{For a later deployment decision, chronological validation can first be used to derive a family-specific \emph{candidate region} that screens the positive-$\rho$ path before any final operating-point rule is applied. The lower endpoint is only a soft activity-onset marker, while the upper endpoint is a family-specific guardrail against the first robust high-$\rho$ deterioration or nonlinear-structure warning. This screen narrows the configurations passed to downstream selection; it does not select $\rho$, define a universally ``safe'' interval, or imply that raw-$\rho$ endpoints are temporally invariant.} The framework is intentionally targeted: it does not, by itself, determine overall valuation level, eliminate nonlinear or local inequity, or replace the assessor-facing vertical-equity diagnostics used for evaluation.}

% \textcolor{blue}{\sout{Applied to CCAO's residential property-valuation workflow, the proposed methodology substantially reduces measured regressivity while preserving nearly all of the predictive performance of the unpenalized LightGBM. We use the oldest 90\% of eligible 2016--2024 sales (344,607 sales) for model development, evaluate the selected models on the newest 10\% (38,290 sales) as a held-out test set, and use 2025 (26,641 sales) as a secondary forward evaluation aligned with the assessment-year workflow after refitting on all 382,897 eligible 2016--2024 sales. Relative to unpenalized LightGBM, the validation-selected model reduces the deviations of three complementary vertical-equity diagnostics from their neutral values by approximately 54\% to 88\%. At the same time, $R^2$ decreases only from 0.885 to 0.881, while mean absolute error increases by approximately 3.5\%. Thus, the proposed methodology produces a substantial reduction in measured regressivity at a small cost in predictive accuracy. Exploratory experiments across six counties further suggest that the correction can improve value-related assessment patterns in different markets, although the appropriate penalty strength and attainable improvement vary across jurisdictions.}}

{In the primary CCAO application, model development and the complete penalty paths are evaluated chronologically within the development sample, followed by a primary held-out evaluation and a separate 2025 forward evaluation after refitting. The controlled baseline comparison reported in Section~\ref{subsec:baseline_tension} establishes the application-specific motivation: LightGBM improves prediction and ratio dispersion relative to the linear benchmark while exhibiting a stronger regressive vertical-equity pattern. 
% \oldtext{This draft reports the complete Direct and Surrogate regularization paths}
\newtext{We report the complete Direct and Surrogate regularization paths} without choosing a penalty strength. Exploratory experiments across six counties are used separately to examine transferability, heterogeneity, and failure modes rather than to make confirmatory claims about those jurisdictions.}

% \textcolor{blue}{\sout{The method was developed in collaboration with CCAO data scientists and has been incorporated into the office's residential model-development codebase. The CCAO has also used and tested the correction during model development. This implementation provides direct evidence that the proposed approach can be integrated into an existing public-sector mass-appraisal workflow without requiring a fundamental redesign of the underlying valuation pipeline.}}

{The method was developed and tested in collaboration with CCAO data scientists and is designed to integrate with the office's model-development workflow without requiring a fundamental redesign of the surrounding valuation pipeline.}

% \todo{Retain the existing collaboration-and-testing statement unless a public institutional record explicitly supports a stronger implementation/adoption claim.}


% \textcolor{blue}{\st{Accordingly, the method should be understood as targeting vertical equity in the valuation stage rather than equity in final property-tax liabilities.}}

% \textcolor{blue}{\sout{The paper makes three contributions. First, motivated by an important operational problem recognized by CCAO, we connect assessor-defined assessment regressivity with dependence-based constrained learning by translating the value-related error pattern into a tractable model-training objective. Second, we develop a direct squared-covariance penalty and a sample-additive surrogate that can be implemented through the standard LightGBM interface and evaluated within an assessor's existing temporal-validation and ratio-study workflow. Third, we demonstrate that the correction can substantially reduce regressivity while largely preserving predictive accuracy, and we document its incorporation into CCAO's residential model-development codebase.}}

% \oldtext{Our contribution is application-driven. Motivated by the regressive pattern in the controlled CCAO baseline comparison, we study whether that pattern can be reduced through the model-training objective without redesigning the surrounding valuation architecture. We use the natural first-order association between log valuation errors and log sale prices as a training mechanism, characterize the cross-observation coupling induced by the direct squared-covariance penalty, and develop a sample-additive upper-bound surrogate that can be implemented through standard gradient-boosting interfaces. We then evaluate whether these corrections improve assessor-facing vertical-equity diagnostics under chronological out-of-sample evaluation while preserving predictive performance. Supplementary multi-jurisdiction experiments examine transferability and the limits of the first-order correction.}

% \oldtext{Our contribution is application-driven and has three parts. First, we translate an assessor-facing regressivity pattern into a trainable first-order objective based on the association between log valuation errors and log sale prices, and characterize the cross-observation coupling induced by the Direct squared-covariance penalty. Second, we derive a sample-additive upper-bound Surrogate and show that it is exactly a weighted squared-error objective whose weights increase with squared log-price distance from the training-sample center, making both its implementation and its distinct mechanism transparent. Third, within a translated CCAO residential workflow, we trace complete chronological-validation, held-out, and forward regularization paths and evaluate them jointly using predictive performance, assessor-facing vertical-equity measures, mechanism and residual-structure diagnostics, value-group ratio profiles, and a centered-recalibration benchmark. Exploratory multi-jurisdiction analyses separately examine transferability, heterogeneity, and failure modes.}
% \newtext{Our contribution is application-driven and has three parts. First, we translate an assessor-facing regressivity pattern into a trainable first-order objective based on the association between log valuation errors and log sale prices, and characterize the cross-observation coupling and fixed-space rank-one geometry induced by the Direct squared-covariance penalty. Second, we derive a sample-additive upper-bound Surrogate and show that it is exactly a weighted squared-error objective whose weights increase with squared log-price distance from the training-sample center; a parallel fixed-space analysis characterizes it as a weighted projection with a generally multi-directional correction path. These results make clear that the Surrogate is related to, but is not merely an implementation approximation of, the Direct penalty. Third, within a translated CCAO residential workflow, we trace complete chronological-validation, held-out, and forward regularization paths and evaluate them jointly using predictive performance, assessor-facing vertical-equity measures, mechanism and residual-structure diagnostics, and value-group ratio profiles. Exploratory multi-jurisdiction analyses separately examine transferability, heterogeneity, and failure modes.}
\newtext{
Our contribution is application-driven and has three parts. First, we turn the value-related pattern associated with regressivity into a training objective based on the overall linear association between log valuation errors and log sale prices. For the Direct squared-covariance penalty, we show how observations are linked through the sample-wide covariance and, in a theoretical benchmark that holds the prediction space fixed, why the correction follows a single direction (a rank-one path). Second, we derive a Surrogate based on an upper bound on squared covariance that separates across observations. We show that the resulting training objective is exactly weighted squared error: the farther a sale's log price is from the training-sample average, the more weight its squared error receives. In the same fixed-space benchmark, the Surrogate is a weighted projection whose correction can combine multiple directions. The Surrogate is therefore related to the Direct penalty but is a distinct objective, not simply an approximation used to implement it. Third, in a research translation of the CCAO residential workflow, we trace the full Direct and Surrogate regularization paths through chronological validation and evaluate them on a held-out sample and a separate forward sample. We assess predictive performance and assessor-facing vertical equity, and use additional diagnostics and value-group ratio profiles to examine how the targeted association changes and what value-related structure remains. Separately, exploratory analyses across multiple jurisdictions examine transferability, heterogeneity, and failure modes.
}
% \oldtext{We report complete regularization paths rather than a selected penalty.}
\newtext{Throughout the primary CCAO analysis, we report the complete regularization paths rather than select a deployment operating point.}

% \oldtext{The remainder of this section reviews the related literature on machine learning in mass appraisal, vertical equity, and dependence-based fair regression. The rest of the paper is organized as follows. Section~\ref{sec:setting} describes the mass-appraisal setting, the CCAO valuation workflow, and the evaluation framework. Section~\ref{sec:method} develops the proposed covariance-based regressivity correction. The remaining sections present the empirical evaluation, examine the robustness and practical implementation of the approach, and discuss its limitations and extensions.}
\newtext{The remainder of this section reviews related work on machine learning in mass appraisal, vertical equity, and dependence-based fair regression. Section~\ref{sec:setting} then defines the CCAO setting and evaluation framework, Section~\ref{sec:method} develops the covariance-guided correction, and the remaining sections describe the empirical design, report the regularization-path evidence, and discuss practical implications, limitations, and extensions.}

% The baseline motivation figure is introduced after the notation and diagnostics in Section~\ref{subsec:baseline_tension}; no separate Introduction figure is needed.


% % \newpage
% \subsection{Related Work}
% \label{subsec:related_work}

% % \paragraph{Machine learning in mass appraisal.}
% % Research and assessor practice have moved from traditional regression toward tree ensembles and other nonlinear methods \cite{WangLi2019,Rootetal2023,Sevgen2024100111,Zilli2024}. Their use does not remove the need for sales verification, temporal testing, ratio studies, interpretation, and administrative review \cite{IAAO2020SalesVerification,IAAOAITaskForce2022,BidansetRakow2022}. The CCAO pipeline follows this sequence: model training is followed by assessment, evaluation, interpretation, and final review \cite{ccao_data_2026}.

% % {
% % \color{red}
% \paragraph{Machine Learning in Mass Appraisal} Recent reviews document the growing use of automated valuation models and machine learning (ML) in mass appraisal and real-estate valuation. Multiple regression-based models remain common, alongside with tree ensembles (e.g., Random Forest and Gradient Boosting), and neural-network-based methods \cite{WangLi2019,Rootetal2023,Zilli2024,purdueAssessingTechniques}. Most studies compare these methods primarily through predictive criteria, such as out-of-sample accuracy and performance relative to alternative algorithms. On the other hand, equity is less often incorporated directly into model training, model selection, and performance assessment \cite{WangLi2019,Rootetal2023,Zilli2024}. In terms of predictive  power alone, tree ensembles are among the most frequently studied and often strongest-performing nonlinear methods in this literature \cite{Rootetal2023,Sevgen2024100111,Zilli2024}. Therefore, a practical equity-aware intervention and assessment should be compatible with these already widely used model classes rather than restrict solutions to a specialized class of valuation models.

% Even though adoption of ML models is increasing, it is important to note that their adoption by assessment offices remains partial and uneven. Among the perceived adoption barriers, \textcite{BidansetRakow2022} identify
% factors such as: communication challenges, limited technical capacity, and budget constraints. Further, for assessor's offices, the predictive modeling stage is only one part, of many, involved in the assessment process; the use of ML does not replace sales verification, out-of-time evaluation, ratio studies, model interpretation, administrative review, and further adjustments \cite{IAAO2020SalesVerification,IAAOAITaskForce2022}. As one of the main assessors organizations, the \textcite{IAAOAITaskForce2022} emphasizes explainability, defensibility, administrative feasibility, and integration into existing workflows as central conditions for operational use of advanced ML models. These considerations favor equity-aware methods that can be incorporated into already existing regression objectives and model-development pipelines used by assessors, with limited disruption. This practical requirement motivates the approach developed in this paper.
% % }

% % \paragraph{Vertical equity and regressivity.}
% % \advisorchange{The assessor literature distinguishes horizontal dispersion from systematic value-related bias \cite{IAAO2013Ratio,cornia2006property,McMillenSingh2023}. Assessor standards use several complementary diagnostics because no single statistic fully describes vertical equity. Large samples can hide local pockets of bias, and regression-based diagnostics can behave differently when assessed values are themselves produced by regression models \cite{IAAO2013Ratio,IAAO2023VerticalReview,McMillenSingh2023}.} Proposed corrections include spatial and geographically weighted methods, stratified procedures, ratio-study-based validation, and fairness-aware valuation models \cite{Quintos2014,BidansetEtAl2019,BidansetEtAl2025,HermansEtAl2022,HermansEtAl2025,CandoganHanLu2023}. \advisorchange{The correction developed here adds a global price-related target during training} and leaves the standard ratio-study review in place.

% % {\color{red}
% \paragraph{Vertical equity and regressivity.}
% Equity has long been a central objective in mass appraisal practice and standards. The assessor literature distinguishes \emph{horizontal equity}, which concerns the uniform treatment of similar properties, from \emph{vertical equity}, which concerns systematic differences in assessment outcomes across property-value levels \cite{IAAO2013Ratio,cornia2006property,McMillenSingh2023}.  A common form of vertical inequity is \emph{regressivity}: assessment-to-sale ratios tend to decline with property value, so lower-valued properties tend to be relatively over-assessed and higher-valued properties tend to be relatively under-assessed. Ratio studies evaluate these patterns through stratified analyses and complementary measures of uniformity and price-related bias, such as the Coefficient of Dispersion (COD), the Price-Related Differential (PRD), and the Price-Related Bias (PRB) \cite{IAAO2013Ratio,IAAO2023VerticalReview}.

% From the ratio studies measures definitions, one can note that no single statistic fully characterizes vertical equity. Aggregate results of these metrics can conceal bias within geographic or property subgroups, and some regression-based diagnostics can behave poorly, in terms of equity quantification, when estimated values are generated by regression models \cite{IAAO2023VerticalReview,McMillenSingh2023}. Accordingly, assessor guidance favors examining several diagnostics all together, instead of relying on single global measures. Empirical studies further show that regressivity is widespread, that its consequences are of substantial importance, and that limitations in assessment data and valuation models themselves are among its likely causes \cite{Berry2021,AvenancioLeonHoward2022}.

% A smaller body of literature moves beyond measurement toward explicit correction of inequities within the valuation models. Proposed approaches include spatial and geographically weighted methods, stratified procedures, ratio-study-based model validation, and fairness-aware valuation models \cite{Quintos2014,BidansetEtAl2019,BidansetEtAl2025,HermansEtAl2022,HermansEtAl2025,CandoganHanLu2023}. The method developed here contributes to this strand by adding a global price-related equity-aware target during training procedures, while retaining standard ratio studies and subgroup review for its out-of-sample evaluation.
% % }

% % \paragraph{Fair regression and dependence penalties.}
% % Fair-regression methods use constrained estimation, post-processing, or dependence penalties for continuous outcomes \cite{AgarwalDudikWu2019,FitzsimonsEtAl2019,ChzhenEtAl2020Wasserstein,ChzhenEtAl2020Recalibration,WeiEtAl2023}. Prior work has used covariance, correlation, kernel dependence, and maximal-correlation penalties \cite{KomiyamaEtAl2018,MaryEtAl2019,PerezSuayEtAl2017,LiEtAl2022Fairness,ScutariPaneroProissl2022,lee2022maximal}. We use that machinery for an assessor-specific target: the value-related pattern in assessment ratios. The paper then studies implementation, temporal calibration, and failure modes in mass-appraisal models.

% % {
% % \color{red}
% \paragraph{Fair regression and dependence-based constrained learning.}
% Although much of the fairness-in-ML literature initially focused on classification, \emph{fair regression} studies how fairness requirements can be incorporated when predictions are continuous, as it is in mass appraisal. Existing approaches include constrained or regularized empirical risk minimization and post-processing methods. These methods impose fairness-aware criteria to the ML pipelines, such as statistical parity, bounded group loss, distributional parity through optimization hard constraints, soft-constraint through penalties, and recalibration procedures \cite{AgarwalDudikWu2019,FitzsimonsEtAl2019,ChzhenEtAl2020Wasserstein,ChzhenEtAl2020Recalibration,WeiEtAl2023}. The broader methodological principle is that predictive error loss (e.g., sum of the squares of the residuals) does not need to be the only single training objective: an application-specific requirement can be added as a hard constraint or soft penalty to the learning procedure, which in most cases produces an explicit tradeoff between predictive performance and the relevant notion of fairness.

% A particularly relevant literature for this paper represents unfairness as statistical dependence between model outputs and a variable of concern. Prior work has controlled this dependence using covariance or correlation, coefficients of determination, or kernel-based measures such as Hilbert-Schmidt Independence Criterio (HSIC), and maximal correlation \cite{KomiyamaEtAl2018,MaryEtAl2019,PerezSuayEtAl2017,LiEtAl2022Fairness,ScutariPaneroProissl2022,lee2022maximal}. These methods are especially useful when the relevant variable is continuous and can often be incorporated into standard regression objectives. However, our application differs from conventional group-fairness settings: property value is not treated as a protected attribute since it is the target itself. Instead, we use dependence regularization to encode an assessor-specific structural requirement by penalizing the property value-related pattern in assessment ratios associated with regressivity. In this sense, the proposed method is a form of domain-informed constrained learning that translates a standard ratio-study concern into a trainable objective.
% % }







\subsection{Related Work}
\label{subsec:related_work}

\paragraph{Machine Learning in Mass Appraisal}
Recent reviews document the {growing research on} automated valuation models and machine learning (ML) in mass appraisal and {real estate} valuation. {Multiple regression remains common, alongside tree ensembles such as random forests and gradient boosting, as well as neural-network-based methods} \cite{WangLi2019,Rootetal2023,Zilli2024,purdueAssessingTechniques}. Most studies compare these methods primarily through predictive criteria, such as out-of-sample accuracy and performance relative to alternative algorithms. {Vertical equity, although routinely evaluated in assessor practice, is much less often incorporated directly into model training or model selection} \cite{WangLi2019,Rootetal2023,Zilli2024}. {When predictive performance is the primary criterion, tree ensembles are among the most frequently studied and often strong-performing nonlinear methods in this literature} \cite{Rootetal2023,Sevgen2024100111,Zilli2024}. {A practical training-time equity intervention should therefore be compatible with these commonly studied and used model classes rather than require a specialized valuation architecture.}

 {Operational adoption by assessment offices nevertheless remains partial and uneven.} \textcite{BidansetRakow2022} identify {communication challenges, limited technical capacity, and budget constraints among perceived barriers to wider adoption}. {Moreover, predictive modeling is only one component of the assessment workflow; the use of ML does not replace sales verification, out-of-time evaluation, ratio studies, model interpretation, or administrative review} \cite{IAAO2020SalesVerification,IAAOAITaskForce2022}. {The} \textcite{IAAOAITaskForce2022} emphasizes explainability, defensibility, administrative feasibility, and integration into existing workflows as central conditions for operational use of advanced ML models. {These considerations favor equity-aware methods that can be incorporated into existing model-training, validation, and review pipelines with minimal disruption.}

\paragraph{Vertical equity and regressivity.}
Equity has long been a central objective in mass appraisal practice and standards. The assessor literature distinguishes \emph{horizontal equity}, which concerns the uniform treatment of similar properties, from \emph{vertical equity}, which concerns systematic differences in assessment outcomes across property-value levels \cite{IAAO2013Ratio,cornia2006property,McMillenSingh2023}. A common form of vertical inequity is \emph{regressivity}: assessment-to-sale ratios tend to decline with property value, so lower-valued properties tend to be relatively overassessed and higher-valued properties tend to be relatively underassessed. {Ratio studies evaluate assessment uniformity using complementary diagnostics. For example, the Coefficient of Dispersion (COD) summarizes dispersion in assessment ratios, whereas the Price-Related Differential (PRD) and Price-Related Bias (PRB) are commonly used to diagnose price-related vertical inequity} \cite{IAAO2013Ratio,IAAO2023VerticalReview}.

% \textcolor{blue}{\st{From the ratio studies measures definitions, one can note that no single statistic fully characterizes vertical equity. Aggregate results of these metrics can conceal bias within geographic or property subgroups, and some regression-based diagnostics can behave poorly, in terms of equity quantification, when estimated values are generated by regression models.}}

{No single diagnostic fully characterizes vertical equity. Methodological studies show that vertical-equity measures can differ materially in their statistical behavior, and that some regression-based diagnostics can be biased when assessed values are themselves generated by regression-based valuation models} \cite{IAAO2023VerticalReview,McMillenSingh2023}. {Accordingly, the assessor literature recommends examining multiple complementary diagnostics rather than relying on a single global measure.}  {Empirical studies further document substantial regressivity and consequential distributional disparities across jurisdictions} \cite{Berry2021,AvenancioLeonHoward2022}. {Recent large-scale evidence also shows that predictive accuracy and assessment equity need not be intrinsically opposed: improvements in the information available to valuation models can, in many settings, improve both simultaneously \cite{SmithEtAl2026}.}

{A smaller literature moves beyond aggregate diagnosis toward model specification, localized validation, and explicit correction intended to improve assessment equity.} Proposed approaches include spatial and geographically weighted methods, stratified procedures, ratio-study-based model validation, and fairness-aware valuation models \cite{Quintos2014,BidansetEtAl2019,BidansetEtAl2025,HermansEtAl2022,HermansEtAl2025,CandoganHanLu2023}. 
% \textcolor{blue}{\sout{Most closely related to our work, Candogan, Han, and Lu develop a K-segment property-valuation model that explicitly targets regressivity and studies the associated accuracy--fairness relationship. Our approach addresses the same assessor-side objective through a different mechanism: rather than introducing a segmented valuation architecture, we incorporate a dependence-based penalty directly into the training objective of existing regression learners. The method developed here therefore contributes to this strand by directly targeting a global price-related dependence associated with regressivity during training, while retaining standard ratio studies and subgroup review for out-of-sample evaluation.}}
{A closely related approach is \textcite{CandoganHanLu2023}, who develop a \emph{K-segment} property-valuation model that explicitly targets regressivity and studies the associated accuracy--fairness relationship. Our approach addresses the same assessor-side objective through a different mechanism: rather than introducing a segmented valuation architecture, we use a first-order log-error and log-sale-price association as a training target within an existing regression learner. The method therefore targets a deliberately narrow first-order mechanism during training while retaining standard ratio studies, proxy-aware vertical-equity diagnostics, and subgroup review for out-of-sample evaluation.}

\paragraph{Fair regression and dependence-based constrained learning.}

% \oldtext{Although much of the fairness-in-ML literature initially focused on classification, \emph{fair regression} studies how fairness requirements can be incorporated when predictions are continuous, as in property valuation. Existing approaches include constrained or regularized empirical risk minimization and post-processing methods. Different works consider fairness notions such as statistical parity, bounded group loss, mean parity, or distributional parity. Methodologically, these requirements can be imposed through constrained or regularized empirical risk minimization or through post-processing and recalibration procedures \cite{AgarwalDudikWu2019,FitzsimonsEtAl2019,ChzhenEtAl2020Wasserstein,ChzhenEtAl2020Recalibration,WeiEtAl2023}.}

\newtext{Although much of the fairness-in-ML literature initially focused on classification, \emph{fair regression} studies how application-specific fairness requirements can be incorporated when predictions are continuous, as in property valuation. Existing approaches consider criteria such as statistical parity, bounded group loss, mean parity, or distributional parity and implement them through constrained or regularized empirical risk minimization, post-processing, or recalibration \cite{AgarwalDudikWu2019,FitzsimonsEtAl2019,ChzhenEtAl2020Wasserstein,ChzhenEtAl2020Recalibration,WeiEtAl2023}.} The broader methodological principle is that predictive risk can be optimized subject to, or jointly with, an application-specific fairness criterion. Varying the strength of this requirement produces a family of solutions that makes the relationship between predictive performance and the targeted fairness criterion explicit, without requiring that improvements in one necessarily worsen the other.

{A particularly relevant strand expresses fairness requirements as restrictions on statistical dependence between predictions and a sensitive variable.} {Prior work has quantified or controlled such dependence using correlation and coefficients of determination \cite{KomiyamaEtAl2018}, kernel-based measures such as the Hilbert--Schmidt Independence Criterion (HSIC) \cite{PerezSuayEtAl2017,LiEtAl2022Fairness}, and maximal-correlation criteria \cite{MaryEtAl2019,lee2022maximal}. Related regularization approaches also control the contribution of sensitive attributes to regression predictions \cite{ScutariPaneroProissl2022}.} 
% \textcolor{blue}{\sout{These approaches are particularly relevant when the variable of concern is continuous and can often be incorporated directly into regression objectives. Our application, however, differs from conventional group-fairness settings. Sale price, our observed proxy for market value, is the prediction target rather than a protected attribute, and our objective is not demographic parity. We instead adapt the general dependence-control principle to a different structural relationship: dependence between valuation errors and property value. The assessor literature therefore supplies the domain-specific target---price-related regressivity---while the fair-regression literature supplies the general dependence-control mechanism. The proposed method connects the two by translating value-related dependence in valuation errors into a trainable regularization objective.}}
{These approaches are particularly relevant when the variable of concern is continuous and can often be incorporated directly into regression objectives. Our application, however, differs from conventional group-fairness settings. Sale price is the observed transaction-based proxy for market value and the prediction target, not a protected attribute, and our objective is not demographic parity. We adapt the broader dependence-control principle to a narrower structural target: the signed first-order association between log valuation errors and log sale prices. The assessor literature supplies the domain-specific objective---reducing value-related regressivity---while the fair-regression literature supplies a general mechanism for incorporating statistical associations into model training.}










\section{Preliminaries: The CCAO Baseline and Empirical Motivation}
\label{sec:setting}

% \textcolor{blue}{\sout{This section fixes the CCAO application setting, prediction notation, and evaluation criteria used throughout the paper. We first describe the residential valuation workflow and baseline model classes, then distinguish predictive performance from assessor-facing ratio-study diagnostics, and finally present the controlled baseline comparison that motivates the proposed correction.}}

% \textcolor{blue}{\sout{This section fixes the CCAO application setting, prediction notation, and evaluation framework used throughout the paper. We first describe the residential valuation workflow and baseline model classes. We then distinguish predictive performance, assessor-facing measures of assessment level and ratio dispersion, complementary vertical-equity diagnostics, and supplemental checks for nonlinear or localized value-related patterns. Finally, we present the controlled baseline comparison that motivates the proposed correction.}}
% \oldtext{This section fixes the CCAO application setting, prediction notation, and evaluation framework used throughout the paper. We first describe the residential valuation workflow and baseline model classes. We then distinguish predictive performance, assessor-facing measures of valuation level and ratio dispersion, complementary vertical-equity diagnostics, and supplemental checks for first-order and nonlinear sale-price-related patterns. Finally, we present the controlled baseline comparison that motivates the proposed correction.}

\newtext{This section fixes the CCAO application setting, prediction notation, and evaluation framework used throughout the paper. We first describe the residential valuation workflow and baseline model classes. We then distinguish predictive performance, assessor-facing measures of valuation level and ratio dispersion, complementary vertical-equity diagnostics, and mechanism and residual-structure diagnostics for first-order, nonlinear conditional-mean, and broader dependence patterns. Finally, we present the controlled baseline comparison that motivates the proposed correction.}


\subsection{The CCAO Residential Valuation Workflow}
\label{subsec:ccao_workflow}

The CCAO values a large and heterogeneous residential inventory on a triennial reassessment cycle \cite{CCAOAssessmentCalendar2026}. Its residential modeling evolved from township-level linear-regression models, with linear-regression and gradient-boosting alternatives used in 2019--2020, to a county-wide LightGBM model beginning in 2021 \cite{CCAOModelResAVM2026}. LightGBM is an implementation of gradient-boosted decision trees \cite{lightgbm}.

The CCAO's current published residential automated valuation model (AVM) uses verified sales and a broad set of property, spatial, neighborhood, market, and contextual predictors to estimate residential market values \cite{CCAOModelResAVM2026}. The prediction model is embedded in a broader operational workflow comprising data preparation, model training, valuation, evaluation, interpretation, and final review. Candidate valuations are evaluated using both predictive-performance measures and assessor-specific ratio-study statistics across multiple geographic areas and property classes.

% \textcolor{blue}{\sout{For the empirical analysis, we compare a linear-regression benchmark motivated by CCAO's earlier models with standard LightGBM, representing its current model class, while holding fixed the prediction target, underlying feature information, development sample, temporal design, and evaluation procedures.}}

{For the empirical analysis, we compare a linear-regression benchmark motivated by CCAO's earlier models with unpenalized LightGBM, representing its current model class, while holding fixed the prediction target, underlying feature information, development sample, temporal design, and evaluation procedures.}
The comparison is therefore a controlled comparison of model classes rather than a reconstruction of official historical assessments. The exact sample construction, transaction filters, and temporal validation procedures are described in Section~\ref{subsec:ccao_design}.


\subsection{Prediction Setting and Baseline Models}
\label{subsec:baseline_models}

Let $\mathcal{I}$ denote the set of residential sales retained after the CCAO sales-validation and transaction-eligibility procedures described in Section~\ref{subsec:ccao_design}. For each $i\in\mathcal{I}$, we observe
\[
(x_i,P_i,t_i),
\]
where $x_i$ contains property and parcel characteristics; tract-level socioeconomic and demographic variables from the U.S. Census Bureau's American Community Survey (ACS); and geographic, neighborhood, environmental, accessibility, market, and temporal features.

% \textcolor{blue}{\sout{The variable $P_i>0$ is the verified sale price and serves as the transaction-based benchmark for market value, while $t_i$ is the sale date used to construct chronological samples and time-related predictors.}}

% \textcolor{blue}{\sout{The variable $P_i>0$ is the verified sale price and is treated as the observed transaction-based proxy for the property's market value, while $t_i$ is the sale date used to construct chronological samples and time-related predictors. Market value itself is latent: a verified sale therefore provides an observed market-value indicator rather than a directly observed market value. The prediction problem is to estimate market value from the information available in $x_i$.}}

{The variable $P_i>0$ is the verified sale price and is treated as the observed transaction-based proxy for the property's market value, while $t_i$ denotes the sale date used to construct the chronological samples \cite{CCAOModelResAVM2026,IAAO2026ExposureRatio}. Any time-derived predictors used by the models are included in $x_i$. Market value itself is latent: a verified sale provides an observed market-value indicator rather than a directly observed market value. The prediction problem is to estimate market value from the information available in $x_i$.}

Since residential sale prices are positive and highly right-skewed, the baseline models are trained on the log-price scale. Define
\begin{equation}
\label{eq:prediction_scale}
y_i=\log P_i,
\qquad
\widehat{y}_i=f(x_i),
\qquad
e_i(f)=\widehat{y}_i-y_i,
\end{equation}
% \textcolor{blue}{\sout{where $f\in\mathcal{H}$ is a candidate prediction function.}}
{where $f\in\mathcal{H}$ is a candidate prediction function and $e_i(f)$ is the \emph{log valuation error}, defined with prediction minus observed log price so that it coincides with the log valuation ratio below.}
The corresponding prediction on the original price scale is
\begin{equation}
\label{eq:price_prediction}
\widehat{P}_i=\exp\bigl(f(x_i)\bigr).
\end{equation}

% \textcolor{blue}{\sout{The log-scale is the estimation scale, while $\widehat{P}_i$ is the market-value prediction used for price-scale predictive evaluation and ratio-study diagnostics. Because both $\widehat{P}_i$ and $P_i$ are expressed on the market-value scale, define}}

% \textcolor{blue}{\sout{The log-price scale is the estimation scale, whereas $\widehat{P}_i$ is the directly retransformed valuation used for price-scale predictive evaluation and ratio-study diagnostics. Throughout the analysis, $\widehat{P}_i$ is defined operationally by Eq.~\eqref{eq:price_prediction}. Note that the direct exponentiation should not be interpreted as a claim that $\widehat{P}_i$ is a bias-corrected estimate of the expected value of property $i$ given $x_i$. We define the corresponding valuation-to-sale ratio as}}

{The log-price scale is the estimation scale, whereas $\widehat{P}_i$ is the directly retransformed valuation used for price-scale predictive evaluation and ratio-study diagnostics. Throughout the analysis, $\widehat{P}_i$ is defined operationally by~\eqref{eq:price_prediction}; direct exponentiation is not presented as a bias correction for the conditional mean on the original price scale. We define the corresponding \emph{valuation-to-sale ratio} as}
\begin{equation}
\label{eq:valuation_ratio}
r_i(f)
=
\frac{\widehat{P}_i}{P_i}
=
\exp\bigl(e_i(f)\bigr).
\end{equation}

% \textcolor{blue}{\sout{For brevity and consistency with ratio-study terminology, we refer to $r_i(f)$ as the assessment ratio throughout the paper. Its numerator is the model-estimated market value, not Cook County's statutory assessed value after downstream adjustments.}}

% \textcolor{blue}{\sout{For brevity and consistency with ratio-study terminology, we subsequently refer to $r_i(f)$ as the assessment ratio. Its numerator is the model-estimated market value, not Cook County's statutory assessed value after downstream adjustments. When a fitted model is fixed, we suppress the dependence on $f$ and write simply $r_i$ and $e_i$. The log-scale residual and the assessment ratio are therefore linked exactly by}}

{For brevity, we subsequently refer to $r_i(f)$ as the \emph{valuation ratio}. Its numerator is the model-estimated market value, not Cook County's statutory assessed value after downstream adjustments. When a fitted model is fixed, we suppress the dependence on $f$ and write simply $r_i$ and $e_i$. The log valuation error and valuation ratio are linked exactly by}
\begin{equation}
\label{eq:log_ratio}
e_i(f)
=
\log \widehat{P}_i-\log P_i
=
\log r_i(f).
\end{equation}
Thus, $e_i(f)>0$ (equivalently, $r_i(f)>1$) corresponds to a prediction above the observed sale price, whereas $e_i(f)<0$ (equivalently, $r_i(f)<1$) corresponds to a prediction below it.
% \textcolor{blue}{\sout{Estimation on the log-price scale therefore converts multiplicative price errors into additive residuals, and the training residual is exactly the log assessment ratio.}}
{Estimation on the log-price scale therefore converts multiplicative price errors into additive log valuation errors, and $e_i(f)$ is exactly the log valuation ratio.}

For the predictive models, we consider two baseline model classes: linear regression and gradient-boosted decision trees. As described above, these choices provide a historically motivated linear benchmark and a representation of CCAO's current model class. Between models, we hold fixed the prediction target, development observations, underlying feature information, temporal design, and evaluation protocol; the models differ in functional form and model-specific estimation procedure.

For a training sample of size $n$, both baselines are fitted to reduce mean squared log-price error, i.e.,
\begin{equation}
\label{eq:baseline_learning}
L(f)
=
\frac{1}{n}\sum_{i=1}^n e_i(f)^2,
\end{equation}
while differing in functional form, estimation procedure, and model-specific complexity control. This loss is the predictive component of the training objective modified in Section~\ref{sec:method}.

For linear regression,
\begin{equation}
\label{eq:linear_model}
f_{\mathrm{lin}}(x)
=
\beta_0+x^\top\beta.
\end{equation}
The hypothesis class is therefore restricted to functions that are linear in the supplied predictors. In the ordinary linear-regression baseline, no additional regressivity term is included in the training objective. This specification provides the controlled linear benchmark motivated by CCAO's earlier linear-regression models. Its linear structure imposes a comparatively rigid relationship between the predictors and log sale price.

% \textcolor{blue}{%
  % \sout{LightGBM, the modeling foundation of CCAO's current published residential AVM}%%
  % \sout{, constructs an additive ensemble of regression trees.}%
% }

{On the other hand, LightGBM, the modeling foundation of CCAO's current published residential AVM, constructs an additive ensemble of regression trees \cite{CCAOModelResAVM2026,lightgbm}. A gradient-boosted tree predictor can be written as}
\begin{equation}
\label{eq:boosted_tree_model}
f_{\mathrm{gbm}}^{(M)}(x)
=
f^{(0)}(x)
+
\sum_{m=1}^{M}\nu b_m(x),
\end{equation}
where $f^{(0)}$ is an initial prediction, $b_m$ is the regression tree added at boosting iteration $m$, $M$ is the number of boosting iterations, and $\nu>0$ is the learning rate. Unlike the linear model, the tree ensemble can represent nonlinearities, threshold effects, and interactions without requiring them to be specified in advance.

LightGBM fits the ensemble in~\eqref{eq:boosted_tree_model} sequentially. The gradient and Hessian calculations required for the proposed custom objectives are introduced with the penalties in Section~\ref{sec:method}.

The two baselines therefore solve the same prediction problem under a common empirical design but use substantially different model classes.
% \textcolor{blue}{\sout{Both baseline objectives minimize squared log-price error and contain no term that directly controls systematic value-related patterns in residuals or assessment ratios. We next distinguish predictive performance from the assessor-facing criteria used to evaluate assessment level, ratio dispersion, and regressivity.}}
{Both baseline objectives minimize squared log-price error and contain no term that directly controls systematic value-related patterns in log valuation errors or valuation ratios. We next distinguish predictive performance from the assessor-facing criteria used to evaluate valuation level, ratio dispersion, and regressivity.}


\subsection{
\texorpdfstring{
% \oldtext{Predictive and Assessor-Facing Evaluation Criteria}
\newtext{Predictive, Assessor-Facing, and Residual-Structure Diagnostics}}{Predictive, Assessor-Facing, and Residual-Structure Diagnostics}}
\label{subsec:metrics}

% \textcolor{blue}{\sout{The two scales play distinct roles in evaluation. Predictive performance is assessed both on the log-price scale used for estimation and on the original price scale on which market values are reported. The assessor-facing ratio-study diagnostics are computed on the original price scale.}}

{The two scales play distinct roles in evaluation. Predictive performance is assessed both on the log-price scale used for estimation and on the original price scale on which valuations are reported. Assessor-facing ratio-study diagnostics are computed on the original price scale. The resulting measures answer different questions and are therefore interpreted jointly rather than collapsed into a single performance criterion.}

\paragraph{Predictive performance.}
% \textcolor{blue}{\sout{For an evaluation sample $S\subseteq\mathcal{I}$, let $n_S=|S|$, $\bar P_S=n_S^{-1}\sum_{i\in S}P_i$, and $\bar y_S=n_S^{-1}\sum_{i\in S}y_i$. We use three complementary headline measures: price-scale mean absolute error (MAE), price-scale $R^2$, and log-scale root mean squared error (RMSE).}}
{For an evaluation sample $S\subseteq\mathcal{I}$, let $n_S=|S|$ and $\bar P_S=\frac{1}{n_S}\sum_{i\in S}P_i$. We use four complementary headline accuracy-related measures: price-scale mean absolute error (MAE), price-scale mean absolute percentage error (MAPE), price-scale $R^2$, and log-scale root mean squared error (RMSE):}
\begin{align}
\label{eq:predictive_metrics}
\operatorname{MAE}_{P}(S)
&=
\frac{1}{n_S}\sum_{i\in S}\left|\widehat P_i-P_i\right|,
\\
{\operatorname{MAPE}_{P}(S)}
&{=100\times
\frac{1}{n_S}
\sum_{i\in S}
\left|
\frac{\widehat P_i-P_i}{P_i}
\right|
=100\times 
\frac{1}{n_S}\sum_{i\in S}|r_i-1|},
\\
R^2_{P}(S)
&=
1-
\frac{\sum_{i\in S}(\widehat P_i-P_i)^2}
{\sum_{i\in S}(P_i-\bar P_S)^2},
\\
\operatorname{RMSE}_{\log P}(S)
&=
\left(
\frac{1}{n_S}\sum_{i\in S}e_i(f)^2
\right)^{1/2}.
\end{align}
% \textcolor{blue}{\sout{The price-scale MAE gives the average absolute prediction error in dollars, while $R^2_P$ summarizes predictive fit on the market-value scale. The log-scale RMSE directly evaluates the loss scale used for estimation and, by~\eqref{eq:log_ratio}, can equivalently be viewed as the root mean squared log-ratio error. When useful, we also report the supplementary log-scale coefficient of determination $R^2_{\log P}$.}}

% \textcolor{blue}{\sout{The price-scale MAE gives the average absolute prediction error in dollars, while MAPE gives the average absolute error relative to the observed sale price. In particular, MAPE measures deviation from the economically neutral ratio $r_i=1$ and is therefore distinct from COD below, which measures dispersion around the sample median ratio. The statistic $R^2_P$ summarizes predictive fit on the original price scale. Finally, $\operatorname{RMSE}_{\log P}$ directly evaluates the scale and squared-error loss used for estimation and, by Eq.~\eqref{eq:log_ratio}, is equivalently the root mean squared log-ratio error.}}

{The price-scale MAE gives the average absolute prediction error in dollars, while MAPE gives the average absolute error relative to the observed sale price. In particular, MAPE measures absolute deviation from the unit-ratio benchmark $r_i=1$ and is therefore distinct from COD below, which measures dispersion around the sample median ratio. The statistic $R^2_P$ summarizes predictive fit on the original price scale. Finally, $\operatorname{RMSE}_{\log P}$ directly evaluates the squared-error scale used for estimation and, by~\eqref{eq:log_ratio}, is equivalently the root mean squared log valuation ratio.}
% \textcolor{blue}{\sout{Because $R^2_{\log P}$ and $\operatorname{RMSE}_{\log P}$ are driven by the same log-scale sum of squared errors on a fixed evaluation sample, $R^2_{\log P}$ is treated as supplementary rather than as a separate headline criterion. Percentage errors are also not used as headline measures: for example, on the price scale, the mean absolute percentage error (MAPE) equals $\frac{100}{n_S}\sum_{i\in S}|r_i-1|$ and therefore overlaps directly with the ratio-based evaluation below, while percentage errors of log-prices do not have a stable scale interpretation.}}
% \textcolor{impGreen}{We retain $R^2_{\log P}$ only as a supplementary implementation diagnostic. On a fixed evaluation sample it is driven by the same log-scale sum of squared errors as $\operatorname{RMSE}_{\log P}$ and is therefore not treated as an additional headline criterion.}


% \paragraph{\textcolor{blue}{\sout{Assessment level.}}
\paragraph{Valuation level.}
Recall from~\eqref{eq:valuation_ratio} that $r_i=\widehat P_i/P_i$. A ratio above one indicates that the predicted value exceeds the observed sale price, whereas a ratio below one indicates that it falls below the sale price. Define the median, arithmetic mean, and price-weighted mean ratios as
\begin{equation}
\label{eq:level_ratios}
m_S
=
\operatorname{median}_{i\in S}\{r_i\},
\qquad
\bar r_S
=
\frac{1}{n_S}\sum_{i\in S}r_i,
\qquad
\bar r_{W,S}
=
\frac{\sum_{i\in S}\widehat P_i}{\sum_{i\in S}P_i}
=
\frac{\sum_{i\in S}r_iP_i}{\sum_{i\in S}P_i}.
\end{equation}

% \textcolor{blue}{\sout{These quantities summarize aggregate assessment level. Because both numerator and denominator of $r_i$ are expressed on the market-value scale in this analysis, the neutral reference value for the three statistics is a value of one. The median is relatively insensitive to extreme ratios, the arithmetic mean assigns equal weight to each sale, and the weighted mean places greater weight on higher-value properties. Agreement among these measures indicates broad aggregate calibration but does not establish either horizontal or vertical equity.}}

{These quantities summarize aggregate valuation level. Because both numerator and denominator of $r_i$ are expressed on the market-value scale in this analysis, the neutral reference value for all three statistics is one. The median is relatively insensitive to extreme ratios, the arithmetic mean assigns equal weight to each sale, and the weighted mean places greater weight on higher-value properties. Agreement among these measures indicates broad aggregate calibration but does not establish either horizontal uniformity or vertical equity.}

% \textcolor{impGreen}{Any numerical assessment-level bands used as model-selection guardrails are specified separately in Section~\ref{sec:model_selection} \todo{we erased the model selection section. Update this reference}; they are not treated here as universal substitutes for the neutral target of one.}

\paragraph{Ratio dispersion (Horizontal uniformity).}
Ratio dispersion is an important assessor-facing dimension of horizontal uniformity. We summarize it using the coefficient of dispersion (COD) and coefficient of variation (COV). The COD is
\begin{equation}
\label{eq:cod}
\operatorname{COD}(S)
=
100\times
\frac{1}{n_S}
\sum_{i\in S}
\frac{|r_i-m_S|}{m_S},
\end{equation}
and the COV is
\begin{equation}
\label{eq:cov}
\operatorname{COV}(S)
=
100\times
\frac{s_{r,S}}{\bar r_S},
\qquad
s_{r,S}
=
\left(
\frac{1}{n_S-1}
\sum_{i\in S}(r_i-\bar r_S)^2
\right)^{1/2}.
\end{equation}

% \textcolor{blue}{\sout{The COD measures the average absolute deviation from the median ratio, whereas the COV measures standard deviation relative to the mean. Lower values indicate less aggregate ratio dispersion. Applicable reference ranges depend on property class and prevailing market conditions.}}

{The COD measures the average absolute deviation from the median ratio, whereas the COV measures standard deviation relative to the mean. Lower values indicate less aggregate ratio dispersion, but the reference ranges depend on property class and prevailing market conditions, and unusually low dispersion can itself warrant diagnostic review \cite{IAAO2013Ratio,IAAO2026ExposureRatio}.}
To avoid ambiguity with the covariance penalty introduced later, $\operatorname{COV}$ always denotes the coefficient of variation, whereas statistical covariance is written as $\operatorname{Cov}(\cdot,\cdot)$.

\paragraph{Regressivity (Vertical equity).}
% \textcolor{blue}{\sout{Vertical equity concerns whether assessment ratios vary systematically with property value. A regressive pattern arises when lower-value properties tend to have higher assessment ratios than higher-value properties. A progressive pattern arises in the opposite direction. No single statistic fully characterizes vertical equity: different diagnostics summarize different aspects of the relationship between valuations, ratios, and property value and can therefore disagree. We consequently use a complementary suite of ratio-based, regression-based, distributional, and grouped diagnostics rather than treating any single statistic as definitive. The May 2026 IAAO exposure draft places primary emphasis on the grouped VEI and identifies PRB, MKI, and PRD as complementary measures. Given that the document is still in draft form, we use its proposed procedures and ranges as descriptive guidance rather than as final adopted compliance standards.}}
{Vertical equity concerns whether valuation ratios vary systematically across underlying market value. A \emph{regressive} pattern arises when lower-value properties tend to have higher valuation ratios than higher-value properties; a \emph{progressive} pattern arises in the opposite direction. Because true market value is latent, empirical vertical-equity procedures necessarily use observed sales or other market-value proxies, and the diagnostics below differ partly in how they address this measurement problem \cite{IAAO2026ExposureRatio}. No single statistic fully characterizes vertical equity: different diagnostics summarize different aspects of the relationship between valuations, ratios, and market-value proxies and can therefore disagree \cite{IAAO2023VerticalReview,IAAO2026ExposureRatio}. We consequently use a complementary suite of ratio-based, regression-based, distributional, and grouped diagnostics rather than treating any single statistic as definitive. The May 2026 IAAO exposure draft places primary emphasis on the grouped VEI and identifies PRB, MKI, and PRD as complementary measures. Because the document remains an exposure draft, we use its proposed procedures and ranges as descriptive guidance rather than as final adopted compliance standards.}

The \emph{price-related differential} (PRD) is
\begin{equation}
\label{eq:prd}
\operatorname{PRD}(S)
=
\frac{\bar r_S}{\bar r_{W,S}}.
\end{equation}
Note that, since the weighted mean ratio assigns greater weight to higher-value properties, a PRD above one tends to arise when higher-value properties have lower valuation ratios than lower-value properties. The 2013 IAAO standard provides a reference range of $[0.98,1.03]$ for the PRD: values above this range tend to indicate regressivity, whereas values below it tend to indicate progressivity \cite{IAAO2013Ratio}.

{PRD is a simple and widely used ratio-based diagnostic, but it can be particularly sensitive to influential high-value observations and extreme ratios. Therefore, we interpret it jointly with the other vertical-equity measures rather than as a standalone test \cite{IAAO2026ExposureRatio}.}

% \textcolor{blue}{\sout{We also report the price-related bias (PRB). Following the IAAO PRB methodology, define the equal-weight market-value proxy}}

{The following regression- and group-based vertical-equity diagnostics require a proxy for the property's value level. Ranking properties only by sale price can mechanically bias ratio-based vertical-equity tests toward finding regressivity because sale price also appears in the denominator of $r_i$, while using valuation alone can induce bias in the opposite direction. Following IAAO guidance, we therefore use the equal-weight market-value proxy}
\begin{equation}
\label{eq:market_value_proxy}
V_i^{\mathrm{proxy}}
=
\frac{1}{2}
\left(
P_i+\frac{\widehat P_i}{m_S}
\right),
\end{equation}
{which combines the verified sale price with the median-normalized predicted value to reduce this source of proxy-induced bias \cite{IAAO2013Ratio,IAAO2026ExposureRatio}.}

\newtext{The 2013 IAAO Standard explicitly defines the equal-weight proxy as $0.50(\mathrm{AV}/\mathrm{Median})+0.50\,\mathrm{SP}$. As of August 2026, the May 2026 ratio-studies document remains an exposure draft and the 2013 Standard remains the recommended IAAO guidance. The exposure draft's VEI appendix describes an equal-weight proxy but displays the normalized valuation term without an explicit $0.50$ multiplier. For reproducibility, all PRB and VEI results in this paper use the single equal-weight convention in Eq.~\eqref{eq:market_value_proxy}, consistent with the 2013 Standard and the exposure draft's stated equal-weight interpretation.}
\latesttext{The executed canonical metric code uses this equal-weight proxy exactly. PRB in \texttt{utils/motivation\_utils.py} computes $0.5\,\mathrm{SP}+0.5\,(\mathrm{AV}/\mathrm{Med})$, and VEI computes $0.5\,y_{\mathrm{true}}+0.5\,(y_{\mathrm{pred}}/m_S)$. Both match Eq.~\eqref{eq:market_value_proxy}. SHA256 of that file at population time is \texttt{4219a9d68e5266e300b319660f93d37ca184189ee340edb0785f40c9f3c39beb}.}

{We first report the \textit{price-related bias} (PRB).}
Let $\widetilde V_S^{\mathrm{proxy}}$ denote the median of $V_i^{\mathrm{proxy}}$ over $S$. The PRB is the slope in the simple linear regression
\begin{equation}
\label{eq:prb}
\frac{r_i-m_S}{m_S}
=
a_S
+
\operatorname{PRB}(S)\cdot
\log_2\left(
\frac{V_i^{\mathrm{proxy}}}
{\widetilde V_S^{\mathrm{proxy}}}
\right)
+
\varepsilon_i.
\end{equation}

Negative PRB values indicate that valuation ratios tend to decline as the market-value proxy increases and are therefore associated with regressivity. Values between $-0.05$ and $0.05$ are conventionally preferred \cite{IAAO2013Ratio}.

{As a complementary distributional diagnostic, we report the \emph{Modified Kakwani Index} (MKI). Let $\pi(1),\allowbreak\ldots,\pi(n_S)$ index the observations in $S$ in nondecreasing order of sale price, and let
\[
q_j=\frac{j-\tfrac12}{n_S},
\qquad
\overline{\widehat P}_S
=
\frac{1}{n_S}\sum_{i\in S}\widehat P_i.
\]
Define the \emph{concentration index} of predicted values when observations are ordered by sale price and the corresponding \emph{Gini coefficient} of sale prices as
\begin{equation}
\label{eq:mki_components}
C_{\widehat P\mid P}(S)
=
\frac{2}{n_S\overline{\widehat P}_S}
\sum_{j=1}^{n_S}
\widehat P_{\pi(j)}q_j
-1,
\qquad 
G_P(S)
=
\frac{2}{n_S\bar P_S}
\sum_{j=1}^{n_S}
P_{\pi(j)}q_j
-1,
\end{equation}
respectively. Then, the MKI is
\begin{equation}
\label{eq:mki}
\operatorname{MKI}(S)
=
\frac{C_{\widehat P\mid P}(S)}{G_P(S)}.
\end{equation}
When there is sufficient price variation, an MKI value of one indicates vertical neutrality, whereas values below one indicate a regressive direction, and values above one indicate a progressive direction \cite{IAAO2026ExposureRatio}. Because MKI orders observations using sale price, it can itself be affected by errors in the market-value proxy, which is another reason to interpret it as one component of the broader diagnostic suite rather than as a definitive test.}


% \textcolor{impGreen}{}

% \textcolor{blue}{\sout{Finally, we report the Vertical Equity Indicator (VEI), a high-versus-low value-group comparison proposed in the May 2026 IAAO exposure draft.}}

Finally, we report the \emph{Vertical Equity Indicator} (VEI), the quantile-grouped vertical-equity measure emphasized in the May 2026 IAAO exposure draft \cite{IAAO2026ExposureRatio}.
The draft varies the number of percentile groups with sample size: for the large evaluation samples used in the main countywide analysis ($n_S\geq 501$), observations are divided into deciles according to $V_i^{\mathrm{proxy}}$. Let $S_{q_1}$ and $S_{q_{10}}$ denote the lowest and highest deciles, and let $m_{S_{q_1}}$ and $m_{S_{q_{10}}}$ denote their median valuation ratios. The VEI is
\begin{equation}
\label{eq:vei}
\operatorname{VEI}(S)
=
100\times
\frac{m_{S_{q_{10}}}-m_{S_{q_1}}}{m_S}.
\end{equation}

A negative VEI indicates that higher-value properties have lower median valuation ratios than lower-value properties. The exposure draft proposes a reference band of $[-10\%,10\%]$, together with confidence-interval-based testing for point estimates outside that band.

% \textcolor{blue}{\sout{Because the May 2026 document remains an exposure draft, we use VEI and the proposed band as supplementary descriptive guidance rather than as a final adopted compliance standard.}}
{Again, since the May 2026 document remains an exposure draft, we use VEI and the proposed band as descriptive guidance rather than as a final adopted compliance standard. Moreover, note that the first-versus-last quantile comparison in~\eqref{eq:vei} does not necessarily reveal non-monotone patterns within the median ratios distribution. Hence, we also report the median valuation ratio and its 90\% confidence interval for every VEI percentile group, together with the overall median ratio.} 

% For smaller stratified samples, the number of groups follows the corresponding sample-size guidance in the exposure draft rather than imposing deciles mechanically \cite{IAAO2026ExposureRatio}.}

% \textcolor{blue}{\sout{The previous summary table, which omitted MAPE and MKI and summarized vertical equity only through PRD, PRB, and VEI, is replaced by the revised table below.}}

% \begin{table}[!ht]
% \centering
% \footnotesize
% \setlength{\tabcolsep}{5pt}
% \renewcommand{\arraystretch}{1.10}

% % Indentation for individual measures within each family
% \newcommand{\metriclabel}[1]{\hspace*{0.9em}#1}

% \caption{Predictive and assessor-facing evaluation measures.}
% \label{tab:assessment_metrics_summary}

% \begin{tabularx}{\textwidth}{
% @{}
% >{\raggedright\arraybackslash}p{3.8cm}
% >{\centering\arraybackslash}p{1.8cm}%{1.5cm}
% >{\centering\arraybackslash}p{2.9cm}
% >{\raggedright\arraybackslash}X
% @{}
% }
% \toprule
% \textbf{Measure}
% & \textbf{\shortstack{Ideal value}}
% & \oldtext{\textbf{\shortstack{Reference range}}}\newtext{\textbf{\shortstack{Reference /\\ guidance range}}}
% & \textbf{Interpretation} \\
% \midrule

% \textbf{Predictive performance}
% & & & \\
% \cmidrule(r){1-1}

% \metriclabel{$R^2_P$}
% & $1$
% & --
% & Original price-scale goodness of fit. \\

% \metriclabel{$\operatorname{MAE}_P$}
% & $0$
% & --
% & Mean absolute prediction error in dollars. \\

% \metriclabel{$\operatorname{MAPE}_P$}
% & $0$
% & --
% & Mean absolute prediction error relative to sale price. \\

% \metriclabel{$\operatorname{RMSE}_{\log P}$}
% & $0$
% & --
% & RMSE on the log-price estimation scale. \\


% \midrule
% \addlinespace[5pt]
% \textbf{{Valuation level}
% % \quad\textcolor{impGreen}{Valuation level}
% }
% & & & \\
% \cmidrule(r){1-1}

% \metriclabel{Median ratio $m_S$}
% & $1$
% & --
% & Robust aggregate valuation level. \\

% \metriclabel{Mean ratio $\bar r_S$}
% & $1$
% & --
% & Unweighted aggregate valuation level. \\

% \metriclabel{Weighted mean $\bar r_{W,S}$}
% & $1$
% & --
% & Value-weighted valuation level. \\


% \midrule
% \addlinespace[6pt]
% \textbf{{Horizontal uniformity}
% % \quad\textcolor{impGreen}{Ratio dispersion}
% }
% & & & \\
% \cmidrule(r){1-1}

% \metriclabel{COD}
% & --
% & $[5.0,15.0]$\textsuperscript{$\dagger$}
% & Absolute dispersion relative to the median ratio. \\

% \metriclabel{COV}
% & 0
% & --
% & Relative dispersion around the mean ratio. \\


% \midrule
% \addlinespace[5pt]
% \textbf{Vertical equity}
% & & & \\
% \cmidrule(r){1-1}

% \metriclabel{PRD}
% & $1$
% & $[0.98,1.03]$
% & Mean-to-weighted-mean ratio comparison. \\

% \metriclabel{PRB}
% & $0$
% & $[-0.05,0.05]$
% & Regression-based value-related ratio trend. \\

% \metriclabel{MKI}
% & $1$
% & --
% & Distributional concentration relative to sale prices. \\

% \metriclabel{VEI}
% & $0$
% & $[-10\%,10\%]$\textsuperscript{$\dagger$}
% & High-versus-low value-group median-ratio difference. \\

% \bottomrule
% \end{tabularx}

% \vspace{1mm}
% \begin{minipage}{\textwidth}
% \scriptsize
% \emph{Notes.}
% ``--'' denotes that no universal ideal or reference range is imposed.
% PRD and PRB reference ranges follow the 2013 IAAO Standard \cite{IAAO2013Ratio}.
% \textsuperscript{$\dagger$}The May 2026 IAAO exposure draft gives $5.0$--$15.0$
% as the general COD range for improved residential property and proposes
% $[-10\%,10\%]$ for VEI; applicable COD ranges vary by property and market profile
% \cite{IAAO2026ExposureRatio}. The May 2026 document remains an exposure draft
% under review; the 2013 Standard remains the recommended IAAO guidance pending
% approval of a final revision.
% \end{minipage}

% \end{table}


% \begin{table}[!ht]
% \centering
% \footnotesize
% \setlength{\tabcolsep}{5pt}
% \renewcommand{\arraystretch}{1.10}

% % Indentation for individual measures within each family
% \newcommand{\metriclabel}[1]{\hspace*{0.9em}#1}

% \caption{Predictive and assessor-facing evaluation measures.}
% \label{tab:assessment_metrics_summary}

% \begin{tabularx}{\textwidth}{
% @{}
% >{\raggedright\arraybackslash}p{3.8cm}
% >{\centering\arraybackslash}p{1.8cm}%{1.5cm}
% >{\centering\arraybackslash}p{2.9cm}
% >{\raggedright\arraybackslash}X
% @{}
% }
% \toprule
% \textbf{Measure}
% & \textbf{\shortstack{Ideal value}}
% & \textbf{\shortstack{Reference range}}
% & \textbf{Interpretation} \\
% \midrule

% \textbf{Predictive performance}
% & & & \\
% \cmidrule(r){1-1}

% \metriclabel{$R^2_P$}
% & $1$
% & --
% & Original price-scale goodness of fit. \\

% \metriclabel{$\operatorname{MAE}_P$}
% & $0$
% & --
% & Mean absolute prediction error in dollars. \\

% \metriclabel{$\operatorname{MAPE}_P$}
% & $0$
% & --
% & Mean absolute prediction error relative to sale price. \\

% \metriclabel{$\operatorname{RMSE}_{\log P}$}
% & $0$
% & --
% & RMSE on the log-price estimation scale. \\


% \midrule
% \addlinespace[5pt]
% \textbf{{Valuation level}
% % \quad\textcolor{impGreen}{Valuation level}
% }
% & & & \\
% \cmidrule(r){1-1}

% \metriclabel{Median ratio $m_S$}
% & $1$
% & $[0.90,1.10]$
% & Robust aggregate valuation level. \\

% \metriclabel{Mean ratio $\bar r_S$}
% & $1$
% & $[0.90,1.10]$
% & Unweighted aggregate valuation level. \\

% \metriclabel{Weighted mean $\bar r_{W,S}$}
% & $1$
% & $[0.90,1.10]$
% & Value-weighted valuation level. \\


% \midrule
% \addlinespace[6pt]
% \textbf{{Horizontal uniformity}
% % \quad\textcolor{impGreen}{Ratio dispersion}
% }
% & & & \\
% \cmidrule(r){1-1}

% \metriclabel{COD}
% & --
% & $[5.0,15.0]$
% & Absolute dispersion relative to the median ratio. \\

% \metriclabel{COV}
% & 0
% & --
% & Relative dispersion around the mean ratio. \\


% \midrule
% \addlinespace[5pt]
% \textbf{Vertical equity}
% & & & \\
% \cmidrule(r){1-1}

% \metriclabel{PRD}
% & $1$
% & $[0.98,1.03]$
% & Mean-to-weighted-mean ratio comparison. \\

% \metriclabel{PRB}
% & $0$
% & $[-0.05,0.05]$
% & Regression-based value-related ratio trend. \\

% \metriclabel{MKI}
% & $1$
% & $[0.95,1.05]$
% & Distributional concentration relative to sale prices. \\

% \metriclabel{VEI}
% & $0$
% & $[-10\%,10\%]$
% & High-versus-low value-group median-ratio difference. \\

% \bottomrule
% \end{tabularx}

% \vspace{1mm}
% \begin{minipage}{\textwidth}
% \scriptsize
% \emph{Notes.}
% ``--'' denotes that no universal ideal or reference range is imposed.
% The median-ratio, PRD, PRB, and COD ranges follow IAAO ratio-study and
% mass-appraisal guidance \cite{IAAO2013Ratio,IAAO2018AVM,IAAO2025Mass}.
% The MKI range $[0.95,1.05]$ is used by CCAO's ratio-study implementation
% and is also identified in recent IAAO exposure-draft guidance as an
% optimal range for lower-variability strata.
% The May 2026 IAAO exposure draft proposes $[-10\%,10\%]$ for VEI.
% The exposure draft remains under review; the 2013 Standard remains the
% recommended IAAO guidance until a final revision is approved.
% \end{minipage}

% \end{table}

\begin{table}[!ht]
\centering
\footnotesize
\setlength{\tabcolsep}{5pt}
\renewcommand{\arraystretch}{1.10}

% Indentation for individual measures within each family
\newcommand{\metriclabel}[1]{\hspace*{0.9em}#1}

\caption{Predictive and assessor-facing evaluation measures.}
\label{tab:assessment_metrics_summary}

\begin{tabularx}{\textwidth}{
@{}
>{\raggedright\arraybackslash}p{3.8cm}
>{\centering\arraybackslash}p{1.8cm}%{1.5cm}
>{\centering\arraybackslash}p{2.9cm}
>{\raggedright\arraybackslash}X
@{}
}
\toprule
\textbf{Measure}
& \textbf{\shortstack{Ideal value}}
& 
% \oldtext{\textbf{\shortstack{Reference range}}}
\newtext{\textbf{\shortstack{Reference range}}}
& \textbf{Interpretation} \\
\midrule

\textbf{Predictive performance}
& & & \\
\cmidrule(r){1-1}

\metriclabel{$R^2_P$}
& $1$
& --
& Original price-scale goodness of fit. \\

\metriclabel{$\operatorname{MAE}_P$}
& $0$
& --
& Mean absolute prediction error in dollars. \\

\metriclabel{$\operatorname{MAPE}_P$}
& $0$
& --
& Mean absolute prediction error relative to sale price. \\

\metriclabel{$\operatorname{RMSE}_{\log P}$}
& $0$
& --
& RMSE on the log-price estimation scale. \\


\midrule
\addlinespace[5pt]
\textbf{{Valuation level}
% \quad\textcolor{impGreen}{Valuation level}
}
& & & \\
\cmidrule(r){1-1}

\metriclabel{Median ratio $m_S$}
& $1$
& $[0.90,1.10]$
& Robust aggregate valuation level. \\

\metriclabel{Mean ratio $\bar r_S$}
& $1$
& $[0.90,1.10]$
& Unweighted aggregate valuation level. \\

\metriclabel{Weighted mean $\bar r_{W,S}$}
& $1$
& $[0.90,1.10]$
& Value-weighted valuation level. \\


\midrule
\addlinespace[6pt]
\textbf{{Horizontal uniformity}
% \quad\textcolor{impGreen}{Ratio dispersion}
}
& & & \\
\cmidrule(r){1-1}

\metriclabel{COD}
& --
& $[5.0,15.0]$
& Absolute dispersion relative to the median ratio. \\

\metriclabel{COV}
& --
& $[6.25\%,18.75\%]$
& Relative dispersion around the mean ratio. \\


\midrule
\addlinespace[5pt]
\textbf{Vertical equity}
& & & \\
\cmidrule(r){1-1}

\metriclabel{PRD}
& $1$
& $[0.98,1.03]$
& Mean-to-weighted-mean ratio comparison. \\

\metriclabel{PRB}
& $0$
& $[-0.05,0.05]$
& Regression-based value-related ratio trend. \\

\metriclabel{MKI}
& $1$
& $[0.95,1.05]$
& Distributional concentration relative to sale prices. \\

\metriclabel{VEI}
& $0$
& $[-10\%,10\%]$
& High-versus-low value-group median-ratio difference. \\

\bottomrule
\end{tabularx}

\vspace{1mm}
\begin{minipage}{\textwidth}
\scriptsize
\emph{Notes.}
% \oldtext{``--'' denotes that no universal ideal or reference range is imposed.}
\newtext{``--'' denotes that no universal ideal or guidance range is imposed. The ranges below combine adopted IAAO guidance, explicitly identified derived approximations, and supplementary exposure-draft or implementation guidance as detailed next.}
The median-ratio, mean-ratio, weighted-mean-ratio, PRD, PRB, and COD ranges
follow IAAO ratio-study and mass-appraisal guidance
\cite{IAAO2013Ratio,IAAO2018AVM,IAAO2025Mass}.
For COV, the approximate range $[6.25\%,18.75\%]$ is derived from the
IAAO relationship $\operatorname{COV}\approx1.25\,\operatorname{COD}$ under
approximately normal ratio distributions and the residential-improved COD
reference range $[5,15]$.
The MKI range $[0.95,1.05]$ is used by CCAO's ratio-study implementation
and is also identified in recent IAAO exposure-draft guidance as an
optimal range for lower-variability strata.
The May 2026 IAAO exposure draft proposes $[-10\%,10\%]$ for VEI.
The exposure draft remains under review; the 2013 Standard remains the
recommended IAAO guidance until a final revision is approved.
\end{minipage}

\end{table}
% \textcolor{blue}{\sout{Table~\ref{tab:assessment_metrics_summary} summarizes the principal evaluation criteria used throughout the paper. Predictive metrics quantify valuation error, assessment-level statistics summarize aggregate calibration, COD and COV summarize ratio dispersion, and PRD, PRB, and VEI diagnose price-related patterns. These dimensions should be interpreted jointly: strong predictive performance or low aggregate dispersion does not by itself imply neutral vertical-equity behavior.}}

% \textcolor{blue}{\sout{Table~\ref{tab:assessment_metrics_summary} summarizes the principal predictive and assessor-facing criteria used throughout the paper. Predictive metrics quantify different aspects of valuation error, assessment-level statistics summarize aggregate calibration, COD and COV summarize ratio dispersion, and VEI, PRB, MKI, and PRD provide complementary grouped, regression-based, distributional, and ratio-based views of vertical equity. Strong predictive performance or low aggregate dispersion does not by itself imply neutral vertical-equity behavior.}}

{Table~\ref{tab:assessment_metrics_summary} summarizes the principal predictive and assessor-facing criteria used throughout the paper. Predictive metrics quantify different aspects of valuation error, valuation-level statistics summarize aggregate calibration, COD and COV summarize ratio dispersion, and VEI, PRB, MKI, and PRD provide complementary grouped, regression-based, distributional, and ratio-based views of vertical equity. Strong predictive performance or low aggregate dispersion does not by itself imply neutral vertical-equity behavior.}
\newtext{The mechanism and residual-structure diagnostics below are deliberately excluded from the reference-range table because they are not assessor-standard compliance measures.}


\paragraph{
% \oldtext{Supplemental diagnostics.}
\newtext{Mechanism and residual-structure diagnostics.}}
{\color{impGreen}
The assessor-facing measures above evaluate valuation outcomes. The following three
diagnostics instead examine the residual--price structure most directly relevant to
the proposed training mechanism. They form an increasingly general sequence:
$\beta_{\log}$ measures the targeted first-order trend, $\Delta_{\mathrm{NL}}$
measures non-affine conditional-mean structure beyond that trend, and distance
correlation measures broader statistical dependence. None is treated as an
assessor-standard compliance statistic.

\emph{First-order mechanism.}
We first report the slope from regressing log valuation ratios on log sale prices:
\begin{equation}
\label{eq:log_ratio_slope}
\beta_{\log}(S)
=
\frac{
\widehat{\operatorname{Cov}}_S(e,y)
}{
\widehat{\operatorname{Var}}_S(y)
}
=
\frac{
\widehat{\operatorname{Cov}}_S(\log r,\log P)
}{
\widehat{\operatorname{Var}}_S(\log P)
}.
\end{equation}
A value of zero indicates no first-order log-ratio--log-sale-price trend. Negative
values indicate the regressive direction and positive values the progressive
direction. For a fixed evaluation sample, $\widehat{\operatorname{Var}}_S(y)$ is
common across models, so $\beta_{\log}(S)$ is the signed first-order slope
corresponding directly to the covariance targeted during training. Because sale
price is an imperfect proxy for latent market value and also appears in the
denominator of $r_i$, we use $\beta_{\log}$ as a mechanism diagnostic rather than
as a standalone definition of vertical equity.

\emph{Nonlinear conditional-mean structure.}
Let
\[
Z=\frac{y-\E[y]}{\sqrt{\Var(y)}}.
\]
Following the nonparametric coefficient-of-determination literature, define Pearson's
correlation ratio
\[
\eta^2(e\mid Z)
=
\frac{\Var\!\left(\E[e\mid Z]\right)}{\Var(e)},
\]
which measures the fraction of log-error variance explained by the unrestricted
conditional mean of $e$ given $Z$ \cite{DoksumSamarov1995}. The squared Pearson
correlation $\Corr(e,Z)^2$ is the corresponding fraction explained by the best
affine relationship. We therefore define the \emph{correlation-ratio nonlinearity
gap}
\begin{equation}
\label{eq:nonlinearity_gap}
\Delta_{\mathrm{NL}}
:=
\eta^2(e\mid Z)-\Corr(e,Z)^2.
\end{equation}
The descriptive name and the symbol $\Delta_{\mathrm{NL}}$ are paper-specific;
the underlying correlation-ratio construction is established in the statistical
literature. The gap is nonnegative at the population level and equals zero when the
conditional mean of the log valuation error is affine in log sale price. Positive
values therefore measure conditional-mean explanatory power beyond the best affine
trend. This distinction is especially relevant along a regularization path:
$\lvert\beta_{\log}\rvert$ may continue to fall even if stronger regularization
induces an S-shaped, U-shaped, or otherwise non-affine residual--price pattern.
Appendix~\ref{app:metrics} gives the equivalent projection representation, relates
the construction to parametric-versus-nonparametric regression-fit comparisons
\cite{HardleMammen1993}, and specifies the fixed cross-fitted estimator used in the
held-out and 2025 evaluations.
\latesttext{The same fixed estimator was later applied, fold by fold, to the already-frozen chronological-validation predictions as a retrospective diagnostic; those reconstructed CV values do not define the frozen prediction/COD transition span, the $\rho$ grid, or a model-selection rule.}

\emph{Broader dependence.}
Finally, we report distance correlation,
\begin{equation}
\label{eq:dcor_diagnostic}
\operatorname{dCor}(e,y)
=
\operatorname{dCor}(\log r,\log P),
\end{equation}
as an omnibus diagnostic of dependence between log valuation errors and log sale
prices \cite{SzekelyRizzoBakirov2007}. At the population level and under the usual
moment conditions, distance correlation is zero if and only if the two variables
are independent. It can therefore respond to nonlinear mean relationships,
price-related dispersion, and other forms of dependence that are intentionally
outside the narrower scope of $\beta_{\log}$ and $\Delta_{\mathrm{NL}}$. Distance
correlation does not identify a regressive or progressive direction, and
distributional independence is stronger than vertical neutrality as conventionally
evaluated in ratio studies. We use it as a broader residual-structure diagnostic,
not as an assessor-facing measure of vertical equity.
}

% \paragraph{\textcolor{impGreen}{Supplemental diagnostics.}} \todo{Note: I think we need this or an equivalent metric to also account for the non-linear distributional independence. We already talked about this in the past, but I am making it explicit here. IAAO metrics are not enough IMO. GOing more granular to locaton or other sensitive attributes buys a lot of credibility too.
% We can sill move it to Appendix if needed...}
% \textcolor{impGreen}{The assessor-facing statistics are complemented by three checks designed around failure modes of the proposed first-order correction. First, we examine the full VEI percentile-group profile described above rather than only its first-versus-last summary. Second, we report the \emph{distance correlation}.}
% \begin{equation}
% \label{eq:dcor_diagnostic}
% \textcolor{impGreen}{
% \operatorname{dCor}\bigl(e,y\bigr)
% =
% \operatorname{dCor}\bigl(\log r,\log P\bigr),
% }
% \end{equation}
% \textcolor{impGreen}{as an omnibus diagnostic of residual value dependence \cite{SzekelyRizzoBakirov2007}. Under standard moment conditions, population distance correlation equals zero if and only if the variables are distributionally independent ($r \perp P$). Note that independence is stronger than vertical neutrality as conventionally evaluated in ratio studies. For example, distance correlation can respond to price-related changes in residual dispersion even when group medians are flat. Hence, we use $\operatorname{dCor}$ only to test whether a (linear) covariance-based correction induces nonlinear dependencies, not as an assessor compliance metric. Third, where sample size and value variation are sufficient, we recompute the principal vertical-equity diagnostics previously defined within prespecified geographic and property strata rather than relying exclusively on countywide averages \cite{IAAO2026ExposureRatio}.}

% \textcolor{blue}{\sout{Two qualifications are important for interpretation. First, throughout the empirical model comparisons we use these point estimates as descriptive diagnostics rather than as formal determinations of compliance with IAAO standards; formal compliance testing can require additional uncertainty analysis. Second, countywide statistics can conceal geographic, property-class, temporal, or nonlinear heterogeneity. We therefore evaluate models using several complementary statistics and ratio plots rather than drawing conclusions from any single measure.}}

% \textcolor{blue}{\sout{Three benchmark principles govern the empirical use of these measures. First, point estimates are treated as descriptive model diagnostics rather than as formal determinations of compliance with IAAO standards: uncertainty is reported where it is material to substantive comparisons, including the confidence intervals accompanying the VEI percentile-group profile. Second, countywide statistics are complemented by prespecified geographic and property-stratum analyses when the available sample and value variation support meaningful comparison. Third, during model development the evaluation statistics are computed on chronological validation folds before the aggregation used for model selection. The full evaluation suite is not itself imposed as a collection of constraints over the model, but as a general diagnostic of its estimates.}}

{Three benchmark principles govern the empirical use of these measures. First, point estimates are treated as descriptive model diagnostics rather than as formal determinations of compliance with IAAO standards; uncertainty is reported where it is material to substantive comparisons, including the confidence intervals accompanying the VEI percentile-group profile \cite{IAAO2013Ratio,IAAO2026ExposureRatio}. Second, countywide statistics can conceal geographic and property-stratum heterogeneity; such subgroup analyses remain necessary complements for later operational validation when sample size and value variation support meaningful comparison. Third, during model development, evaluation statistics are computed on chronological validation folds before fold-level results are aggregated to summarize each regularization path. The full suite is used for evaluation rather than imposed jointly as a collection of model constraints.}





\subsection{Baseline Comparison and Empirical Motivation}
\label{subsec:baseline_tension}

% \textcolor{blue}{\sout{In this subsection we compare the historically motivated linear regression benchmark with the current model class used by the CCAO: unpenalized LightGBM.}}
{We compare the historically motivated linear-regression benchmark with unpenalized LightGBM, representing CCAO's current model class.} Both models are estimated and evaluated under the same prediction target,
underlying feature information, sample-construction rules, temporal-validation
framework, and evaluation procedure. The comparison is therefore a controlled
comparison of model classes rather than a reconstruction of official
assessments produced under different historical workflows.
Section~\ref{subsec:ccao_design} provides the complete empirical design.

{The primary comparison below uses the held-out evaluation sample defined in Section~\ref{subsec:ccao_design}. We also report a separate 2025 forward evaluation after each baseline specification is refit using all eligible 2016--2024 sales; the 2025 sample is not used for model selection.}

% Table~\ref{tab:ccao_baseline_results} and Figure~\ref{fig:baseline_motivation} regenerated from output/paper_v6_preselection baseline-report predictions under the corrected metric code and the frozen Section-2 LightGBM vector.

\oldtext{The previous baseline presentation combined prediction, valuation level, horizontal uniformity, vertical equity, and mechanism/residual-structure diagnostics in a single table.}
\oldtext{For readability and to keep the empirical hierarchy explicit, the baseline comparison is split into a primary table for prediction and assessor-facing vertical equity and a complementary table for valuation level, horizontal uniformity, and mechanism/residual structure.}
\latesttext{For readability and to keep the empirical hierarchy explicit, the main text retains the primary baseline table for prediction and assessor-facing vertical equity, while the complementary valuation-level, horizontal-uniformity, and mechanism/residual-structure baseline diagnostics are reported in Appendix~\ref{app:ccao_path_details}.}

{\color{latestPurple}
\begin{table}[!ht]
\centering
\scriptsize
\setlength{\tabcolsep}{2.4pt}
\renewcommand{\arraystretch}{1.06}
\newcommand{\baselineprimarymetric}[1]{\hspace*{0.6em}#1}
\caption{Primary baseline comparison: prediction and vertical equity on the held-out and 2025 forward evaluations.}
\label{tab:ccao_baseline_results}
\begin{tabularx}{\textwidth}{@{} >{\raggedright\arraybackslash}p{2.80cm} >{\centering\arraybackslash}X >{\centering\arraybackslash}X >{\centering\arraybackslash}X >{\centering\arraybackslash}X @{}}
\toprule
\textbf{Measure} & \multicolumn{2}{c}{\textbf{Held-out evaluation}} & \multicolumn{2}{c}{\textbf{2025 forward evaluation}} \\
\cmidrule(lr){2-3}\cmidrule(l){4-5}
& \textbf{Linear} & \textbf{\shortstack{Light\\GBM}} & \textbf{Linear} & \textbf{\shortstack{Light\\GBM}} \\
\midrule
\multicolumn{5}{@{}l}{\textbf{Prediction}} \\
\cmidrule(r){1-1}
\baselineprimarymetric{$R^2_P$} & 0.799 & \textbf{0.894} & 0.799 & \textbf{0.904} \\
\baselineprimarymetric{$\operatorname{MAE}_P$} & \$90,092 & \textbf{\$75,655} & \$99,371 & \textbf{\$78,484} \\
\baselineprimarymetric{$\operatorname{MAPE}_P$} & 24.1\% & \textbf{21.2\%} & 24.9\% & \textbf{20.8\%} \\
\baselineprimarymetric{$\operatorname{RMSE}_{\log P}$} & 0.322 & \textbf{0.289} & 0.313 & \textbf{0.278} \\
\addlinespace[5pt]
\multicolumn{5}{@{}l}{\textbf{Vertical equity}} \\
\cmidrule(r){1-1}
\baselineprimarymetric{PRD} & \textbf{1.042} & 1.069 & \textbf{1.052} & 1.079 \\
\baselineprimarymetric{PRB} & \textbf{-0.016} & -0.091 & \textbf{-0.029} & -0.106 \\
\baselineprimarymetric{MKI} & \textbf{0.975} & 0.923 & \textbf{0.954} & 0.907 \\
\baselineprimarymetric{VEI} & \textbf{-11.6\%} & -26.5\% & \textbf{-17.3\%} & -28.6\% \\
\bottomrule
\end{tabularx}
\vspace{1mm}
\begin{minipage}{\textwidth}
\scriptsize
\emph{Notes.}
Preferred values and applicable reference ranges are defined in Table~\ref{tab:assessment_metrics_summary}.
MAPE and VEI are reported in percent; MAE is reported in dollars.
For the held-out evaluation, both models are fit on the 344,607-sale development pool.
For the 2025 forward evaluation, each baseline specification is refit on all 382,897 eligible 2016--2024 sales.
Boldface follows the preferred direction of each measure and does not denote statistical significance or formal compliance.
\end{minipage}
\end{table}

}

\begin{oldrevisionblock}
\noindent\textbf{Previous main-text complementary baseline table.}\par\smallskip
\centering
\scriptsize
\setlength{\tabcolsep}{2.4pt}
\renewcommand{\arraystretch}{1.06}
\newcommand{\baselinecompmetric}[1]{\hspace*{0.6em}#1}
\begin{tabularx}{\textwidth}{@{} >{\raggedright\arraybackslash}p{2.80cm} >{\centering\arraybackslash}X >{\centering\arraybackslash}X >{\centering\arraybackslash}X >{\centering\arraybackslash}X @{}}
\toprule
\textbf{Measure} & \multicolumn{2}{c}{\textbf{Held-out evaluation}} & \multicolumn{2}{c}{\textbf{2025 forward evaluation}} \\
\cmidrule(lr){2-3}\cmidrule(l){4-5}
& \textbf{Linear} & \textbf{\shortstack{Light\\GBM}} & \textbf{Linear} & \textbf{\shortstack{Light\\GBM}} \\
\midrule
\multicolumn{5}{@{}l}{\textbf{Valuation level}} \\
\cmidrule(r){1-1}
\baselinecompmetric{Median ratio $m_S$} & \textbf{0.969} & 0.929 & \textbf{1.020} & 0.950 \\
\baselinecompmetric{Mean ratio $\bar r_S$} & 1.020 & \textbf{0.989} & 1.081 & \textbf{1.015} \\
\baselinecompmetric{Weighted mean $\bar r_{W,S}$} & \textbf{0.979} & 0.924 & \textbf{1.027} & 0.941 \\
\addlinespace[5pt]
\multicolumn{5}{@{}l}{\textbf{Horizontal uniformity}} \\
\cmidrule(r){1-1}
\baselinecompmetric{COD} & 24.7\% & \textbf{21.6\%} & 24.3\% & \textbf{21.3\%} \\
\baselinecompmetric{COV} & 45.2\% & \textbf{39.7\%} & 42.1\% & \textbf{37.2\%} \\
\addlinespace[5pt]
\multicolumn{5}{@{}l}{\textbf{Mechanism and residual structure}} \\
\cmidrule(r){1-1}
\baselinecompmetric{$\beta_{\log}$} & \textbf{-0.092} & -0.150 & \textbf{-0.109} & -0.164 \\
\baselinecompmetric{$\Delta_{\mathrm{NL}}$} & 0.131 & \textbf{0.119} & 0.122 & \textbf{0.121} \\
\baselinecompmetric{$\operatorname{dCor}(e,y)$} & \textbf{0.250} & 0.387 & \textbf{0.269} & 0.422 \\
\bottomrule
\end{tabularx}
\vspace{1mm}
\begin{minipage}{\textwidth}
\scriptsize
\emph{Notes.}
Median, mean, weighted-mean ratio, COD, and COV are assessor-facing valuation-level or horizontal-uniformity diagnostics defined in Section~\ref{subsec:metrics}.
The first-order mechanism slope $\beta_{\log}$, correlation-ratio nonlinearity gap $\Delta_{\mathrm{NL}}$, and distance correlation are defined in Eqs.~\eqref{eq:log_ratio_slope}, \eqref{eq:nonlinearity_gap}, and~\eqref{eq:dcor_diagnostic}.
For the three mechanism/residual-structure diagnostics, boldface follows the direction toward less first-order, non-affine conditional-mean, and broader residual--price structure, respectively; it should not be read as an assessor-standard equity ranking.
Boldface does not denote statistical significance or formal compliance.
\end{minipage}
\end{oldrevisionblock}

\latesttext{The corresponding complementary baseline diagnostics are retained in Appendix Table~\ref{tab:ccao_baseline_complementary}.}




\begin{figure}[!htbp]
\centering
\safeincludegraphics[width=0.8\textwidth]{img/generated_v12_994/baseline_models_motivation_2024_2025.pdf}
\caption{Descriptive valuation-ratio patterns against sale price for the two baseline models in the primary held-out evaluation and the 2025 forward evaluation. Curves show median ratios in 30 equal-count sale-price bins; shaded regions show the corresponding interquartile ranges.}
\label{fig:baseline_motivation}
\end{figure}


The corresponding VEI percentile-group profile, with deterministic 90\% bootstrap intervals, is reported in Appendix~\ref{app:ccao_path_details} (Figure~\ref{fig:vei_group_profile_placeholder}).

% \textcolor{blue}{\sout{Table~\ref{tab:ccao_baseline_results} and Figure~\ref{fig:baseline_value_groups} apply the evaluation framework in Section~\ref{subsec:metrics} to the two baseline model classes. Assessment level and vertical equity are interpreted separately: the former concerns aggregate calibration of the ratios, whereas the latter concerns how those ratios vary across property-value levels.}}

\oldtext{Table~\ref{tab:ccao_baseline_results} and Figure~\ref{fig:baseline_motivation} apply the evaluation framework in Section~\ref{subsec:metrics} to the two baseline model classes. Valuation level and vertical equity are interpreted separately: the former concerns aggregate calibration of the ratios, whereas the latter concerns how ratios vary across property-value levels.}\oldtext{Tables~\ref{tab:ccao_baseline_results} and~\ref{tab:ccao_baseline_complementary}, together with Figure~\ref{fig:baseline_motivation}, apply the evaluation framework in Section~\ref{subsec:metrics} to the two baseline model classes. The primary table isolates prediction and vertical equity, while the complementary table separates valuation level, horizontal uniformity, and mechanism/residual structure.}\latesttext{Table~\ref{tab:ccao_baseline_results} and Figure~\ref{fig:baseline_motivation} present the primary baseline comparison in the main text. Appendix Table~\ref{tab:ccao_baseline_complementary} reports the complementary valuation-level, horizontal-uniformity, and mechanism/residual-structure diagnostics.} 

{On the primary held-out evaluation, LightGBM improves all four predictive measures relative to linear regression: $R^2_P$ rises from $0.799$ to $0.894$, while MAE, MAPE, and log-scale RMSE decrease. PRD moves from $1.042$ to $1.069$, and VEI from $-11.6\%$ to $-26.5\%$. The 2025 forward evaluation shows the same qualitative ordering.}

% \textcolor{blue}{\sout{The comparison is application-specific and should not be interpreted as evidence that machine-learning or nonlinear valuation models are intrinsically more regressive than linear models. Its purpose is to determine whether the predictive advantages of the more flexible baseline are accompanied by stronger departures from vertical neutrality under the common evaluation framework. If the more flexible baseline provides meaningful predictive or ratio-dispersion gains while exhibiting a stronger value-related vertical-equity pattern, the model-design question is not simply whether to return to the linear benchmark. It is whether those predictive advantages can be retained while reducing the value-related pattern through the training objective.}}

{The comparison is application-specific and should not be interpreted as evidence that machine-learning or nonlinear valuation models are intrinsically more regressive than linear models. In this controlled comparison, however, the more flexible baseline provides substantial predictive gains and lower aggregate ratio dispersion while exhibiting a stronger regressive vertical-equity pattern. The model-design question is therefore not simply whether to return to the linear benchmark, but whether the predictive advantages of LightGBM can be retained while reducing that pattern through the training objective.}

% \textcolor{blue}{\sout{The assessor-facing diagnostics above define how the resulting valuations are evaluated; they need not be optimized directly during training. A first-order value-related pattern can therefore be represented on the training scale through covariance, which is used as a tractable training proxy rather than as a replacement for the assessor-facing vertical-equity measures.}}

{The assessor-facing diagnostics above determine how the resulting valuations are evaluated; they need not be optimized directly during training. The first-order mechanism slope in Eq.~\eqref{eq:log_ratio_slope} isolates the signed log-ratio--log-sale-price trend:
\[
\beta_{\log}(S)
=
\frac{\widehat{\operatorname{Cov}}_S(e,y)}
{\widehat{\operatorname{Var}}_S(y)}.
\]
Within a fixed training sample, $\widehat{\operatorname{Var}}(y)$ does not depend on the fitted model. Penalizing squared covariance therefore controls the same first-order slope up to a fixed scaling. Using $e_i(f)=\log r_i(f)$ and $y_i=\log P_i$,
\begin{equation}
\label{eq:training_covariance_bridge}
\widehat{\operatorname{Cov}}\bigl(e(f),y\bigr)
=
\widehat{\operatorname{Cov}}\bigl(\log r(f),\log P\bigr).
\end{equation}
We use covariance as the trainable mechanism, while the assessor-facing measures and the residual-structure diagnostics remain separate out-of-sample checks on what the correction does and does not remove. Section~\ref{sec:method} develops the correction.}



\section{Covariance-Guided Regressivity Correction}
\label{sec:method}

% \textcolor{blue}{\sout{This section present translates the vertical-equity concern identified in Section~\ref{sec:setting} into a trainable objective. The construction proceeds in five steps. We first express assessment ratios as residuals on the log-price scale and use covariance to measure their global linear relationship with property values. We then introduce a direct squared-covariance penalty and derive a sample-additive, covariance-guided surrogate. Next, we instantiate both formulations for linear regression and second-order tree boosting. Finally, we clarify what the penalties can and cannot control and describe how the penalty level is selected using temporal validation and assessor-facing diagnostics.}}

{This section translates the value-related regressivity identified in Section~\ref{sec:setting} into a training-time correction penalty. We first define a global first-order covariance-based regressivity target on the log-price scale, then introduce a direct squared-covariance penalty, and finally derive a sample-additive upper-bound surrogate that is compatible with standard boosting software. The correction penalty leaves the prediction target, feature information, and base model class unchanged. That is, for a fixed penalty formulation, the only additional tuning parameter is the penalty strength, which we denote $\rho$.}

{Let $J_0(f)$ denote the baseline training objective, whose data-fit component is the squared log-price loss $L(f)$ defined in Section~\ref{subsec:baseline_models}. All quantities below are computed from the ``current'' training sample. In temporal validation, the centering quantities and penalties are recomputed using the training block of each fold only.}


\subsection{{Log-Scale Regressivity Target}}
\label{subsec:training_regressivity}

Let
\begin{equation}
\label{eq:centered_log_price}
\bar y=\frac{1}{n}\sum_{i=1}^n y_i,
\qquad
c_i=y_i-\bar y,
\qquad
\bar e(f)=\frac{1}{n}\sum_{i=1}^n e_i(f).
\end{equation}
The empirical covariance between log-residuals and log-prices is
\begin{equation}
\label{eq:training_covariance}
C(f)
=
\widehat{\operatorname{Cov}}\bigl(e(f),y\bigr)
=
\frac{1}{n}\sum_{i=1}^n
\bigl(e_i(f)-\bar e(f)\bigr)(y_i-\bar y)
=
\frac{1}{n}\sum_{i=1}^n e_i(f)c_i.
\end{equation}
The last equality follows from $\sum_{i=1}^n c_i=0$, so centering the residuals does not change the covariance.

This covariance has a direct slope interpretation. Let
\[
V_y=\frac{1}{n}\sum_{i=1}^n c_i^2.
\]
Provided $V_y>0$, the least-squares slope from regressing log valuation ratios on log prices is
\begin{equation}
\label{eq:covariance_slope}
\operatorname{slope}\bigl(e(f)\sim y\bigr)
=
\frac{C(f)}{V_y}.
\end{equation}
Because $V_y$ is fixed by the training sample, reducing $|C(f)|$ is equivalent to reducing the magnitude of this global linear trend. In particular, $C(f)<0$ indicates that log valuation ratios tend to decline with log sale price, which is the first-order pattern associated with regressivity.

\newtext{A qualification is essential for interpreting this training quantity. Sale price is both the prediction target and the observed value benchmark, so residual--outcome association can be generated partly by the geometry of squared-error fitting itself. Appendix~\ref{app:theory} makes this explicit: in an exact fixed linear prediction space containing an intercept, the ordinary least-squares projection satisfies $C_0=-\|f_0-y\|_n^2\leq0$. Negative in-sample covariance is therefore not, by itself, evidence of latent vertical inequity. We use it as a controllable first-order training mechanism and rely on chronological out-of-sample, proxy-aware assessor diagnostics to evaluate whether the resulting valuations actually improve.}

{This slope is a training-scale mechanism measure, not the assessor-facing PRB statistic defined in Section~\ref{subsec:metrics}. We use covariance rather than correlation because $V_y$ is fixed within the training sample, whereas correlation additionally normalizes by the model-dependent residual variance. Covariance therefore provides a direct training target for the global value-related slope.}

{Sale price enters this construction as the observed market-value benchmark for verified training sales; it is not introduced as a protected attribute. The broader connection to dependence-based fair regression is discussed in Section~\ref{subsec:related_work}.}


\subsection{{Direct Squared-Covariance Penalty}}
\label{subsec:direct_penalty}

% \textcolor{blue}{\sout{We therefore define the direct covariance-penalized learner as}}
{We first penalize the chosen first-order target directly. The \emph{covariance-penalized learner} is}

% \begin{center}
% % \textcolor{blue}{\sout{\mbox{$\displaystyle
% % f_\rho^{\mathrm{cov}}
% % \in
% % \argminop_{f\in\Hcal}
% \left\{
% L(f)+\Omega(f)+\frac{\rho}{2}C(f)^2
% \right\},
% \qquad
% \rho\ge 0.
% $}}}
% \end{center}

{
\begin{equation}
\label{eq:direct_penalty}
f_\rho^{\mathrm{cov}}
\in
\argminop_{f\in\Hcal}
\left\{
J_0(f)+\frac{\rho}{2}C(f)^2
\right\},
\qquad
\rho\geq 0.
\end{equation}
}

When $\rho=0$, formulation~\eqref{eq:direct_penalty} reduces to the baseline learner. As $\rho$ increases, the optimization places greater weight on reducing the global residual--price trend. Squaring the covariance treats negative and positive trends symmetrically. This is important because the objective is not to replace regressivity with progressivity, but to move the first-order relationship toward zero.

% \textcolor{blue}{\sout{The direction of the correction can be seen from the gradient. Let $e=(e_1,\ldots,e_n)^\top$ and $c=(c_1,\ldots,c_n)^\top$. Excluding $\Omega(f)$, the gradient and Hessian with respect to the prediction vector $\widehat y$ are}}

% \begin{center}
% \textcolor{blue}{\sout{\mbox{$\displaystyle
% \nabla_{\widehat y}
% \left(L(f)+\frac{\rho}{2}C(f)^2\right)
% =
% \frac{2}{n}e+\frac{\rho}{n}C(f)c
% $}}}

% \textcolor{blue}{\sout{\mbox{$\displaystyle
% \nabla_{\widehat y}^2
% \left(L(f)+\frac{\rho}{2}C(f)^2\right)
% =
% \frac{2}{n}I+\frac{\rho}{n^2}cc^\top.
% $}}}
% \end{center}

{For interpretation, it is enough to consider the contribution of the covariance penalty to the gradient of observation $i$:}
{
\begin{equation}
\label{eq:direct_penalty_gradient}
\frac{\partial}{\partial \widehat y_i}
\left[
\frac{\rho}{2}C(f)^2
\right]
=
\frac{\rho}{n}C(f)c_i.
\end{equation}
}

% \textcolor{blue}{\sout{Suppose, for example, that the fitted model is regressive, so $C(f)<0$. For a lower-value property, $c_i<0$, and the covariance component of the gradient is positive; a descent step therefore lowers its predicted value. For a higher-value property, $c_i>0$, and the covariance component is negative; a descent step raises its predicted value. The penalty consequently acts in the direction needed to flatten the global ratio--value relationship. The direction reverses automatically when the model is progressive.}}


% \textcolor{blue}{\sout{Suppose, for example, that the fitted model is regressive due to $C(f)<0$. For a sale below the geometric-mean price (lower-tier property), $c_i<0$, and thus the covariance component of the gradient is positive, i.e., a descent step lowers its predicted value. For a sale above the geometric-mean price (higher-tier property), $c_i>0$, and thus the covariance component is negative, i.e., a descent step raises its predicted value. The penalty therefore acts in the direction required to flatten the global log-ratio--log-price trend. Note that the direction and interpretation reverse automatically when the fitted trend is progressive.}}

{Suppose, for example, that the fitted model is regressive due to $C(f)<0$. For a sale below the geometric-mean price, $c_i<0$, and thus the covariance component of the gradient is positive, i.e., a descent step lowers its predicted value. For a sale above the geometric-mean price, $c_i>0$, and thus the covariance component is negative, i.e., a descent step raises its predicted value. The penalty therefore acts in the direction required to flatten the global log-ratio--log-price trend. Note that the direction and interpretation reverse automatically when the fitted trend is progressive.}


Note that the direct penalty targets how valuation ratios vary with sale price, but not their 
% \oldtext{overall assessment level}
\newtext{overall valuation level}. Multiplying every predicted value by the same positive constant shifts all log valuation ratios by the same amount but leaves the covariance $C(f)$ unchanged. Thus, even a model with $C(f)=0$ can have a median or mean valuation ratio different from one. Hence, valuation level must be evaluated separately.

% {The direct penalty controls this first-order slope separately from \oldtext{overall assessment level}\newtext{overall valuation level}. Adding the same constant to every log prediction leaves $C(f)$ unchanged, i.e., multiplying every price prediction by the same positive constant does not change the targeted covariance. Assessment-level calibration therefore remains a separate modeling and evaluation dimension.}

% \textcolor{blue}{\sout{The direct formulation most faithfully represents the statistical target, but it creates a computational complication. The rank-one term $cc^\top$ in~\eqref{eq:direct_gradient_hessian} couples all observations. Standard second-order boosting interfaces instead expect one gradient and one scalar Hessian for each observation. Section~\ref{subsec:model_specific_implementations} therefore develops both a diagonal Hessian approximation for standard boosting software and an exact leaf-weight solver for a fixed tree structure.}}

{For second-order models (e.g., LightGBM), the corresponding Hessian of the squared log-price loss plus direct penalty is}
{
\begin{equation}
\label{eq:direct_gradient_hessian}
\nabla_{\widehat y}^2
\left(
L(f)+\frac{\rho}{2}C(f)^2
\right)
=
\frac{2}{n}I+\frac{\rho}{n^2}cc^\top.
\end{equation}
}
The rank-one Hessian term $\frac{\rho}{n^{2}}cc^\top$ introduces cross-observation curvature. Standard LightGBM custom objectives provide one gradient and one Hessian value per observation, so this dense term cannot be represented exactly through the usual interface. An exact second-order implementation therefore requires additional machinery: using only the diagonal is possible, but discards the cross-observation curvature. The next subsection proposes an alternative penalty term to handle this caveat. Further, the model-specific derivations are given in Appendix~\ref{app:implementation}.

% . The exact second-order implementation therefore requires additional machinery: a diagonal approximation is possible but discards these cross-observation terms. The model-specific derivations are given in Appendix~\ref{app:implementation}.}

% \textcolor{blue}{\sout{The fixed-prediction-space analysis in Appendix~\ref{app:theory} provides additional intuition. Within a fixed linear prediction space, the magnitude of the covariance decreases smoothly with $\rho$, while the increase in in-sample squared error is second order near $\rho=0$. This explains why a modest penalty can, in principle, meaningfully reduce the targeted trend at a relatively small accuracy cost. For boosted trees, this result is only a local benchmark because changing $\rho$ can also change splits, leaf structures, and stopping behavior.}}


\subsection{{Sample-Additive Covariance Upper-Bound Surrogate}}
\label{subsec:surrogate}

% \textcolor{blue}{\sout{The direct penalty presents two practical limitations. First, covariance is an aggregate statistic, so positive and negative local relationships can offset one another. Second, its squared form is not additive across observations and therefore does not fit directly into the standard custom-objective interface used by LightGBM. We address the computational limitation by deriving a sample-additive surrogate motivated by an upper bound on the squared covariance.}}

{For standard second-order boosting software, we seek a penalty that decomposes exactly across observations. Given that~\eqref{eq:training_covariance} already gives $C(f)=\frac{1}{n}\sum_{i=1}^n e_i(f)c_i$, we can make use of Jensen's inequality that yields a sample-additive upper bound to the squared-covariance term directly.}

% \textcolor{blue}{\sout{By Jensen's inequality,}}
% \begin{center}
% \textcolor{blue}{\sout{\mbox{$\displaystyle
% C(f)^2
% =
% \left[
% \frac{1}{n}\sum_{i=1}^n
% \bigl(e_i(f)-\bar e(f)\bigr)c_i
% \right]^2
% \leq
% \frac{1}{n}\sum_{i=1}^n
% \bigl(e_i(f)-\bar e(f)\bigr)^2c_i^2.
% $}}}
% \end{center}

{
\begin{equation}
\label{eq:surrogate_upper_bound}
C(f)^2
=
\left(
\frac{1}{n}\sum_{i=1}^n e_i(f)c_i
\right)^2
\leq
\frac{1}{n}\sum_{i=1}^n e_i(f)^2c_i^2.
\end{equation}
}
% \textcolor{blue}{\sout{Although the right-hand side is a sum, it remains coupled through the model-dependent mean residual $\bar e(f)$. To obtain a fully observation-additive objective, we replace the centered residual $e_i(f)-\bar e(f)$ with $e_i(f)$ and define}}
{This inequality motivates the fully observation-additive surrogate penalty}
\begin{equation}
\label{eq:surrogate}
\Psi^{\mathrm{surr}}(f)
=
\frac{1}{n}\sum_{i=1}^n e_i(f)^2c_i^2
=
\frac{1}{n}\sum_{i=1}^n
e_i(f)^2(y_i-\bar y)^2.
\end{equation}
The resulting \emph{surrogate-penalized learner} is therefore
\begin{equation}
\label{eq:surrogate_learning}
f_\rho^{\mathrm{surr}}
\in
\argminop_{f\in\Hcal}
\left\{
{J_0(f)}
+
\rho\Psi^{\mathrm{surr}}(f)
\right\}.
\end{equation}
% \textcolor{blue}{\sout{After replacing the centered residual, $\Psi^{\mathrm{surr}}$ is not, in general, an upper bound on $C(f)^2$. It should therefore be interpreted as a \emph{covariance-guided} regularizer rather than exact covariance control. Its connection to the direct penalty is that both emphasize residuals according to their position in the value distribution. The surrogate makes this principle observation-specific and computationally separable.}}
% \textcolor{blue}{\sout{Thus, $\Psi^{\mathrm{surr}}(f)$ is a genuine upper bound on the squared empirical covariance in~\eqref{eq:training_covariance}. Nevertheless, the bound does not make surrogate minimization equivalent to direct covariance minimization: the upper bound can be loose when the signed contributions $e_i(f)c_i$ cancel in aggregate covariance. The surrogate instead penalizes the squared magnitude of each contribution before aggregation, so opposite-signed contributions cannot cancel within the penalty.}}
{
Thus, $\Psi^{\mathrm{surr}}(f)$ is a genuine upper bound on the squared empirical covariance in~\eqref{eq:training_covariance}. Its exact relationship with the direct-covariance penalty can be seen by making the observation
% \[
% q_i(f)=e_i(f)c_i.
% \]
% Since $C(f)=n^{-1}\sum_i q_i(f)$,
\begin{equation}
\label{eq:surrogate_decomposition}
\Psi^{\mathrm{surr}}(f)
=
C(f)^2
+
\frac{1}{n}
\sum_{i=1}^n
\left(e_i(f)c_i-C(f)\right)^2.
\end{equation}
% \oldtext{Hence, the surrogate controls squared covariance together with the dispersion of the observation-level residual--price contributions. It is therefore generally more conservative than direct covariance minimization: opposite-signed contributions can cancel in $C(f)$ but not in $\Psi^{\mathrm{surr}}(f)$.} 
\newtext{Hence, $\Psi^{\mathrm{surr}}(f)$ exceeds $C(f)^2$ by the nonnegative dispersion term in~\eqref{eq:surrogate_decomposition}: opposite-signed observation-level contributions can cancel in $C(f)$ but not in the surrogate. This upper-bound relation should not be interpreted as control of a broader class of nonlinear residual--price relationships; neither the bound nor its weighted-loss representation directly penalizes nonlinear conditional-mean lack-of-fit.}
}


Combining the baseline loss with the surrogate gives
\begin{equation}
\label{eq:surrogate_weighted_loss}
L(f)+\rho\Psi^{\mathrm{surr}}(f)
=
\frac{1}{n}\sum_{i=1}^n
\left[1+\rho(y_i-\bar y)^2\right]e_i(f)^2.
\end{equation}
Thus, the surrogate is equivalent to a weighted squared-error objective with fixed observation weights
\begin{equation}
\label{eq:surrogate_weights}
w_i(\rho)=1+\rho(y_i-\bar y)^2.
\end{equation}
Properties near the center of the log-price distribution receive weights close to the baseline weight, whereas properties in the lower and upper tails receive larger weights.

% \textcolor{blue}{\sout{Because $\bar y$ is the logarithm of the geometric mean sale price, the weighting is symmetric in proportional distance from the center of the price distribution. For example, equally sized multiplicative price deviations above and below this center receive the same additional weight.}}

% {Since $\exp(\bar y)$ is the geometric mean sale price, the weight function is symmetric in log-distance from this center: prices that differ from the geometric mean by reciprocal multiplicative factors receive the same additional weight. This symmetry concerns the weighting function only: the empirical lower and upper tails, and the resulting prediction changes, need not be symmetric.}

% \textcolor{blue}{\sout{Since $\exp(\bar y)$ is the geometric mean sale price, the weight function is symmetric in log-distance from this center: prices that differ from the geometric mean by reciprocal multiplicative factors receive the same additional weight. This symmetry concerns the weighting function only: the empirical lower and upper tails, and the resulting prediction changes, need not be symmetric.}}

{Since $\exp(\bar y)$ is the geometric mean sale price, the weight function is symmetric in log-distance from this center: prices that differ from the geometric mean by reciprocal multiplicative factors receive the same additional weight. This symmetry concerns the weighting function only: the empirical lower and upper tails, and the resulting prediction changes, need not be symmetric. Moreover, because the additional weight grows quadratically with log-price distance, observations far from the center can receive disproportionate influence. The log transformation compresses this leverage relative to an analogous raw-price construction, but does not bound it.}

% \textcolor{blue}{\sout{This representation clarifies the mechanism without imposing a directional assumption on individual observations. The surrogate raises the cost of residuals farther from the center of the log-price distribution. Because it uses $e_i(f)^2c_i^2$, however, it does not encode whether the aggregate residual--price trend is regressive or progressive. Its effect on covariance and on assessor-facing vertical-equity diagnostics must therefore be evaluated out of sample.}}

{This representation clarifies the mechanism without imposing a directional assumption on individual observations. The surrogate raises the cost of residuals farther from the center of the log-price distribution. Unlike the direct penalty, its observation-level correction is not determined by the sign of the aggregate covariance: it depends on the residual $e_i(f)$ and the squared log-price deviation $c_i^2$. Its resulting effect on the signed covariance and on assessor-facing vertical-equity diagnostics must therefore be evaluated out of sample.}

% \textcolor{blue}{\sout{This representation also clarifies why the surrogate can affect the ratio--value trend. Under a regressive baseline, the largest systematic residuals often occur toward the lower and upper ends of the value distribution. By increasing the cost of errors in both tails, the surrogate discourages the model from obtaining strong average predictive performance by allowing systematically positive residuals at the low end and negative residuals at the high end. Unlike covariance itself, the pointwise squared terms cannot cancel across observations.}}

% {This representation clarifies the mechanism without imposing a directional assumption on individual observations. The surrogate raises the cost of residuals farther from the center of the log-price distribution. Because it uses $e_i(f)^2c_i^2$, however, it does not encode whether the aggregate residual--price trend is regressive or progressive. Its effect on covariance and on assessor-facing vertical-equity diagnostics must therefore be evaluated out of sample.}

For squared loss, the observation-wise gradient and Hessian are
\begin{align}
\label{eq:surrogate_derivatives}
g_i
&=
\frac{2}{n}e_i(f)
\left[1+\rho(y_i-\bar y)^2\right],
\\
h_i
&=
\frac{2}{n}
\left[1+\rho(y_i-\bar y)^2\right].
\end{align}
% \oldtext{Both quantities depend only on observation $i$ once $\bar y$ has been computed. They can therefore be supplied through a standard second-order boosting custom-objective interface.}
\newtext{Both quantities depend only on observation $i$ once $\bar y$ has been computed. They can therefore be supplied through a standard second-order boosting custom-objective interface. Equation~\eqref{eq:surrogate_weighted_loss} also shows that the same mathematical objective can, in principle, be expressed as weighted squared loss in regression software that supports fixed observation weights. The CCAO paths reported in this paper use the custom-objective implementation below; an alternative native-weighted implementation should not be assumed to reproduce the identical finite boosting path without an explicit implementation-equivalence check.}

% \textcolor{blue}{\sout{The implementation uses the corresponding convention with the common normalization absorbed into the objective scaling.}}

{For the LightGBM implementation, we multiply the complete objective by the common positive factor $n/2$. This leaves the minimizer and the relative weighting indexed by $\rho$ unchanged while putting the $\rho=0$ derivatives on LightGBM's native squared-error scale. The supplied surrogate derivatives are therefore
\[
g_i=e_i(f)\left[1+\rho c_i^2\right],
\qquad
h_i=1+\rho c_i^2.
\]
The corresponding direct-objective implementation is given in Appendix~\ref{app:implementation}.}

% \textcolor{blue}{\sout{The direct and surrogate penalties consequently serve different methodological roles. The direct penalty is the clearest representation of the statistical quantity being controlled, but its curvature couples observations. The surrogate sacrifices exact covariance control in exchange for sample additivity, standard-software compatibility, and reduced sensitivity to cancellation across different parts of the value distribution. Moreover, the numerical values of $\rho$ under the two formulations are not directly comparable because the penalties have different scales. In the CCAO application, the surrogate provides the strongest combined validation performance and is therefore the principal practical specification.}}

{Unlike the direct covariance penalty, the surrogate is not invariant to a common shift in all log predictions because it penalizes the magnitude of $e_i(f)$ itself. It can therefore affect 
% \oldtext{overall assessment level}
\newtext{overall valuation level} in addition to the value-related trend. %Assessment-level diagnostics consequently remain part of validation and out-of-sample evaluation.
}


\subsection{{Interpretation, Validation, and Scope}}
\label{subsec:correction_scope}

% \textcolor{blue}{\sout{The two formulations serve complementary roles. The direct penalty targets the chosen global first-order association between log valuation ratios and log sale prices. The surrogate replaces direct covariance minimization with a sample-additive upper-bound objective that is directly compatible with standard boosting software. Because the two penalties have different scales, the numerical values of $\rho$ are not directly comparable across formulations. Their appropriate scale can also vary across training samples and jurisdictions, so the two regularization paths are traced separately through chronological validation.}}

{
% \oldtext{The two formulations intervene differently. The direct penalty responds to the aggregate signed first-order association between log valuation ratios and log sale prices. The surrogate instead reallocates squared-error importance across observations according to their distance from the center of the log-price distribution. Although this weighted objective upper-bounds squared covariance, it generally controls more than the targeted first-order trend, as shown in~\eqref{eq:surrogate_decomposition}.}
\newtext{The two formulations intervene differently. The Direct penalty responds to the aggregate signed first-order association between log valuation ratios and log sale prices. The Surrogate instead changes the data-fit objective itself by upweighting squared errors farther from the center of the log-price distribution. Although this observation-additive objective upper-bounds squared covariance, it is not merely an implementation approximation to the Direct penalty: its additional dispersion term in~\eqref{eq:surrogate_decomposition} gives it a distinct mechanism and can produce different path behavior.} Because the two penalties have different scales, the numerical values of $\rho$ are not directly comparable across formulations. Their appropriate scale can also vary across training samples and jurisdictions, so the two regularization paths are traced separately through chronological validation.}

% \textcolor{impGreen}{The two formulations serve complementary roles. The direct penalty targets the chosen global first-order association between log valuation ratios and log sale prices. The surrogate replaces direct covariance minimization with a sample-additive upper-bound objective that is directly compatible with standard boosting software. Because the two penalties have different scales, the numerical values of $\rho$ are not directly comparable across formulations. Their appropriate scale can also vary across training samples and jurisdictions, so the two regularization paths are traced separately through chronological validation.}

% \oldtext{Neither formulation directly optimizes PRD, PRB, VEI, COD, or assessment level. Nor does zero training covariance imply the absence of nonlinear, geographic, property-class, temporal, or other local value-related patterns. The training covariance is therefore treated as a mechanism target rather than as a sufficient measure of vertical equity. Predictive performance, valuation level, PRD, PRB, VEI, and ratio curves remain separate validation and out-of-sample evaluation criteria; subgroup diagnostics remain necessary complements for later operational validation.}

{\color{impGreen}Neither formulation directly optimizes PRD, PRB, VEI, COD, or valuation level. Nor does zero training covariance imply the absence of nonlinear, geographic, property-class, temporal, or other local value-related patterns. The training covariance is therefore treated as a mechanism target rather than as a sufficient measure of vertical equity. Predictive performance, valuation level, PRD, PRB, VEI, ratio curves, the correlation-ratio nonlinearity gap $\Delta_{\mathrm{NL}}$ in~\eqref{eq:nonlinearity_gap}, and the broader distance-correlation diagnostic remain separate validation and out-of-sample evaluation criteria; subgroup diagnostics remain necessary complements for later operational validation.}

{Sale price is required to construct the correction only for verified sales used during model fitting. Once a model has been trained, the resulting valuation function continues to generate predictions from $x$ alone, as in the baseline AVM; the correction does not require a current sale price for an unsold property.}

% \oldtext{Section~\ref{sec:empirical_design} specifies the chronological validation and regularization-path procedure used in the empirical application. Appendix~\ref{app:theory} provides complementary intuition for the direct penalty: within a fixed linear prediction space, increasing $\rho$ shrinks the targeted covariance smoothly while the increase in squared error is second order near $\rho=0$. This result is an interpretive benchmark, not a guarantee for fully retrained boosted trees, whose splits, leaves, and stopping behavior can change with the penalty.}
% \oldtext{Section~\ref{sec:empirical_design} specifies the chronological validation and regularization-path procedure used in the empirical application. Appendix~\ref{app:theory} gives a fixed-prediction-space benchmark for the Direct penalty: increasing $\rho$ shrinks the targeted covariance smoothly; when the space contains an intercept, the exact path is a one-parameter centered spread rescaling of the unpenalized fit; and the increase in squared error is second order near $\rho=0$. This rescaling interpretation motivates the simple post-hoc comparator in Section~\ref{subsec:comparators}. It remains an interpretive benchmark rather than a guarantee for fully retrained boosted trees, whose splits, leaves, and stopping behavior can change with the penalty.}
\newtext{Section~\ref{sec:empirical_design} specifies the chronological validation and regularization-path procedure used in the empirical application. Appendix~\ref{app:theory} gives parallel fixed-prediction-space benchmarks for the two penalties. For the Direct penalty, increasing $\rho$ shrinks the targeted covariance smoothly and, when the space contains an intercept, the exact path is a one-parameter centered spread rescaling of the unpenalized fit. For the Surrogate, the exact path is instead a log-price-distance-weighted projection of the outcome onto the same prediction space; locally, it moves in the learnable component of the baseline residuals weighted by squared log-price distance from the sample center. In both cases, the increase in baseline in-sample squared error is second order near $\rho=0$. These results clarify the rank-one geometry of the Direct correction and the generally multi-directional geometry of the Surrogate, while remaining interpretive benchmarks rather than guarantees for fully retrained boosted trees, whose splits, leaves, and stopping behavior can change with the objective.}




% \section{Covariance-Guided Regressivity Correction}
% \label{sec:method}

% This section present translates the vertical-equity concern identified in Section~\ref{sec:setting} into a trainable objective. The construction proceeds in five steps. We first express assessment ratios as residuals on the log-price scale and use covariance to measure their global linear relationship with property values. We then introduce a direct squared-covariance penalty and derive a sample-additive, covariance-guided surrogate. Next, we instantiate both formulations for linear regression and second-order tree boosting. Finally, we clarify what the penalties can and cannot control and describe how the penalty level is selected using temporal validation and assessor-facing diagnostics.


% \subsection{Direct Squared-Covariance Penalty}
% \label{subsec:direct_penalty}

% Let
% \begin{equation}
% \label{eq:centered_log_price}
% \bar y=\frac{1}{n}\sum_{i=1}^n y_i,
% \qquad
% c_i=y_i-\bar y,
% \qquad
% \bar e(f)=\frac{1}{n}\sum_{i=1}^n e_i(f).
% \end{equation}
% The empirical covariance between log residuals and log prices is
% \begin{align}
% \label{eq:training_covariance}
% C(f)=
% \widehat{\Cov}\bigl(e(f),y\bigr) =
% \frac{1}{n}\sum_{i=1}^n
% \bigl(e_i(f)-\bar e(f)\bigr)(y_i-\bar y) =
% \frac{1}{n}\sum_{i=1}^n e_i(f)c_i.
% \end{align}
% The last equality follows from \(\sum_i c_i=0\), so centering the residuals does not change the covariance.

% This covariance has a direct slope interpretation. Let
% \[
% V_y=\frac{1}{n}\sum_{i=1}^n c_i^2.
% \]
% Provided \(V_y>0\), the least-squares slope from regressing log valuation ratios on log prices is
% \begin{equation}
% \label{eq:covariance_slope}
% \operatorname{slope}\bigl(e(f)\sim y\bigr)
% =
% \frac{C(f)}{V_y}.
% \end{equation}
% Because \(V_y\) is fixed by the training sample, reducing \(|C(f)|\) is equivalent to reducing the magnitude of this global linear trend. In particular, \(C(f)<0\) indicates that log valuation ratios tend to decline with log price, which is the first-order pattern associated with regressivity.

% We therefore define the direct covariance-penalized learner as
% \begin{equation}
% \label{eq:direct_penalty}
% f_\rho^{\mathrm{cov}}
% \in
% \argminop_{f\in\Hcal}
% \left\{
% L(f)+\Omega(f)+\frac{\rho}{2}C(f)^2
% \right\},
% \qquad
% \rho\ge 0.
% \end{equation}
% When \(\rho=0\), formulation~\eqref{eq:direct_penalty} reduces to the baseline learner. As \(\rho\) increases, the optimization places greater weight on reducing the global residual--price trend. Squaring the covariance treats negative and positive trends symmetrically. This is important because the objective is not to replace regressivity with progressivity, but to move the first-order relationship toward zero.

% The direction of the correction can be seen from the gradient. Let \(e=(e_1,\ldots,e_n)^\top\) and \(c=(c_1,\ldots,c_n)^\top\). Excluding \(\Omega(f)\), the gradient and Hessian with respect to the prediction vector \(\widehat y\) are
% \begin{align}
% \label{eq:direct_gradient_hessian}
% \nabla_{\widehat y}
% \left(
% L(f)+\frac{\rho}{2}C(f)^2
% \right)
% &=
% \frac{2}{n}e+\frac{\rho}{n}C(f)c,
% \\
% \nabla_{\widehat y}^2
% \left(
% L(f)+\frac{\rho}{2}C(f)^2
% \right)
% &=
% \frac{2}{n}I+\frac{\rho}{n^2}cc^\top.
% \end{align}
% Suppose, for example, that the fitted model is regressive, so \(C(f)<0\). For a lower-value property, \(c_i<0\), and the covariance component of the gradient is positive; a descent step therefore lowers its predicted value. For a higher-value property, \(c_i>0\), and the covariance component is negative; a descent step raises its predicted value. The penalty consequently acts in the direction needed to flatten the global ratio--value relationship. The direction reverses automatically when the model is progressive.

% The direct formulation most faithfully represents the statistical target, but it creates a computational complication. The rank-one term \(cc^\top\) in~\eqref{eq:direct_gradient_hessian} couples all observations. Standard second-order boosting interfaces instead expect one gradient and one scalar Hessian for each observation. Section~\ref{subsec:model_specific_implementations} therefore develops both a diagonal Hessian approximation for standard boosting software and an exact leaf-weight solver for a fixed tree structure.

% The fixed-prediction-space analysis in Appendix~\ref{app:theory} provides additional intuition. Within a fixed linear prediction space, the magnitude of the covariance decreases smoothly with \(\rho\), while the increase in in-sample squared error is second order near \(\rho=0\). This explains why a modest penalty can, in principle, meaningfully reduce the targeted trend at a relatively small accuracy cost. For boosted trees, this result is only a local benchmark because changing \(\rho\) can also change splits, leaf structures, and stopping behavior.

% \subsection{Sample-Additive Covariance-Guided Surrogate}
% \label{subsec:surrogate}

% The direct penalty presents two practical limitations. First, covariance is an aggregate statistic, so positive and negative local relationships can offset one another. Second, its squared form is not additive across observations and therefore does not fit directly into the standard custom-objective interface used by LightGBM. We address the computational limitation by deriving a sample-additive surrogate motivated by an upper bound on the squared covariance.

% By Jensen's inequality,
% \begin{align}
% \label{eq:surrogate_upper_bound}
% C(f)^2
% &=
% \left[
% \frac{1}{n}\sum_{i=1}^n
% \bigl(e_i(f)-\bar e(f)\bigr)c_i
% \right]^2 \le
% \frac{1}{n}\sum_{i=1}^n
% \bigl(e_i(f)-\bar e(f)\bigr)^2c_i^2.
% \end{align}
% Although the right-hand side is a sum, it remains coupled through the model-dependent mean residual \(\bar e(f)\). To obtain a fully observation-additive objective, we replace the centered residual \(e_i(f)-\bar e(f)\) with \(e_i(f)\) and define
% \begin{equation}
% \label{eq:surrogate}
% \Psi^{\mathrm{surr}}(f)
% =
% \frac{1}{n}\sum_{i=1}^n e_i(f)^2c_i^2
% =
% \frac{1}{n}\sum_{i=1}^n
% e_i(f)^2(y_i-\bar y)^2.
% \end{equation}
% The resulting learner is
% \begin{equation}
% \label{eq:surrogate_learning}
% f_\rho^{\mathrm{surr}}
% \in
% \argminop_{f\in\Hcal}
% \left\{
% L(f)+\Omega(f)+\rho\Psi^{\mathrm{surr}}(f)
% \right\}.
% \end{equation}

% After replacing the centered residual, \(\Psi^{\mathrm{surr}}\) is not, in general, an upper bound on \(C(f)^2\). It should therefore be interpreted as a \emph{covariance-guided} regularizer rather than exact covariance control. Its connection to the direct penalty is that both emphasize residuals according to their position in the value distribution. The surrogate makes this principle observation-specific and computationally separable.

% Combining the baseline loss with the surrogate gives
% \begin{equation}
% \label{eq:surrogate_weighted_loss}
% L(f)+\rho\Psi^{\mathrm{surr}}(f)
% =
% \frac{1}{n}\sum_{i=1}^n
% \left[1+\rho(y_i-\bar y)^2\right]e_i(f)^2.
% \end{equation}
% Thus, the surrogate is equivalent to a weighted squared-error objective with fixed observation weights
% \[
% w_i(\rho)=1+\rho(y_i-\bar y)^2.
% \]
% Properties near the center of the log-price distribution receive weights close to the baseline weight, whereas properties in the lower and upper tails receive larger weights. Because \(\bar y\) is the logarithm of the geometric mean sale price, the weighting is symmetric in proportional distance from the center of the price distribution. For example, equally sized multiplicative price deviations above and below this center receive the same additional weight.

% This representation also clarifies why the surrogate can affect the ratio--value trend. Under a regressive baseline, the largest systematic residuals often occur toward the lower and upper ends of the value distribution. By increasing the cost of errors in both tails, the surrogate discourages the model from obtaining strong average predictive performance by allowing systematically positive residuals at the low end and negative residuals at the high end. Unlike covariance itself, the pointwise squared terms cannot cancel across observations.

% For squared loss, the observation-wise gradient and Hessian are
% \begin{align}
% \label{eq:surrogate_derivatives}
% g_i
% &=
% \frac{2}{n}e_i(f)
% \left[1+\rho(y_i-\bar y)^2\right],
% \\
% h_i
% &=
% \frac{2}{n}
% \left[1+\rho(y_i-\bar y)^2\right].
% \end{align}
% Both quantities depend only on observation \(i\) once \(\bar y\) has been computed. They can therefore be supplied directly through a standard second-order boosting custom-objective interface. The implementation uses the corresponding convention with the common normalization absorbed into the objective scaling.

% The direct and surrogate penalties consequently serve different methodological roles. The direct penalty is the clearest representation of the statistical quantity being controlled, but its curvature couples observations. The surrogate sacrifices exact covariance control in exchange for sample additivity, standard-software compatibility, and reduced sensitivity to cancellation across different parts of the value distribution. Moreover, the numerical values of \(\rho\) under the two formulations are not directly comparable because the penalties have different scales. In the CCAO application, the surrogate provides the strongest combined validation performance and is therefore the principal practical specification.












\section{Application and Evaluation Design}
\label{sec:empirical_design}

{This section specifies the empirical design used to evaluate the proposed correction. We first describe the CCAO application, data, and implementation; then define the chronological validation and regularization-path protocol; next summarize the model comparators; and finally describe the exploratory six-county benchmark. All model comparisons reported in the next section use the evaluation criteria defined in Section~\ref{subsec:metrics}.}


\subsection{{The CCAO Residential Pipeline: Application, Data, and Implementation}}
\label{subsec:ccao_design}

% \textcolor{blue}{\sout{For the empirical analysis, we implement in Python the components of CCAO's residential valuation workflow required for model development and evaluation. The implementation closely follows CCAO's data-preparation, model-fitting, and ratio-study procedures, preserving the operational structure of the existing workflow while enabling controlled evaluation of alternative training objectives.}}

{For the primary application, we implement in Python the components of CCAO's residential valuation workflow required for model development and evaluation. The implementation follows the data, feature, model-fitting, and ratio-study structure of CCAO's published residential pipeline, but it is a research translation rather than a direct production run. The reported results should therefore be interpreted as evidence within this translated workflow, not as official CCAO assessment outputs.}

% \textcolor{blue}{\sout{The proposed methodology modifies only the model-training objective. We hold fixed the sale filters, predictor set, temporal data splits, transformation of predictions to the price scale, ratio-study procedures, and surrounding review process. This restricted intervention serves two purposes. First, it isolates the effect of incorporating vertical equity into model estimation. Second, it demonstrates how the methodology can be integrated into an existing mass-appraisal workflow without redesigning the broader valuation system.}}

{Within the LightGBM comparison, the intervention is deliberately restricted to the training objective. We hold fixed the sale filters, log-price target, underlying predictor information, temporal design, transformation back to the price scale, and evaluation procedures. The 
% \oldtext{frozen Section-2 LightGBM hyperparameter configuration}
\newtext{fixed Section-2 LightGBM hyperparameter configuration} is held fixed across the standard and penalized fits, so $\rho$ is the only dimension varied within each penalized family. 
% \oldtext{This design isolates the practical effect of changing the objective while preserving the surrounding valuation workflow.} 
% \oldtext{Holding the surrounding model specification fixed makes the training-objective change as controlled as possible; the separate $\rho=0$ implementation check below measures any residual difference introduced by the custom-objective code path.}
\newtext{Holding the surrounding model specification fixed makes the training-objective change as controlled as possible. Implementation equivalence of the custom objectives at $\rho=0$ is treated as a reproducibility audit and reported in Appendix~\ref{app:implementation}; until the canonical parity audit is frozen, attribution of positive-penalty changes is made within each custom-objective family.}}

% \oldtext{The CCAO application uses verified Cook County residential sales from the CCAO residential modeling data used for the final evaluation. The residential scope includes the property types represented in the CCAO residential pipeline, including single-family homes, townhomes, small multifamily properties, and condominiums. Following the office pipeline, we exclude multicard PINs and sales marked as outliers.}

\newtext{The primary CCAO application uses verified Cook County residential sales from the final residential-modeling extract used to construct the evaluation samples. The residential scope follows the property types represented in the translated CCAO residential pipeline, including single-family homes, townhomes, small multifamily properties, and condominiums; following the office pipeline, we exclude multicard PINs and sales marked as outliers.}
\todo{Insert the exact final CCAO extract/version, build or retrieval date, and citable source into the first sentence before submission.}

% \textcolor{blue}{\sout{The pre-2024 modeling sample contains 349,661 observations. Sales from 2016--2022 form the development period; 2023 is held out for final testing. The 2024 sales block is reserved for later assessment-year analysis and is not used for the results reported here.}}

% \oldtext{The sales universe is 2016--2024. The current training file contains 2016 sales, so the first observed sale is on 2016-01-01. The oldest 90\% of eligible 2016--2024 sales (344,607 observations, through 2023-11-09) form the development/CV pool used for model fitting and all chronological validation. The newest 10\% (38,290 observations, 2023-11-09 through 2024-12-31) form the primary held-out evaluation. The 2025 block contains 26,641 sales and is reported separately as a secondary forward evaluation after each specification is refit on all 382,897 eligible 2016--2024 sales. The complete prespecified regularization paths are evaluated on the held-out and forward samples for descriptive out-of-sample tradeoff analysis; these samples are not used to modify the baseline configuration, penalty grid, or training design. Because earlier exploratory versions of the project examined 2024 outputs, we describe the test as a held-out evaluation rather than as an untouched or preregistered confirmatory test.}
\newtext{Eligible 2016--2024 sales are ordered chronologically. The oldest 90\% (344,607 sales, through 2023-11-09) form the development/CV pool, and the newest 10\% (38,290 sales, 2023-11-09 through 2024-12-31) form the primary held-out evaluation. The separate 2025 block contains 26,641 sales and is evaluated after each specification is refit on all 382,897 eligible 2016--2024 sales.}
\oldtext{Neither out-of-time sample is used to alter the baseline configuration, penalty grid, objective definitions, or training design.}
\latesttext{The original out-of-time paths had already been inspected during paper development. For the later lower-$\rho$ resolution extension described in Section~\ref{subsec:path_design}, however, the augmented chronological-CV result was frozen before the newly added lower-tail points were evaluated on the held-out and 2025 samples.}
\newtext{Because earlier project development inspected 2024 outputs, we describe the primary sample as held out rather than as an untouched or preregistered confirmatory test.}

\newtext{Because the 90/10 split is defined on chronologically ordered observations rather than on distinct calendar dates, sales occurring on the cutoff date can appear on both sides of the split; no sale observation appears in both samples.}

\begin{table}[!ht]
\centering
\caption{Final CCAO temporal samples.}
\label{tab:ccao_samples}
\begin{tabular}{llr}
\toprule
Sample & Period & Observations \\
\midrule
Development / validation & 2016-01-01--2023-11-09 & 344,607 \\
\makecell[l]{
% \oldtext{Primary held-out test}
\\\newtext{Primary held-out evaluation}} & 2023-11-09--2024-12-31 & 38,290 \\
\makecell[l]{
% \oldtext{Production training}
\\\newtext{Full 2016--2024 refit sample}} & 2016-01-01--2024-12-31 & 382,897 \\
Secondary forward evaluation & 2025-01-01--2025-12-29 & 26,641 \\
\bottomrule
\end{tabular}
\end{table}
% \oldtext{\footnotesize The current training file contains 2016 sales; observed sales begin on 2016-01-01.}
\newtext{\footnotesize Observed eligible sales in the final 2016--2024 modeling universe begin on 2016-01-01.}

The model uses 95 predictors, including property characteristics, location and proximity variables, Census measures, time variables, parcel-shape measures, neighborhood identifiers, and other assessor variables. Twenty-three are treated as categorical.
% \textcolor{blue}{\sout{All predictors are constructed upstream; the penalized models use the same feature set as the LightGBM baseline.}}
{All predictors are constructed upstream by the CCAO pipeline. Standard and penalized LightGBM fits use the same 95 predictors, with the same 23 variables treated as categorical; the penalty introduces no additional predictors or feature engineering. The linear benchmark uses the same underlying predictor information with learner-specific representation where required.}

{Table~\ref{tab:feature_groups} summarizes the predictor information used by the models. The full feature list is fixed before model fitting and should be archived with the replication materials.}

{
% \color{impGreen}
\begin{table}[!ht]
\centering
\caption{CCAO predictor groups used in the empirical models.}
\label{tab:feature_groups}
\begin{tabular}{lrl}
\toprule
Feature group & Count & Examples \\
\midrule
Assessor/property variables & 30 & Class, building area, land area, age, rooms, neighborhood \\
Location variables & 9 & Coordinates, census tract, school district, municipality \\
Proximity variables & 26 & Transit, parks, roads, hospitals, water, airports \\
ACS 5-year measures & 16 & Income, education, poverty, employment, tenure \\
Time variables & 8 & Sale year, month, quarter, day, post-COVID indicator \\
Parcel-shape variables & 6 & Edge/angle dispersion, bounding ratios, vertex count \\
\midrule
Total & 95 & 23 categorical predictors in LightGBM \\
\bottomrule
\end{tabular}
\end{table}
}


\subsection{Temporal Validation and Regularization-Path Design}
\label{subsec:path_design}

Within the 344,607-sale development/CV pool, we use seven expanding-window
rolling-origin folds. The cumulative windows advance in 15-month steps anchored at
2016-01-01, and within each cumulative window the newest 10\% of observations form
validation. For fold $v$, let $\mathcal T_v$ and $\mathcal V_v$ denote its training
and validation samples. The training window expands across folds, and all
training-dependent quantities used by the proposed objectives, including $\bar y$
and the centered log-price deviations, are recomputed from $\mathcal T_v$ only.
% \oldtext{The exact seven-fold date, validation-rule, and sample-size table was displayed in the main text.}
\newtext{The exact fold dates and sample sizes are reported in Appendix~\ref{app:ccao_path_details}, Table~\ref{tab:fold_structure}; because the newest-10\% validation rule is identical in every fold, it is stated once rather than repeated as a table column.}

The experiment traces the full regularization paths rather than choosing a penalty
strength at this stage. 
% \oldtext{The candidate set contains standard LightGBM, a custom-objective implementation control with $\rho=0$, and the direct and sample-additive covariance-penalized LightGBM families.}
\newtext{The substantive comparison contains standard LightGBM and the positive-penalty Direct and Surrogate LightGBM families.}
\oldtext{For each positive-penalty family, we evaluate 50 geometrically spaced values between $0.1$ and $100$.}
\latesttext{Each family was first evaluated on 50 geometrically spaced positive values from $0.1$ to $100$. The initial transition analysis left the lower tail unresolved: the Direct $\operatorname{RMSE}_{\log P}$ event occurred at the first positive grid point, while the Surrogate $\operatorname{RMSE}_{\log P}$ event occurred at $\rho=0$. We therefore fixed a one-time lower-tail extension of 32 additional positive values from $0.0010985$ to $0.0868511$, aligned to the same logarithmic spacing as the original grid. The final augmented path contains 82 positive values per family, and the previously fitted $0.1$--$100$ points are retained unchanged.}
Because the direct and surrogate objectives have different numerical scales, their $\rho$
values are not interpreted as comparable magnitudes. The LightGBM hyperparameters were tuned on the same seven chronological folds before any penalty path was estimated. The tuning criterion was the equal-weight mean of fold validation price RMSE, $\operatorname{RMSE}_P(S)=\sqrt{n_S^{-1}\sum_{i\in S}(\widehat P_i-P_i)^2}$. The adopted configuration uses 994 trees and was then frozen for ordinary LightGBM, Direct, and Surrogate fits. The held-out and 2025 samples were not used for this hyperparameter choice. With that vector held fixed, $\rho$ is the only dimension varied within each penalized family.

\todo{Archive the exact frozen 994-tree LightGBM parameter vector and configuration hash in the replication materials; the main text need not list every parameter.}

For each configuration and fold, we retain the complete set of predictive,
valuation-level, ratio-dispersion, vertical-equity, and mechanism/residual-structure diagnostics.
We report equal-weight fold means and standard deviations to describe the
chronological behavior of the regularization paths.
\oldtext{No $\rho$, model family, or penalized configuration is selected in the present analysis.}
\latesttext{No deployment operating point is selected in the present analysis. The chronological-CV path is used only to characterize the path and, in the screening analysis below, to narrow the positive-$\rho$ configurations that would be passed to a later decision rule.}

\paragraph{\latesttext{CV-derived candidate-region screen.}}
{\color{latestPurple}
The screening rule is generic and is applied on the observed positive grid using $x=\log_{10}\rho$; $\rho=0$ is retained only as the within-family path origin. The benefit family is
$|\mathrm{PRD}-1|$, $|\mathrm{PRB}|$, $|\mathrm{MKI}-1|$, $|\mathrm{VEI}|$, and $|\beta_{\log}|$.
Descriptive continuous piecewise-linear regimes are fit on the observed log-$\rho$ grid, with model complexity chosen by BIC and all candidate events checked against the raw path.
The \emph{activity onset} is the first tested $\rho$ by which at least three of these five diagnostics have entered a supported active-improvement regime; it is a soft marker, not an exclusion or pathology threshold.

For Direct, the upper guardrail is the left edge of the earliest robust cluster of sustained predictive-deterioration breakpoints among $R^2_P$, MAE, MAPE, and $\operatorname{RMSE}_{\log P}$, requiring at least two metrics within two median log-grid spacings and leave-one-fold-out support.
For Surrogate, the first predictive bend is descriptive rather than automatically binding. A predictive-harm event additionally requires that the deteriorating metric has exhausted its advantage relative to the within-family $\rho=0$ control. The Surrogate upper guardrail is the earlier of a stable cross-metric predictive-harm cluster and a stable post-valley rebound in $\Delta_{\mathrm{NL}}$. The latter is interpreted as a \emph{nonlinear-structure caution}---an early warning that non-affine conditional-mean structure has begun worsening---not as the exact onset of a visible S-shaped ratio profile. Distance correlation and the ratio profiles are used to interpret the severity of that warning but do not move the breakpoint by themselves.
The complete construction is repeated in seven leave-one-fold-out CV aggregates. The method specification and endpoints are frozen using chronological CV before the held-out and 2025 paths are used for retrospective portability checks.
}

\paragraph{\latesttext{Descriptive transition diagnostic.}}
\oldtext{After the regularization-path experiment was complete, but before extracting turning events for this diagnostic, we fixed a five-metric rule using $R^2_P$, $\operatorname{MAE}_P$, $\operatorname{MAPE}_P$, $\operatorname{RMSE}_{\log P}$, and COD.}
\latesttext{The five-metric transition rule was fixed in the original path analysis and was left unchanged for the lower-tail extension: it uses $R^2_P$, $\operatorname{MAE}_P$, $\operatorname{MAPE}_P$, $\operatorname{RMSE}_{\log P}$, and COD. On the equal-weight chronological-CV path, the event for $R^2_P$ is its maximum and the events for the other four metrics are their minima, all evaluated on the observed discrete $\rho$ grid. For a penalty family, a common CV transition span is reported only when all five events occur at positive interior grid points; when that condition holds, the span is the smallest interval containing the five event locations. We then recompute the event construction foldwise and in seven leave-one-fold-out aggregations as stability diagnostics. For the v2 lower-tail extension, all 32 new values were first evaluated on the seven chronological folds and the augmented full-CV transition result was frozen before those same new values were evaluated on the held-out and 2025 samples. Because the broader out-of-time paths had already been inspected during paper development, the resulting held-out and 2025 comparison is interpreted as retrospective temporal concordance rather than prospective validation.}

\oldtext{For Direct, we additionally report direction-aware \emph{span regret} on the held-out and 2025 paths: the difference between the split-specific global discrete-grid optimum of a metric and the best value attainable while restricting $\rho$ to the frozen CV span, with the sign defined so that regret is nonnegative. Normalized regret divides this quantity by the metric's full-path range when that range is nonzero. No smoothing, interpolation, near-optimality tolerance, or post-hoc threshold is used. The transition analysis is descriptive and does not select a penalty strength or redefine the prespecified display anchors.}
\oldtext{No smoothing, interpolation, near-optimality tolerance, or post-hoc threshold is used. The transition analysis is descriptive and does not select a penalty strength or redefine the prespecified display anchors.}
\oldtext{For each family with a valid frozen CV transition span, report direction-aware out-of-time span regret as the difference between the split-specific global observed-grid optimum and the best value attainable within that family's frozen CV span, normalized by the full-path metric range when nonzero. No smoothing, interpolation, near-optimality tolerance, or post-hoc threshold is used. The transition analysis is descriptive and does not select a penalty strength or redefine the prespecified display anchors.}
\oldtext{For each family with a valid frozen CV transition span, we report direction-aware out-of-time span regret as the nonnegative difference between the split-specific global observed-grid optimum value and the best value attainable within that family's frozen CV span. We also report normalized regret, defined as raw regret divided by the full-path metric range when that range is nonzero. No smoothing, interpolation, near-optimality tolerance, or post-hoc threshold is used. The transition analysis is descriptive and does not select a penalty strength or redefine the prespecified display anchors.}
\latesttext{For each family with a valid frozen CV transition span, we report direction-aware out-of-time span regret as the nonnegative difference between the split-specific global observed-grid optimum value and the best value attainable within that family's frozen CV span. We also report normalized regret, defined as raw regret divided by the full-path metric range when that range is nonzero. No smoothing, interpolation, near-optimality tolerance, or post-hoc threshold is used. The transition analysis is descriptive, does not select a penalty strength, and is independent of the presentation-only display-anchor convention below.}

% \oldtext{For compact main-text displays only, we use a small set of \emph{prespecified display anchors}: $\rho=0$ together with the positive grid points nearest $0.1$, $1$, $10$, and $100$.}
\oldtext{For compact main-text tables, we use the positive grid points nearest $0.1$, $1$, $10$, and $100$ as prespecified display anchors and show ordinary LightGBM as the unpenalized baseline. These anchors are chosen before inspecting the sweep outcomes and are used only to make tables and ratio-shape figures readable.}
\latesttext{For compact presentation, we use five nominal decade-spaced positive display anchors, $\rho\in\{0.01,0.1,1,10,100\}$, and represent each nominal value by the nearest already-tested point on the frozen 82-point positive grid. The corresponding tested values are $0.0104811$, $0.1$, $0.954095$, $10.4811$, and $100$. This display convention is adopted only to make the tables and ratio-shape figures easy to read; it is not part of the frozen transition rule and does not select a penalty strength. Ordinary LightGBM remains the unpenalized baseline.}
\oldtext{All 50 positive values per family remain in the path figures and machine-readable result table.}
\latesttext{All 82 positive values per family remain in the final path figures and machine-readable result table.}
\oldtext{Ratio-shape figures may additionally show the zero-penalty path origin for visual continuity. Because the direct and surrogate penalties have different scales, the anchors organize each family separately and are not interpreted as equal-strength interventions.}
\latesttext{Ratio-shape figures additionally show the zero-penalty path origin for visual continuity. Because the Direct and Surrogate objectives have different numerical scales, equal nominal display anchors are common presentation coordinates only and must not be interpreted as equal-strength interventions across families.}
\oldtext{Display-anchor refresh: before compilation, extract the four new nominal-$0.01$ rows (Direct/Surrogate on held-out/2025) from the frozen final augmented path at tested $\rho=0.0104811313$, recompute the table-formatting flags from unrounded values, and regenerate \texttt{ratio\_shape\_evolution.pdf} and \texttt{ratio\_shape\_cv\_transition\_span\_only.pdf} using the same nominal anchor convention.}

\oldtext{The same prespecified regularization paths are also evaluated on the two out-of-time samples.}
\latesttext{The same final augmented regularization paths are evaluated on the two out-of-time samples.} For the held-out evaluation, every valid configuration is refit on the full
344,607-sale development pool and evaluated on the 38,290-sale held-out sample. For
the 2025 forward evaluation, the same configuration is refit on all 382,897 eligible
2016--2024 sales and evaluated on the 26,641-sale 2025 sample. \oldtext{These held-out and forward paths are descriptive: they are not used to change the baseline LightGBM configuration, the $\rho$ grid, the objective definitions, or the training design. The chronological CV, held-out, and 2025 paths are therefore three complementary views of the same prespecified experiment rather than sequential model-selection stages.}
\latesttext{These held-out and forward paths are descriptive. For the v2 lower-tail extension, the augmented grid and its CV-derived transition result were frozen before the newly added lower-$\rho$ points were evaluated out of time; those new held-out and forward results were not used to revise the grid, the transition rule, the objective definitions, or the frozen LightGBM configuration. The chronological CV, held-out, and 2025 paths are therefore complementary temporal views of the final augmented experiment, not sequential stages that select a penalty strength.}

All scalar evaluation statistics are computed on the full relevant validation or
evaluation sample after predictions are transformed back to the price scale. For
the ratio-shape figures in the Results, each evaluation sample is partitioned into
30 equal-count sale-price bins and we plot the median valuation ratio in each bin
against sale price on a logarithmic horizontal axis. The binning is used only for
visualization; all scalar metrics are computed on the full evaluation sample.

\begin{table}[t]
\centering
\caption{
% \oldtext{CCAO application and pre-selection regularization-path design.}
\newtext{CCAO application and regularization-path design.}}
\label{tab:ccao_design}
\begin{tabular}{ll}
\toprule
Item & Design \\
\midrule
Population & Verified Cook County residential sales \\
Development / CV & Oldest 90\% of eligible 2016--2024 sales: 344,607, through 2023-11-09 \\
Primary held-out evaluation & Newest 10\%: 38,290 sales, 2023-11-09--2024-12-31 \\
\makecell[l]{
% \oldtext{Production training}
\newtext{Full 2016--2024 refit sample}} & 382,897 eligible 2016--2024 sales \\
2025 forward evaluation & 26,641 sales, 2025-01-01--2025-12-29 \\
Predictors & 95 total; 23 categorical \\
Validation & Seven expanding-window rolling-origin folds; newest 10\% validates \\
Base nonlinear learner & LightGBM on log sale price; 994 trees; hyperparameters frozen before the penalty sweep \\
Corrections evaluated & Direct squared covariance and sample-additive covariance upper bound \\
Penalty path &
\oldtext{50 positive values in $[0.1,100]$ per family}
\latesttext{\makecell[l]{82 positive values per family: original 50-point $[0.1,100]$ path\\plus 32-point lower-tail extension $[0.00110,0.08685]$}} \\
Current analysis & \makecell[l]{\parbox{0.70\textwidth}{\raggedright\oldtext{Full CV, held-out, and forward paths plus descriptive transition-stability analysis; no model selection}}\\\parbox{0.70\textwidth}{\raggedright\latesttext{Full CV, held-out, and forward paths; CV-derived candidate-region screen and descriptive transition diagnostics; no deployment operating-point selection.}}} \\
\bottomrule
\end{tabular}
\end{table}
\todo{V2 experiment update: keep the two-stage grid history explicit in the tracked draft. The lower-tail extension resolves the original near-zero boundary ambiguity; it should not be rewritten later as if all 82 positive values had been prespecified before the first path analysis.}

\subsection{ \newtext{Models and Comparisons}}
\label{subsec:comparators}

% \textcolor{blue}{\sout{Four comparisons are needed to interpret the method.}}
% \oldtext{The empirical comparison separates the baseline model class, the custom-objective implementation path, positive covariance regularization, and a centered post-hoc recalibration benchmark.}
\newtext{The empirical comparison in this paper is deliberately restricted to the baseline model classes, the custom-objective implementation control, and the Direct and Surrogate training-time regularization paths. A systematic comparison with post-hoc calibration is left to follow-up work.}

\paragraph{Linear and nonlinear baselines.}
% \textcolor{blue}{\sout{Linear regression shows the performance of a simpler model; unpenalized LightGBM is the relevant modern AVM baseline.}}
\textcolor{impGreen}{Linear regression provides the simpler historically motivated benchmark; standard LightGBM, trained with the ordinary squared log-price objective, is the relevant nonlinear AVM baseline. Both use the same outcome, development observations, and underlying predictor information under the common temporal design.}

\paragraph{Direct and surrogate penalties.}
% \oldtext{The direct penalty tests the stated covariance target. The surrogate tests the implementation recommended for standard boosting software.}
% \textcolor{blue}{\sout{A custom-objective fit at $\rho=0$ is also needed because LightGBM initialization and training behavior can differ between an ordinary regression objective and a user-supplied objective. This control separates the effect of the penalty from the effect of entering the custom-objective code path.}}
\textcolor{impGreen}{The direct LightGBM comparator uses the diagonal-Hessian approximation described in Appendix~\ref{app:implementation}: its covariance gradient is exact, while the cross-observation second-order curvature is omitted so that the objective can be used through the standard LightGBM interface. The surrogate uses the exact observation-additive objective from Section~\ref{subsec:surrogate}.} 
% \oldtext{The final comparison also includes a custom-objective fit at $\rho=0$. This implementation control separates changes caused by the custom-objective code path from changes caused by a positive penalty.}

% \oldtext{\textbf{Post-hoc rescaling.} The previous draft included a validation-fitted centered rescaling of the baseline predictions as a primary and multi-county comparator.}
\newtext{Post-hoc calibration is outside the empirical scope of the present paper. The fixed-space Direct result in Appendix~\ref{app:theory} makes that comparison scientifically natural, but the current contribution is the design, implementation, mechanism analysis, and evaluation of training-time covariance-guided corrections.}
\todo{Follow-up research: compare in-processing with global and structured post-hoc calibration under matched temporal validation and matched fairness/accuracy targets; do not infer superiority of either family from the present paper.}

% \paragraph{\textcolor{blue}{\sout{Richer information.}}}
% \textcolor{blue}{\sout{Spatial, temporal, neighborhood, and property-characteristic improvements can reduce omitted structure and may improve both accuracy and equity. They are complements to the penalty, not substitutes assumed away by the paper. The preferred sequence is to build the strongest defensible AVM first and then test whether an explicit vertical-equity safeguard still adds value.}}

\textcolor{impGreen}{Better spatial, temporal, neighborhood, and property information can improve both prediction and equity, but this is a modeling extension rather than a comparator in the headline experiment. We therefore discuss richer information as a complementary direction in Section~\ref{sec:discussion} rather than treating it as a completed empirical control here.}


\subsection{Exploratory Six-County ATTOM Benchmark}
\label{subsec:attom_design}

We also run a common ATTOM-derived pipeline in Cook County, IL; Allegheny County, PA; Maricopa County, AZ; King County, WA; Miami-Dade County, FL; and Middlesex County, MA. The benchmark contains 1,231,194 qualified single-family sales and 246,241 held-out sales. Each county uses chronological validation and test blocks and the same modeling code path.

% \textcolor{blue}{\sout{County-specific penalty values are calibrated to the baseline prediction scale.}}
% \oldtext{For each county, eligible recorder transfers are joined to the most recent assessor-history record available before the sale, together with lagged contextual information. The benchmark then orders eligible sales chronologically and reserves the final 20\% as its test block, with model and penalty selection performed only on earlier observations. The benchmark compares standard LightGBM, the custom-objective $\rho=0$ control, positive covariance-penalty values, and validation-fitted centered recalibration of the baseline log predictions. The exact county periods, sample counts, common filters, feature set, and recalibration specification must be archived in Appendix~\ref{app:attom} before final submission.}

\oldtext{For each county, eligible recorder transfers are joined to the most recent assessor-history record available before the sale, together with lagged contextual information. The benchmark then orders eligible sales chronologically and reserves the final 20\% as its test block, with model and penalty selection performed only on earlier observations. The benchmark compares standard LightGBM, the custom-objective $\rho=0$ control, positive covariance-penalty values, and validation-fitted centered recalibration of the baseline log predictions. Appendix~\ref{app:attom} reports the benchmark-specific sample construction, county periods and counts, feature-set and LightGBM-configuration identifiers, chronological validation/test design, and centered-recalibration rule used for the exploratory comparison.}
\oldtext{For each county, eligible recorder transfers are joined to the most recent assessor-history record available before the sale, together with lagged contextual information. The benchmark then orders eligible sales chronologically and reserves the final 20\% as its test block, with model and penalty selection performed only on earlier observations. The benchmark compares standard LightGBM, the custom-objective $\rho=0$ control, and positive covariance-penalty values. Appendix~\ref{app:attom} reports the benchmark-specific sample construction, county periods and counts, feature-set and LightGBM-configuration identifiers, and chronological validation/test design.}
\newtext{For each county, eligible recorder transfers are joined to the most recent assessor-history record available before the sale, together with lagged contextual information. The benchmark then orders eligible sales chronologically and reserves the final 20\% as its test block, with model and penalty selection performed only on earlier observations. The substantive benchmark compares standard LightGBM with positive covariance-penalty values. Appendix~\ref{app:attom} reports the benchmark-specific sample construction, county periods and counts, feature-set and LightGBM-configuration identifiers, and chronological validation/test design.}

This is an exploratory transfer benchmark, not an audit of the six assessor offices. The data are commercially harmonized and do not reproduce each jurisdiction's sales-verification rules, local features, legal setting, or production review. The benchmark was developed iteratively, and its present county configurations were not all frozen under one confirmatory protocol before test evaluation. We therefore use it to examine direction, heterogeneity, and failure modes. It does not provide confirmatory estimates of performance in the six jurisdictions.

\textcolor{impGreen}{Accordingly, the six-county results are kept separate from the primary CCAO application: they test whether the qualitative behavior of the correction transfers across heterogeneous datasets, not whether the six offices would obtain the reported results in their own production systems.}

\section{Results}
\label{sec:results}

% \oldtext{The primary CCAO Results below report three additions to the complete pre-selection analysis: the split-specific $\rho=0$ implementation-control audit, a fixed out-of-time $\Delta_{\mathrm{NL}}$ diagnostic, and a validation-neutral centered-recalibration benchmark. $\Delta_{\mathrm{NL}}$ is computed on the primary held-out and 2025 forward evaluations; the already completed chronological CV is left unchanged.}

\oldtext{The primary CCAO Results report the complete Direct and Surrogate regularization paths together with three complementary checks: the $\rho=0$ implementation control, the fixed cross-fitted correlation-ratio nonlinearity gap $\Delta_{\mathrm{NL}}$, and the validation-fitted centered-recalibration benchmark. $\Delta_{\mathrm{NL}}$ is reported for the primary held-out and 2025 forward evaluations but is not computed retrospectively for the completed chronological CV. No penalty family, penalty strength, or deployment operating point is selected.}
\oldtext{The primary CCAO Results report the complete Direct and Surrogate regularization paths together with two complementary checks: the $\rho=0$ implementation control and the fixed cross-fitted correlation-ratio nonlinearity gap $\Delta_{\mathrm{NL}}$. $\Delta_{\mathrm{NL}}$ is reported for the primary held-out and 2025 forward evaluations but is not computed retrospectively for the completed chronological CV. No penalty family, penalty strength, or deployment operating point is selected.}
\oldtext{The primary CCAO Results report the complete Direct and Surrogate regularization paths together with the fixed cross-fitted correlation-ratio nonlinearity gap $\Delta_{\mathrm{NL}}$. $\Delta_{\mathrm{NL}}$ is reported for the primary held-out and 2025 forward evaluations but is not computed retrospectively for the completed chronological CV. No penalty family, penalty strength, or deployment operating point is selected.}
\oldtext{The primary CCAO Results report the complete Direct and Surrogate regularization paths together with the fixed cross-fitted correlation-ratio nonlinearity gap $\Delta_{\mathrm{NL}}$. Held-out and 2025 $\Delta_{\mathrm{NL}}$ remain the original out-of-time diagnostics. Chronological-CV $\Delta_{\mathrm{NL}}$ was retrospectively computed from the already-frozen validation predictions using the same fixed cross-fitted estimator; it was not used to define the frozen prediction/COD transition span, the $\rho$ grid, or a model-selection rule. No penalty family, penalty strength, or deployment operating point is selected.}
\latesttext{The primary CCAO Results report the complete Direct and Surrogate regularization paths together with the fixed cross-fitted correlation-ratio nonlinearity gap $\Delta_{\mathrm{NL}}$. Held-out and 2025 $\Delta_{\mathrm{NL}}$ remain the original out-of-time diagnostics, while chronological-CV $\Delta_{\mathrm{NL}}$ is reconstructed from the already-frozen validation predictions using the same fixed estimator. In addition to the older prediction/COD transition diagnostic, we report the generic CV-derived candidate-region screen from Section~\ref{subsec:path_design}. That screen removes low-activity positive penalties and places a conservative upper guardrail before the high-$\rho$ deterioration/caution regime, but it does not rank surviving configurations or select a deployment operating point.}
\oldtext{We also report the frozen five-metric transition diagnostic from Section~\ref{subsec:path_design}, including fold and leave-one-fold-out stability, retrospective temporal concordance, and Direct span-regret calculations. This diagnostic characterizes the path and does not define a selected or recommended range.}
\oldtext{We also report the frozen five-metric transition diagnostic from Section~\ref{subsec:path_design}, including fold and leave-one-fold-out stability and retrospective temporal concordance. This diagnostic characterizes the path and does not define a selected or recommended range.}
\latesttext{We also report the frozen five-metric transition diagnostic from Section~\ref{subsec:path_design}, including fold and leave-one-fold-out stability, retrospective exact concordance, and family-symmetric value-based span regret. This diagnostic characterizes the path and does not define a selected or recommended range.}

\subsection{Accuracy Alone Does Not Remove the CCAO Ratio Trend}
\label{subsec:ccao_baseline_results}

Section~\ref{subsec:baseline_tension} reports the Linear versus ordinary LightGBM
comparison on the held-out and 2025 forward samples. Those baseline artifacts are
regenerated under the corrected metric code and the frozen Section-2 LightGBM
configuration; they do not involve a penalized-model choice. The remainder of the
primary CCAO Results treats the \emph{regularization path itself} as the empirical
object of interest. No $\rho$, penalty family, or operating point is selected in the
present analysis.

\oldtext{A dedicated main-text subsection and table previously reported the $\rho=0$ native/custom implementation check.}



\subsection{\texorpdfstring{\oldtext{What the Prespecified Regularization Paths Do}\latesttext{What the Final Augmented Regularization Paths Show}}{What the Final Augmented Regularization Paths Show}}
\label{subsec:path_results}

\oldtext{The Direct and Surrogate objectives are evaluated over their complete prespecified grids.}
\latesttext{The Direct and Surrogate objectives are evaluated over their complete final augmented grids.} Since the two penalties have different numerical scales, their
$\rho$ values are not compared as equal-strength interventions. The principal
question at this stage is descriptive: how do prediction, valuation level,
horizontal uniformity, assessor-facing vertical-equity diagnostics, and the targeted
dependence measures evolve as regularization increases?

\oldtext{Table~\ref{tab:path_anchor_summary} is the compact main-text view. It uses only the prespecified display anchors defined in Section~\ref{subsec:path_design}; the full 82-point positive path for each family remains available in the complete path table and figures. No row in this table is selected because of its observed performance.}
\oldtext{Table~\ref{tab:path_anchor_summary} is the compact main-text view of prediction and assessor-facing vertical equity at the fixed display anchors defined in Section~\ref{subsec:path_design}. Appendix Table~\ref{tab:path_anchor_complementary} reports valuation level, horizontal uniformity, and mechanism/residual-structure diagnostics at the same anchors. The full 82-point positive path for each family remains available in the complete path table and figures, and no anchor row is selected because of its observed performance.}
\latesttext{Table~\ref{tab:path_anchor_summary} is the compact main-text view of prediction and assessor-facing vertical equity at the nominal decade-spaced display anchors defined in Section~\ref{subsec:path_design}. Appendix Table~\ref{tab:path_anchor_complementary} reports valuation level, horizontal uniformity, and mechanism/residual-structure diagnostics at the same nominal anchors. The full 82-point positive path for each family remains the scientific object of analysis; the compact rows are presentation-only snapshots and do not constitute model selection.}

% \begin{table}[!htbp]
% \centering
% \scriptsize
% \caption{Regularization-path summary at prespecified display anchors. Rows are fixed for display before inspecting outcomes and do not constitute model selection.}
% \label{tab:path_anchor_summary}
% \resizebox{\textwidth}{!}{%
% \begin{tabular}{llrrrrrrrr}
% \toprule
% Family & $\rho$ & $R^2_P$ & MAE & PRD & PRB & MKI & VEI & $\beta_{\log}$ & dCor \\
% \midrule
% \multicolumn{11}{l}{\textit{Panel A: held-out evaluation}} \\
% Ordinary LightGBM & -- & 0.894 & \$75,655 & 1.069 & -0.091 & 0.923 & -26.5\% & -0.150 & 0.387 \\
% Direct & 0 & 0.893 & \$75,976 & 1.069 & -0.088 & 0.923 & -26.2\% & -0.147 & 0.382 \\
% Direct & $\approx0.1$ & 0.894 & \$76,235 & 1.068 & -0.088 & 0.926 & -27.0\% & -0.147 & 0.382 \\
% Direct & $\approx0.954$ & 0.899 & \$74,485 & 1.060 & -0.075 & 0.940 & -21.9\% & -0.134 & 0.359 \\
% Direct & $\approx10.481$ & 0.895 & \$75,218 & 1.036 & -0.040 & 0.984 & -10.5\% & -0.096 & 0.302 \\
% Direct & 100 & 0.869 & \$84,918 & 1.029 & -0.015 & 0.998 & 0.7\% & -0.079 & 0.265 \\
% Surrogate & 0 & 0.893 & \$75,976 & 1.069 & -0.088 & 0.923 & -26.2\% & -0.147 & 0.382 \\
% Surrogate & $\approx0.1$ & 0.894 & \$75,897 & 1.064 & -0.079 & 0.931 & -22.7\% & -0.138 & 0.367 \\
% Surrogate & $\approx0.954$ & 0.897 & \$75,468 & 1.049 & -0.045 & 0.953 & -12.6\% & -0.102 & 0.310 \\
% Surrogate & $\approx10.481$ & 0.893 & \$78,110 & 1.024 & 0.013 & 0.988 & 0.2\% & -0.036 & 0.250 \\
% Surrogate & 100 & 0.889 & \$82,094 & 1.010 & 0.042 & 1.013 & 4.0\% & -0.001 & 0.267 \\
% \midrule
% \multicolumn{11}{l}{\textit{Panel B: 2025 forward evaluation}} \\
% Ordinary LightGBM & -- & 0.904 & \$78,484 & 1.079 & -0.106 & 0.907 & -28.6\% & -0.164 & 0.422 \\
% Direct & 0 & 0.904 & \$79,139 & 1.078 & -0.103 & 0.909 & -28.4\% & -0.161 & 0.417 \\
% Direct & $\approx0.1$ & 0.906 & \$78,029 & 1.077 & -0.102 & 0.910 & -27.7\% & -0.160 & 0.416 \\
% Direct & $\approx0.954$ & 0.910 & \$77,139 & 1.069 & -0.090 & 0.926 & -23.9\% & -0.148 & 0.392 \\
% Direct & $\approx10.481$ & 0.902 & \$79,187 & 1.042 & -0.049 & 0.971 & -11.9\% & -0.105 & 0.318 \\
% Direct & 100 & 0.873 & \$89,309 & 1.036 & -0.029 & 0.983 & -4.5\% & -0.093 & 0.281 \\
% Surrogate & 0 & 0.904 & \$79,139 & 1.078 & -0.103 & 0.909 & -28.4\% & -0.161 & 0.417 \\
% Surrogate & $\approx0.1$ & 0.905 & \$78,546 & 1.074 & -0.095 & 0.914 & -24.7\% & -0.153 & 0.402 \\
% Surrogate & $\approx0.954$ & 0.908 & \$78,196 & 1.060 & -0.062 & 0.932 & -15.3\% & -0.120 & 0.343 \\
% Surrogate & $\approx10.481$ & 0.905 & \$81,075 & 1.033 & 0.000 & 0.967 & -2.0\% & -0.050 & 0.258 \\
% Surrogate & 100 & 0.900 & \$85,257 & 1.020 & 0.030 & 0.987 & 1.7\% & -0.015 & 0.266 \\
% \bottomrule
% \end{tabular}}
% \end{table}

% \begin{table}[!htbp]
% \centering
% \scriptsize
% \setlength{\tabcolsep}{3.0pt}
% \renewcommand{\arraystretch}{1.08}

% \caption{Regularization-path summary at prespecified display anchors.
% Rows are fixed for display before inspecting outcomes and do not constitute
% model selection. Boldface identifies point estimates that fall within the
% applicable IAAO reference range.}
% \label{tab:path_anchor_summary}

% \resizebox{\textwidth}{!}{%
% \begin{tabular}{llrrrrrrrr}
% \toprule
% Method & $\rho$
% & $R^2_P$
% & MAE
% & PRD
% & PRB
% & MKI
% & VEI
% & $\beta_{\log}$
% & dCor \\
% \midrule

% \multicolumn{10}{l}{\textit{Panel A: Held-out evaluation}} \\
% \addlinespace[2pt]

% \multicolumn{10}{l}{\textbf{Baseline models}} \\
% \cmidrule(lr){1-10}
% Linear regression
% & --
% & 0.799
% & \$90,092
% & 1.042
% & \textbf{-0.016}
% & 0.975
% & -11.6\%
% & -0.092
% & 0.250 \\

% Ordinary LightGBM
% & --
% & 0.894
% & \$75,655
% & 1.069
% & -0.091
% & 0.923
% & -26.5\%
% & -0.150
% & 0.387 \\

% \addlinespace[3pt]
% \multicolumn{10}{l}{\textbf{Direct penalty}} \\
% \cmidrule(lr){1-10}
% Direct
% & 0
% & 0.893
% & \$75,976
% & 1.069
% & -0.088
% & 0.923
% & -26.2\%
% & -0.147
% & 0.382 \\

% Direct
% & $\approx 0.10$
% & 0.894
% & \$76,235
% & 1.068
% & -0.088
% & 0.926
% & -27.0\%
% & -0.147
% & 0.382 \\

% Direct
% & $\approx 0.95$
% & 0.899
% & \$74,485
% & 1.060
% & -0.075
% & 0.940
% & -21.9\%
% & -0.134
% & 0.359 \\

% Direct
% & $\approx 10.48$
% & 0.895
% & \$75,218
% & 1.036
% & \textbf{-0.040}
% & 0.984
% & -10.5\%
% & -0.096
% & 0.302 \\

% Direct
% & 100
% & 0.869
% & \$84,918
% & \textbf{1.029}
% & \textbf{-0.015}
% & 0.998
% & \textbf{0.7\%}
% & -0.079
% & 0.265 \\

% \addlinespace[3pt]
% \multicolumn{10}{l}{\textbf{Surrogate penalty}} \\
% \cmidrule(lr){1-10}
% Surrogate
% & 0
% & 0.893
% & \$75,976
% & 1.069
% & -0.088
% & 0.923
% & -26.2\%
% & -0.147
% & 0.382 \\

% Surrogate
% & $\approx 0.10$
% & 0.894
% & \$75,897
% & 1.064
% & -0.079
% & 0.931
% & -22.7\%
% & -0.138
% & 0.367 \\

% Surrogate
% & $\approx 0.95$
% & 0.897
% & \$75,468
% & 1.049
% & \textbf{-0.045}
% & 0.953
% & -12.6\%
% & -0.102
% & 0.310 \\

% Surrogate
% & $\approx 10.48$
% & 0.893
% & \$78,110
% & \textbf{1.024}
% & \textbf{0.013}
% & 0.988
% & \textbf{0.2\%}
% & -0.036
% & 0.250 \\

% Surrogate
% & 100
% & 0.889
% & \$82,094
% & \textbf{1.010}
% & \textbf{0.042}
% & 1.013
% & \textbf{4.0\%}
% & -0.001
% & 0.267 \\


% \midrule
% \addlinespace[2pt]
% \multicolumn{10}{l}{\textit{Panel B: 2025 forward evaluation}} \\
% \addlinespace[2pt]

% \multicolumn{10}{l}{\textbf{Baseline models}} \\
% \cmidrule(lr){1-10}
% Linear regression
% & --
% & 0.799
% & \$99,371
% & 1.052
% & \textbf{-0.029}
% & 0.954
% & -17.3\%
% & -0.109
% & 0.269 \\

% Ordinary LightGBM
% & --
% & 0.904
% & \$78,484
% & 1.079
% & -0.106
% & 0.907
% & -28.6\%
% & -0.164
% & 0.422 \\

% \addlinespace[3pt]
% \multicolumn{10}{l}{\textbf{Direct penalty}} \\
% \cmidrule(lr){1-10}
% Direct
% & 0
% & 0.904
% & \$79,139
% & 1.078
% & -0.103
% & 0.909
% & -28.4\%
% & -0.161
% & 0.417 \\

% Direct
% & $\approx 0.10$
% & 0.906
% & \$78,029
% & 1.077
% & -0.102
% & 0.910
% & -27.7\%
% & -0.160
% & 0.416 \\

% Direct
% & $\approx 0.95$
% & 0.910
% & \$77,139
% & 1.069
% & -0.090
% & 0.926
% & -23.9\%
% & -0.148
% & 0.392 \\

% Direct
% & $\approx 10.48$
% & 0.902
% & \$79,187
% & 1.042
% & \textbf{-0.049}
% & 0.971
% & -11.9\%
% & -0.105
% & 0.318 \\

% Direct
% & 100
% & 0.873
% & \$89,309
% & 1.036
% & \textbf{-0.029}
% & 0.983
% & \textbf{-4.5\%}
% & -0.093
% & 0.281 \\

% \addlinespace[3pt]
% \multicolumn{10}{l}{\textbf{Surrogate penalty}} \\
% \cmidrule(lr){1-10}
% Surrogate
% & 0
% & 0.904
% & \$79,139
% & 1.078
% & -0.103
% & 0.909
% & -28.4\%
% & -0.161
% & 0.417 \\

% Surrogate
% & $\approx 0.10$
% & 0.905
% & \$78,546
% & 1.074
% & -0.095
% & 0.914
% & -24.7\%
% & -0.153
% & 0.402 \\

% Surrogate
% & $\approx 0.95$
% & 0.908
% & \$78,196
% & 1.060
% & -0.062
% & 0.932
% & -15.3\%
% & -0.120
% & 0.343 \\

% Surrogate
% & $\approx 10.48$
% & 0.905
% & \$81,075
% & 1.033
% & \textbf{0.000}
% & 0.967
% & \textbf{-2.0\%}
% & -0.050
% & 0.258 \\

% Surrogate
% & 100
% & 0.900
% & \$85,257
% & \textbf{1.020}
% & \textbf{0.030}
% & 0.987
% & \textbf{1.7\%}
% & -0.015
% & 0.266 \\

% \bottomrule
% \end{tabular}
% }

% \vspace{1mm}
% \begin{minipage}{\textwidth}
% \scriptsize
% \emph{Notes.}
% Display anchors are prespecified and are not selected on the basis of observed
% performance. Positive $\rho$ values are shown to two decimal places; Direct and
% Surrogate penalty magnitudes are not numerically comparable across families.
% Boldface indicates that the reported point estimate falls within an applicable
% IAAO reference range: PRD $[0.98,1.03]$ and PRB $[-0.05,0.05]$ follow the
% 2013 IAAO Standard on Ratio Studies, while VEI $[-10\%,10\%]$ follows the
% May 2026 IAAO exposure draft and is used here as descriptive guidance rather
% than as a final adopted compliance standard. MKI has a neutral reference value
% of one but no corresponding IAAO acceptance band used in this paper, so MKI
% values are not bolded. Boldface does not denote statistical significance,
% formal compliance, or model selection.
% \end{minipage}

% \end{table}

\begin{oldrevisionblock}
\noindent\textbf{Previous combined representative-$\rho$ table.}\par\smallskip
\centering
\scriptsize
\renewcommand{\arraystretch}{1.08}
\resizebox{\textwidth}{!}{%
\begin{tabular}{llrrrrrrrrr}
\toprule
Method & $\rho$
& $R^2_P$
& MAE
& PRD
& PRB
& MKI
& VEI
& $\beta_{\log}$
& $\Delta_{\mathrm{NL}}$
& dCor \\
\midrule
\multicolumn{11}{l}{\textit{Panel A: Held-out evaluation}} \\
\addlinespace[2pt]
\multicolumn{11}{l}{\textbf{Baseline models}} \\
\cmidrule(lr){1-11}
Linear regression & -- & 0.799 & \$90,092 & \textbf{1.042} & \textbf{-0.016} & \textbf{0.975} & \textbf{-11.6\%} & \textbf{-0.092} & 0.131 & \textbf{0.250} \\
Ordinary LightGBM & -- & \textbf{0.894} & \textbf{\$75,655} & 1.069 & -0.091 & 0.923 & -26.5\% & -0.150 & \textbf{0.119} & 0.387 \\
\addlinespace[3pt]
\multicolumn{11}{l}{\textbf{Direct penalty}} \\
\cmidrule(lr){1-11}
Direct & $\approx0.1$ & 0.894 & \$76,235 & 1.068 & -0.088 & 0.926 & -27.0\% & -0.147 & 0.119 & 0.382 \\
Direct & $\approx0.954$ & \textbf{0.899} & \textbf{\$74,485} & 1.060 & -0.075 & 0.940 & -21.9\% & -0.134 & 0.122 & 0.359 \\
Direct & $\approx10.481$ & \textbf{0.895} & \textbf{\$75,218} & \textbf{1.036} & -0.040\textsuperscript{*} & \textbf{0.984} & \textbf{-10.5\%} & -0.096 & 0.132 & 0.302 \\
Direct & 100 & 0.869 & \$84,918 & \textbf{1.029} & \textbf{-0.015} & \textbf{0.998} & \textbf{0.7\%} & \textbf{-0.079} & \textbf{0.115} & 0.265 \\
\addlinespace[3pt]
\multicolumn{11}{l}{\textbf{Surrogate penalty}} \\
\cmidrule(lr){1-11}
Surrogate & $\approx0.1$ & 0.894 & \$75,897 & 1.064 & -0.079 & 0.931 & -22.7\% & -0.138 & \textbf{0.113} & 0.367 \\
Surrogate & $\approx0.954$ & \textbf{0.897} & \textbf{\$75,468} & 1.049 & -0.045\textsuperscript{*} & 0.953\textsuperscript{*} & -12.6\% & -0.102 & \textbf{0.099} & 0.310 \\
Surrogate & $\approx10.481$ & 0.893 & \$78,110 & \textbf{1.024} & \textbf{0.013} & \textbf{0.988} & \textbf{0.2\%} & \textbf{-0.036} & \textbf{0.103} & 0.250 \\
Surrogate & 100 & 0.889 & \$82,094 & \textbf{1.010} & 0.042\textsuperscript{*} & \textbf{1.013} & \textbf{4.0\%} & \textbf{-0.001} & 0.124 & 0.267 \\
\midrule
\addlinespace[2pt]
\multicolumn{11}{l}{\textit{Panel B: 2025 forward evaluation}} \\
\addlinespace[2pt]
\multicolumn{11}{l}{\textbf{Baseline models}} \\
\cmidrule(lr){1-11}
Linear regression & -- & 0.799 & \$99,371 & \textbf{1.052} & \textbf{-0.029} & \textbf{0.954} & \textbf{-17.3\%} & \textbf{-0.109} & 0.122 & \textbf{0.269} \\
Ordinary LightGBM & -- & \textbf{0.904} & \textbf{\$78,484} & 1.079 & -0.106 & 0.907 & -28.6\% & -0.164 & \textbf{0.121} & 0.422 \\
\addlinespace[3pt]
\multicolumn{11}{l}{\textbf{Direct penalty}} \\
\cmidrule(lr){1-11}
Direct & $\approx0.1$ & \textbf{0.906} & \textbf{\$78,029} & 1.077 & -0.102 & 0.910 & -27.7\% & -0.160 & 0.121 & 0.416 \\
Direct & $\approx0.954$ & \textbf{0.910} & \textbf{\$77,139} & 1.069 & -0.090 & 0.926 & -23.9\% & -0.148 & 0.127 & 0.392 \\
Direct & $\approx10.481$ & 0.902 & \$79,187 & \textbf{1.042} & -0.049\textsuperscript{*} & \textbf{0.971} & \textbf{-11.9\%} & \textbf{-0.105} & 0.134 & 0.318 \\
Direct & 100 & 0.873 & \$89,309 & \textbf{1.036} & -0.029\textsuperscript{*} & \textbf{0.983} & \textbf{-4.5\%} & \textbf{-0.093} & \textbf{0.119} & 0.281 \\
\addlinespace[3pt]
\multicolumn{11}{l}{\textbf{Surrogate penalty}} \\
\cmidrule(lr){1-11}
Surrogate & $\approx0.1$ & \textbf{0.905} & \$78,546 & 1.074 & -0.095 & 0.914 & -24.7\% & -0.153 & \textbf{0.116} & 0.402 \\
Surrogate & $\approx0.954$ & \textbf{0.908} & \textbf{\$78,196} & 1.060 & -0.062 & 0.932 & \textbf{-15.3\%} & -0.120 & \textbf{0.096} & 0.343 \\
Surrogate & $\approx10.481$ & \textbf{0.905} & \$81,075 & \textbf{1.033} & \textbf{0.000} & \textbf{0.967} & \textbf{-2.0\%} & \textbf{-0.050} & \textbf{0.092} & \textbf{0.258} \\
Surrogate & 100 & 0.900 & \$85,257 & \textbf{1.020} & 0.030\textsuperscript{*} & \textbf{0.987} & \textbf{1.7\%} & \textbf{-0.015} & \textbf{0.111} & \textbf{0.266} \\
\bottomrule
\end{tabular}}
\vspace{1mm}
\begin{minipage}{\textwidth}
\scriptsize
\emph{Notes.}
Display anchors are prespecified and are not selected on the basis of observed
performance. Positive $\rho$ values are shown to two decimal places; Direct and
Surrogate penalty magnitudes are not numerically comparable across families.
Boldface indicates a point estimate equal to or better than both baseline models
for that metric. 
% \oldtext{An asterisk indicates that the reported point estimate falls within the applicable reference range.}
For penalized rows, an asterisk indicates improvement relative to Ordinary LightGBM without equaling or improving upon the better baseline, provided that the point estimate also lies within the applicable reference range: PRD $[0.98,1.03]$, PRB $[-0.05,0.05]$, MKI $[0.95,1.05]$, and VEI $[-10\%,10\%]$. Higher $R^2_P$ and lower MAE are preferred for prediction. PRD and MKI are
compared by distance from one, while PRB and VEI are compared by distance from
zero. For the mechanism/residual-structure diagnostics, smaller
$|\beta_{\log}|$, $\Delta_{\mathrm{NL}}$, and dCor indicate less first-order,
non-affine conditional-mean, and broader residual--price structure, respectively;
these diagnostic directions are not assessor-standard equity targets.
$\Delta_{\mathrm{NL}}$ is the paper-specific shorthand for the correlation-ratio
nonlinearity gap defined in Eq.~\eqref{eq:nonlinearity_gap} and has no external
reference band. Boldface and asterisks do not denote statistical significance,
formal compliance, or model selection.
\end{minipage}
\end{oldrevisionblock}

% {\color{latestPurple}
% \begin{table}[!htbp]
% \centering
% \scriptsize
% \renewcommand{\arraystretch}{1.08}
% \caption{\oldtext{Primary regularization-path summary at prespecified display anchors: prediction and assessor-facing vertical equity. Rows are fixed for display and do not constitute model selection. Boldface identifies point estimates that equal or improve upon both baselines for the corresponding metric. An asterisk identifies penalized point estimates that equal or improve upon Ordinary LightGBM but not the better baseline, and only when the point estimate also lies within the applicable reference range.}\latesttext{Primary regularization-path summary at nominal decade-spaced display anchors: prediction and assessor-facing vertical equity. Rows follow the fixed presentation rule $\rho\in\{0.01,0.1,1,10,100\}$ (nearest tested grid points) and do not constitute model selection. Within the two baseline rows, boldface identifies the baseline closer to the relevant ideal or preferred value. For Direct and Surrogate rows, boldface identifies point estimates that improve upon Ordinary LightGBM for the corresponding metric; an asterisk identifies those that also improve upon Linear Regression.}}
% \label{tab:path_anchor_summary}
% \resizebox{\textwidth}{!}{%
% \begin{tabular}{llrrrrrrrr}
% \toprule
% Method & $\rho$
% & $R^2_P$
% & MAE
% & MAPE
% & $\operatorname{RMSE}_{\log P}$
% & PRD
% & PRB
% & MKI
% & VEI \\
% \midrule
% \multicolumn{10}{l}{\textit{Panel A: Held-out evaluation}} \\
% \addlinespace[2pt]
% \multicolumn{10}{l}{\textbf{Baseline models}} \\
% \cmidrule(lr){1-10}
% Linear regression & -- & 0.799 & \$90,092 & 24.1\% & 0.322 & \textbf{1.042} & \textbf{-0.016} & \textbf{0.975} & \textbf{-11.6\%} \\
% Ordinary LightGBM & -- & \textbf{0.894} & \textbf{\$75,655} & \textbf{21.2\%} & \textbf{0.289} & 1.069 & -0.091 & 0.923 & -26.5\% \\
% \addlinespace[3pt]
% \multicolumn{10}{l}{\textbf{Direct penalty}} \\
% \cmidrule(lr){1-10}
% \oldtext{Direct} & \oldtext{$\approx0.01$} & \oldtext{\textit{pending canonical v2 extraction}} \\
% \latesttext{Direct} & \latesttext{$\approx0.01$} & \latesttext{0.893} & \latesttext{\$76,058} & \latesttext{21.2\%} & \latesttext{0.290} & \latesttext{1.069} & \latesttext{\textbf{-0.089}} & \latesttext{\textbf{0.924}} & \latesttext{-26.9\%} \\
% Direct & \oldtext{$\approx0.1$}\latesttext{$0.1$} & 0.894 & \$76,235 & \oldtext{\textit{source}}\latesttext{21.3\%} & \oldtext{\textit{source}}\latesttext{0.291} & \textbf{1.068} & \textbf{-0.088} & \textbf{0.926} & -27.0\% \\
% Direct & \oldtext{$\approx0.954$}\latesttext{$\approx1$} & \textbf{0.899}\textsuperscript{*} & \textbf{\$74,485}\textsuperscript{*} & \oldtext{\textit{source}}\latesttext{\textbf{21.1\%}\textsuperscript{*}} & \oldtext{\textit{source}}\latesttext{0.290} & \textbf{1.060} & \textbf{-0.075} & \textbf{0.940} & \textbf{-21.9\%} \\
% Direct & \oldtext{$\approx10.481$}\latesttext{$\approx10$} & \textbf{0.895}\textsuperscript{*} & \textbf{\$75,218}\textsuperscript{*} & \oldtext{\textit{source}}\latesttext{21.4\%} & \oldtext{\textit{source}}\latesttext{0.294} & \textbf{1.036}\textsuperscript{*} & \textbf{-0.040} & \textbf{0.984}\textsuperscript{*} & \textbf{-10.5\%}\textsuperscript{*} \\
% Direct & 100 & 0.869 & \$84,918 & \oldtext{\textit{source}}\latesttext{23.9\%} & \oldtext{\textit{source}}\latesttext{0.325} & \textbf{1.029}\textsuperscript{*} & \textbf{-0.015}\textsuperscript{*} & \textbf{0.998}\textsuperscript{*} & \textbf{0.7\%}\textsuperscript{*} \\
% \addlinespace[3pt]
% \multicolumn{10}{l}{\textbf{Surrogate penalty}} \\
% \cmidrule(lr){1-10}
% \oldtext{Surrogate} & \oldtext{$\approx0.01$} & \oldtext{\textit{pending canonical v2 extraction}} \\
% \latesttext{Surrogate} & \latesttext{$\approx0.01$} & \latesttext{0.894} & \latesttext{\$75,748} & \latesttext{\textbf{21.2\%}\textsuperscript{*}} & \latesttext{0.289} & \latesttext{1.069} & \latesttext{\textbf{-0.088}} & \latesttext{\textbf{0.924}} & \latesttext{\textbf{-25.8\%}} \\
% Surrogate & \oldtext{$\approx0.1$}\latesttext{$0.1$} & 0.894 & \$75,897 & \oldtext{\textit{source}}\latesttext{\textbf{21.1\%}\textsuperscript{*}} & \oldtext{\textit{source}}\latesttext{0.292} & \textbf{1.064} & \textbf{-0.079} & \textbf{0.931} & \textbf{-22.7\%} \\
% Surrogate & \oldtext{$\approx0.954$}\latesttext{$\approx1$} & \textbf{0.897}\textsuperscript{*} & \textbf{\$75,468}\textsuperscript{*} & \oldtext{\textit{source}}\latesttext{\textbf{21.0\%}\textsuperscript{*}} & \oldtext{\textit{source}}\latesttext{0.298} & \textbf{1.049} & \textbf{-0.045} & \textbf{0.953} & \textbf{-12.6\%} \\
% Surrogate & \oldtext{$\approx10.481$}\latesttext{$\approx10$} & 0.893 & \$78,110 & \oldtext{\textit{source}}\latesttext{21.6\%} & \oldtext{\textit{source}}\latesttext{0.329} & \textbf{1.024}\textsuperscript{*} & \textbf{0.013}\textsuperscript{*} & \textbf{0.988}\textsuperscript{*} & \textbf{0.2\%}\textsuperscript{*} \\
% Surrogate & 100 & 0.889 & \$82,094 & \oldtext{\textit{source}}\latesttext{22.8\%} & \oldtext{\textit{source}}\latesttext{0.357} & \textbf{1.010}\textsuperscript{*} & \textbf{0.042} & \textbf{1.013}\textsuperscript{*} & \textbf{4.0\%}\textsuperscript{*} \\
% \midrule
% \addlinespace[2pt]
% \multicolumn{10}{l}{\textit{Panel B: 2025 forward evaluation}} \\
% \addlinespace[2pt]
% \multicolumn{10}{l}{\textbf{Baseline models}} \\
% \cmidrule(lr){1-10}
% Linear regression & -- & 0.799 & \$99,371 & 24.9\% & 0.313 & \textbf{1.052} & \textbf{-0.029} & \textbf{0.954} & \textbf{-17.3\%} \\
% Ordinary LightGBM & -- & \textbf{0.904} & \textbf{\$78,484} & \textbf{20.8\%} & \textbf{0.278} & 1.079 & -0.106 & 0.907 & -28.6\% \\
% \addlinespace[3pt]
% \multicolumn{10}{l}{\textbf{Direct penalty}} \\
% \cmidrule(lr){1-10}
% \oldtext{Direct} & \oldtext{$\approx0.01$} & \oldtext{\textit{pending canonical v2 extraction}} \\
% \latesttext{Direct} & \latesttext{$\approx0.01$} & \latesttext{\textbf{0.905}\textsuperscript{*}} & \latesttext{\textbf{\$78,442}\textsuperscript{*}} & \latesttext{\textbf{20.7\%}\textsuperscript{*}} & \latesttext{\textbf{0.278}\textsuperscript{*}} & \latesttext{\textbf{1.078}} & \latesttext{\textbf{-0.105}} & \latesttext{\textbf{0.909}} & \latesttext{-29.8\%} \\
% Direct & \oldtext{$\approx0.1$}\latesttext{$0.1$} & \textbf{0.906}\textsuperscript{*} & \textbf{\$78,029}\textsuperscript{*} & \oldtext{\textit{source}}\latesttext{\textbf{20.7\%}\textsuperscript{*}} & \oldtext{\textit{source}}\latesttext{\textbf{0.278}\textsuperscript{*}} & \textbf{1.077} & \textbf{-0.102} & \textbf{0.910} & \textbf{-27.7\%} \\
% Direct & \oldtext{$\approx0.954$}\latesttext{$\approx1$} & \textbf{0.910}\textsuperscript{*} & \textbf{\$77,139}\textsuperscript{*} & \oldtext{\textit{source}}\latesttext{\textbf{20.7\%}\textsuperscript{*}} & \oldtext{\textit{source}}\latesttext{0.279} & \textbf{1.069} & \textbf{-0.090} & \textbf{0.926} & \textbf{-23.9\%} \\
% Direct & \oldtext{$\approx10.481$}\latesttext{$\approx10$} & 0.902 & \$79,187 & \oldtext{\textit{source}}\latesttext{20.9\%} & \oldtext{\textit{source}}\latesttext{0.281} & \textbf{1.042}\textsuperscript{*} & \textbf{-0.049} & \textbf{0.971}\textsuperscript{*} & \textbf{-11.9\%}\textsuperscript{*} \\
% Direct & 100 & 0.873 & \$89,309 & \oldtext{\textit{source}}\latesttext{23.1\%} & \oldtext{\textit{source}}\latesttext{0.307} & \textbf{1.036}\textsuperscript{*} & \textbf{-0.029} & \textbf{0.983}\textsuperscript{*} & \textbf{-4.5\%}\textsuperscript{*} \\
% \addlinespace[3pt]
% \multicolumn{10}{l}{\textbf{Surrogate penalty}} \\
% \cmidrule(lr){1-10}
% \oldtext{Surrogate} & \oldtext{$\approx0.01$} & \oldtext{\textit{pending canonical v2 extraction}} \\
% \latesttext{Surrogate} & \latesttext{$\approx0.01$} & \latesttext{0.904} & \latesttext{\$78,620} & \latesttext{\textbf{20.7\%}\textsuperscript{*}} & \latesttext{0.279} & \latesttext{\textbf{1.077}} & \latesttext{\textbf{-0.101}} & \latesttext{\textbf{0.910}} & \latesttext{\textbf{-27.4\%}} \\
% Surrogate & \oldtext{$\approx0.1$}\latesttext{$0.1$} & \textbf{0.905}\textsuperscript{*} & \$78,546 & \oldtext{\textit{source}}\latesttext{\textbf{20.6\%}\textsuperscript{*}} & \oldtext{\textit{source}}\latesttext{0.279} & \textbf{1.074} & \textbf{-0.095} & \textbf{0.914} & \textbf{-24.7\%} \\
% Surrogate & \oldtext{$\approx0.954$}\latesttext{$\approx1$} & \textbf{0.908}\textsuperscript{*} & \textbf{\$78,196}\textsuperscript{*} & \oldtext{\textit{source}}\latesttext{\textbf{20.4\%}\textsuperscript{*}} & \oldtext{\textit{source}}\latesttext{0.284} & \textbf{1.060} & \textbf{-0.062} & \textbf{0.932} & \textbf{-15.3\%}\textsuperscript{*} \\
% Surrogate & \oldtext{$\approx10.481$}\latesttext{$\approx10$} & \textbf{0.905}\textsuperscript{*} & \$81,075 & \oldtext{\textit{source}}\latesttext{20.9\%} & \oldtext{\textit{source}}\latesttext{0.310} & \textbf{1.033}\textsuperscript{*} & \textbf{0.000}\textsuperscript{*} & \textbf{0.967}\textsuperscript{*} & \textbf{-2.0\%}\textsuperscript{*} \\
% Surrogate & 100 & 0.900 & \$85,257 & \oldtext{\textit{source}}\latesttext{22.0\%} & \oldtext{\textit{source}}\latesttext{0.336} & \textbf{1.020}\textsuperscript{*} & \textbf{0.030} & \textbf{0.987}\textsuperscript{*} & \textbf{1.7\%}\textsuperscript{*} \\
% \bottomrule
% \end{tabular}}
% \vspace{1mm}
% \begin{minipage}{\textwidth}
% \scriptsize
% \emph{Notes.}
% \oldtext{Display anchors are fixed presentation anchors and are not selected on the basis of observed performance. Positive $\rho$ values are shown at the corresponding tested grid points; Direct and Surrogate penalty magnitudes are not numerically comparable across families.}\latesttext{Display anchors are nominal decade-spaced presentation points $\{0.01,0.1,1,10,100\}$ mapped to the nearest already-tested values $\{0.0104811,0.1,0.954095,10.4811,100\}$. This presentation convention was adopted after the final augmented path analysis solely for readability and does not enter the frozen transition, LOFO, temporal-concordance, span-regret, or model-selection calculations. Direct and Surrogate penalty magnitudes are not numerically comparable across families.}
% \latesttext{Formatting comparisons are evaluated using the unrounded canonical values; entries that appear equal after display rounding can therefore receive different boldface or asterisk formatting.} Within the two baseline rows, boldface marks the baseline closer to the relevant ideal or preferred value. For Direct and Surrogate rows, boldface marks improvement over Ordinary LightGBM, and an asterisk marks improvement over both baselines. Higher $R^2_P$ and lower MAE, MAPE, and $\operatorname{RMSE}_{\log P}$ are preferred for prediction. PRD and MKI are compared by distance from one; PRB and VEI by distance from zero. Applicable assessor-facing reference ranges are PRD $[0.98,1.03]$, PRB $[-0.05,0.05]$, MKI $[0.95,1.05]$, and VEI $[-10\%,10\%]$. Boldface and asterisks do not denote statistical significance, formal compliance, or model selection. Complementary valuation-level, horizontal-uniformity, and mechanism/residual-structure values at the same anchors are reported in Appendix Table~\ref{tab:path_anchor_complementary}.
% \end{minipage}
% \end{table}
% }

{\color{latestPurple}
\begin{table}[!htbp]
\centering
\scriptsize
\renewcommand{\arraystretch}{1.08}
\caption{\latesttext{Primary regularization-path summary at nominal decade-spaced display anchors: prediction and assessor-facing vertical equity. Rows follow the fixed presentation rule $\rho\in\{0.01,0.1,1,10,100\}$ (nearest tested grid points) and do not constitute model selection. Within the two baseline rows, boldface identifies the baseline closer to the relevant ideal or preferred value. For Direct and Surrogate rows, boldface identifies point estimates that strictly improve upon Ordinary LightGBM for the corresponding metric; an asterisk identifies those that also strictly improve upon Linear Regression.}}
\label{tab:path_anchor_summary}
\resizebox{\textwidth}{!}{%
\begin{tabular}{llrrrrrrrr}
\toprule
Method & $\rho$
& $R^2_P$
& MAE
& MAPE
& $\operatorname{RMSE}_{\log P}$
& PRD
& PRB
& MKI
& VEI \\
\midrule
\multicolumn{10}{l}{\textit{Panel A: Held-out evaluation}} \\
\addlinespace[2pt]
\multicolumn{10}{l}{\textbf{Baseline models}} \\
\cmidrule(lr){1-10}
Linear regression & -- & 0.799 & \$90,092 & 24.1\% & 0.322 & \textbf{1.042} & \textbf{-0.016} & \textbf{0.975} & \textbf{-11.6\%} \\
Ordinary LightGBM & -- & \textbf{0.894} & \textbf{\$75,655} & \textbf{21.2\%} & \textbf{0.289} & 1.069 & -0.091 & 0.923 & -26.5\% \\
\addlinespace[3pt]
\multicolumn{10}{l}{\textbf{Direct penalty}} \\
\cmidrule(lr){1-10}
\latesttext{Direct} & \latesttext{$\approx0.01$} & \latesttext{0.893} & \latesttext{\$76,058} & \latesttext{21.2\%} & \latesttext{0.290} & \latesttext{1.069} & \latesttext{\textbf{-0.089}} & \latesttext{\textbf{0.924}} & \latesttext{-26.9\%} \\
Direct & \latesttext{$0.1$} & 0.894 & \$76,235 & \latesttext{21.3\%} & \latesttext{0.291} & \textbf{1.068} & \textbf{-0.088} & \textbf{0.926} & -27.0\% \\
Direct & \latesttext{$\approx1$} & \textbf{0.899}\textsuperscript{*} & \textbf{\$74,485}\textsuperscript{*} & \latesttext{\textbf{21.1\%}\textsuperscript{*}} & \latesttext{0.290} & \textbf{1.060} & \textbf{-0.075} & \textbf{0.940} & \textbf{-21.9\%} \\
Direct & \latesttext{$\approx10$} & \textbf{0.895}\textsuperscript{*} & \textbf{\$75,218}\textsuperscript{*} & \latesttext{21.4\%} & \latesttext{0.294} & \textbf{1.036}\textsuperscript{*} & \textbf{-0.040} & \textbf{0.984}\textsuperscript{*} & \textbf{-10.5\%}\textsuperscript{*} \\
Direct & 100 & 0.869 & \$84,918 & \latesttext{23.9\%} & \latesttext{0.325} & \textbf{1.029}\textsuperscript{*} & \textbf{-0.015}\textsuperscript{*} & \textbf{0.998}\textsuperscript{*} & \textbf{0.7\%}\textsuperscript{*} \\
\addlinespace[3pt]
\multicolumn{10}{l}{\textbf{Surrogate penalty}} \\
\cmidrule(lr){1-10}
\latesttext{Surrogate} & \latesttext{$\approx0.01$} & \latesttext{0.894} & \latesttext{\$75,748} & \latesttext{\textbf{21.2\%}\textsuperscript{*}} & \latesttext{0.289} & \latesttext{1.069} & \latesttext{\textbf{-0.088}} & \latesttext{\textbf{0.924}} & \latesttext{\textbf{-25.8\%}} \\
Surrogate & \latesttext{$0.1$} & 0.894 & \$75,897 & \latesttext{\textbf{21.1\%}\textsuperscript{*}} & \latesttext{0.292} & \textbf{1.064} & \textbf{-0.079} & \textbf{0.931} & \textbf{-22.7\%} \\
Surrogate & \latesttext{$\approx1$} & \textbf{0.897}\textsuperscript{*} & \textbf{\$75,468}\textsuperscript{*} & \latesttext{\textbf{21.0\%}\textsuperscript{*}} & \latesttext{0.298} & \textbf{1.049} & \textbf{-0.045} & \textbf{0.953} & \textbf{-12.6\%} \\
Surrogate & \latesttext{$\approx10$} & 0.893 & \$78,110 & \latesttext{21.6\%} & \latesttext{0.329} & \textbf{1.024}\textsuperscript{*} & \textbf{0.013}\textsuperscript{*} & \textbf{0.988}\textsuperscript{*} & \textbf{0.2\%}\textsuperscript{*} \\
Surrogate & 100 & 0.889 & \$82,094 & \latesttext{22.8\%} & \latesttext{0.357} & \textbf{1.010}\textsuperscript{*} & \textbf{0.042} & \textbf{1.013}\textsuperscript{*} & \textbf{4.0\%}\textsuperscript{*} \\
\midrule
\addlinespace[2pt]
\multicolumn{10}{l}{\textit{Panel B: 2025 forward evaluation}} \\
\addlinespace[2pt]
\multicolumn{10}{l}{\textbf{Baseline models}} \\
\cmidrule(lr){1-10}
Linear regression & -- & 0.799 & \$99,371 & 24.9\% & 0.313 & \textbf{1.052} & \textbf{-0.029} & \textbf{0.954} & \textbf{-17.3\%} \\
Ordinary LightGBM & -- & \textbf{0.904} & \textbf{\$78,484} & \textbf{20.8\%} & \textbf{0.278} & 1.079 & -0.106 & 0.907 & -28.6\% \\
\addlinespace[3pt]
\multicolumn{10}{l}{\textbf{Direct penalty}} \\
\cmidrule(lr){1-10}
\latesttext{Direct} & \latesttext{$\approx0.01$} & \latesttext{\textbf{0.905}\textsuperscript{*}} & \latesttext{\textbf{\$78,442}\textsuperscript{*}} & \latesttext{\textbf{20.7\%}\textsuperscript{*}} & \latesttext{\textbf{0.278}\textsuperscript{*}} & \latesttext{\textbf{1.078}} & \latesttext{\textbf{-0.105}} & \latesttext{\textbf{0.909}} & \latesttext{-29.8\%} \\
Direct & \latesttext{$0.1$} & \textbf{0.906}\textsuperscript{*} & \textbf{\$78,029}\textsuperscript{*} & \latesttext{\textbf{20.7\%}\textsuperscript{*}} & \latesttext{\textbf{0.278}\textsuperscript{*}} & \textbf{1.077} & \textbf{-0.102} & \textbf{0.910} & \textbf{-27.7\%} \\
Direct & \latesttext{$\approx1$} & \textbf{0.910}\textsuperscript{*} & \textbf{\$77,139}\textsuperscript{*} & \latesttext{\textbf{20.7\%}\textsuperscript{*}} & \latesttext{0.279} & \textbf{1.069} & \textbf{-0.090} & \textbf{0.926} & \textbf{-23.9\%} \\
Direct & \latesttext{$\approx10$} & 0.902 & \$79,187 & \latesttext{20.9\%} & \latesttext{0.281} & \textbf{1.042}\textsuperscript{*} & \textbf{-0.049} & \textbf{0.971}\textsuperscript{*} & \textbf{-11.9\%}\textsuperscript{*} \\
Direct & 100 & 0.873 & \$89,309 & \latesttext{23.1\%} & \latesttext{0.307} & \textbf{1.036}\textsuperscript{*} & \textbf{-0.029} & \textbf{0.983}\textsuperscript{*} & \textbf{-4.5\%}\textsuperscript{*} \\
\addlinespace[3pt]
\multicolumn{10}{l}{\textbf{Surrogate penalty}} \\
\cmidrule(lr){1-10}
\latesttext{Surrogate} & \latesttext{$\approx0.01$} & \latesttext{0.904} & \latesttext{\$78,620} & \latesttext{\textbf{20.7\%}\textsuperscript{*}} & \latesttext{0.279} & \latesttext{\textbf{1.077}} & \latesttext{\textbf{-0.101}} & \latesttext{\textbf{0.910}} & \latesttext{\textbf{-27.4\%}} \\
Surrogate & \latesttext{$0.1$} & \textbf{0.905}\textsuperscript{*} & \$78,546 & \latesttext{\textbf{20.6\%}\textsuperscript{*}} & \latesttext{0.279} & \textbf{1.074} & \textbf{-0.095} & \textbf{0.914} & \textbf{-24.7\%} \\
Surrogate & \latesttext{$\approx1$} & \textbf{0.908}\textsuperscript{*} & \textbf{\$78,196}\textsuperscript{*} & \latesttext{\textbf{20.4\%}\textsuperscript{*}} & \latesttext{0.284} & \textbf{1.060} & \textbf{-0.062} & \textbf{0.932} & \textbf{-15.3\%}\textsuperscript{*} \\
Surrogate & \latesttext{$\approx10$} & \textbf{0.905}\textsuperscript{*} & \$81,075 & \latesttext{20.9\%} & \latesttext{0.310} & \textbf{1.033}\textsuperscript{*} & \textbf{0.000}\textsuperscript{*} & \textbf{0.967}\textsuperscript{*} & \textbf{-2.0\%}\textsuperscript{*} \\
Surrogate & 100 & 0.900 & \$85,257 & \latesttext{22.0\%} & \latesttext{0.336} & \textbf{1.020}\textsuperscript{*} & \textbf{0.030} & \textbf{0.987}\textsuperscript{*} & \textbf{1.7\%}\textsuperscript{*} \\
\bottomrule
\end{tabular}}
\vspace{1mm}
\begin{minipage}{\textwidth}
\scriptsize
\emph{Notes.}
\latesttext{Display anchors are nominal decade-spaced presentation points $\{0.01,0.1,1,10,100\}$ mapped to the nearest already-tested values $\{0.0104811,0.1,0.954095,10.4811,100\}$. This presentation convention was adopted after the final augmented path analysis solely for readability and does not enter the frozen transition, LOFO, temporal-concordance, span-regret, or model-selection calculations. Direct and Surrogate penalty magnitudes are not numerically comparable across families.}
\latesttext{Formatting comparisons are evaluated using the unrounded canonical values; entries that appear tied at the displayed precision can therefore still be boldfaced or starred when the underlying canonical value is strictly better.} Within the two baseline rows, boldface marks the baseline closer to the relevant ideal or preferred value. For Direct and Surrogate rows, boldface marks improvement over Ordinary LightGBM, and an asterisk marks improvement over both baselines. Higher $R^2_P$ and lower MAE, MAPE, and $\operatorname{RMSE}_{\log P}$ are preferred for prediction. PRD and MKI are compared by distance from one; PRB and VEI by distance from zero. Applicable assessor-facing reference ranges are PRD $[0.98,1.03]$, PRB $[-0.05,0.05]$, MKI $[0.95,1.05]$, and VEI $[-10\%,10\%]$. Boldface and asterisks do not denote statistical significance, formal compliance, or model selection. Complementary valuation-level, horizontal-uniformity, and mechanism/residual-structure values at the same anchors are reported in Appendix Table~\ref{tab:path_anchor_complementary}.
\end{minipage}
\end{table}
}

\oldtext{Given the current implementation audit, within-family path changes are read relative to the custom-objective $\rho=0$ control. At the prespecified anchor near $\rho=0.954$, the Direct path raises held-out $R^2_P$ from $0.893$ to $0.899$ and lowers MAE from \$75,976 to \$74,485 while PRD, PRB, MKI, VEI, and $\beta_{\log}$ all move toward their neutral directions. The same anchor improves 2025 $R^2_P$ from $0.904$ to $0.910$ and lowers MAE from \$79,139 to \$77,139, again with movement toward neutrality in the vertical-equity diagnostics.}
\oldtext{At the prespecified anchor near $\rho=0.954$, the Direct path has held-out $R^2_P=0.899$ and MAE of \$74,485, compared with $0.894$ and \$75,655 for ordinary LightGBM, while PRD, PRB, MKI, VEI, and $\beta_{\log}$ all move toward their neutral directions. In the 2025 forward evaluation, the same anchor has $R^2_P=0.910$ and MAE of \$77,139, compared with $0.904$ and \$78,484 for ordinary LightGBM, again with movement toward neutrality in the vertical-equity diagnostics.}
\latesttext{At the nominal $\rho=1$ display anchor (tested $\rho=0.954095$), the Direct path has held-out $R^2_P=0.899$ and MAE of \$74,485, compared with $0.894$ and \$75,655 for ordinary LightGBM, while PRD, PRB, MKI, VEI, and $\beta_{\log}$ all move toward their neutral directions. In the 2025 forward evaluation, the same tested configuration has $R^2_P=0.910$ and MAE of \$77,139, compared with $0.904$ and \$78,484 for ordinary LightGBM, again with movement toward neutrality in the vertical-equity diagnostics.} The Surrogate path shows the same qualitative possibility at moderate
regularization. 
% \oldtext{Thus the empirical paths are not a simple monotone exchange of prediction accuracy for vertical equity; at stronger penalties, however, price-scale accuracy eventually deteriorates as the vertical-equity diagnostics continue to move toward neutrality.}
\newtext{Thus the empirical paths are not a simple monotone exchange of prediction accuracy for vertical equity. At stronger penalties, price-scale accuracy eventually deteriorates, while some vertical-equity diagnostics approach or cross their neutral values and others become non-monotone.}
\oldtext{The favorable moderate-penalty anchors in Table~\ref{tab:path_anchor_summary} are therefore descriptive examples, not evidence of a stable multi-metric optimum. Section~\ref{subsec:path_stability_results} evaluates that question separately using the frozen transition rule.}
\oldtext{The decade-spaced rows in Table~\ref{tab:path_anchor_summary} are therefore descriptive snapshots of different penalty scales, not evidence of a stable multi-metric optimum. Section~\ref{subsec:path_stability_results} evaluates path stability separately using the frozen transition rule.}
\latesttext{The decade-spaced rows in Table~\ref{tab:path_anchor_summary} remain descriptive snapshots of the full paths. For downstream use, the separate CV-derived screen identifies a soft activity onset and a conservative upper guardrail without selecting a model inside the surviving region. In CCAO, the Direct screen runs from $\rho=0.355648$ to $2.559548$: all five benefit diagnostics support the activity regime, and the upper bound is the conservative left edge of the first robust predictive-deterioration cluster. The Surrogate screen runs from $\rho=0.202359$ to $2.222996$: the first predictive bend occurs earlier but remains nonbinding while predictive slack remains; the later predictive-harm cluster occurs at $\rho=6.866488$, so the binding upper guardrail is instead the stable post-valley $\Delta_{\mathrm{NL}}$ rebound at $\rho=2.222996$.}

{\color{latestPurple}
\begin{table}[!htbp]
\centering
\small
\setlength{\tabcolsep}{5pt}
\renewcommand{\arraystretch}{1.08}
\caption{Chronological-CV candidate-region screen for the two CCAO penalty families.}
\label{tab:rho_candidate_regions}
\begin{tabularx}{\textwidth}{@{}lrrXl@{}}
\toprule
Family & Activity onset & Upper guardrail & Binding upper condition & LOFO reading \\
\midrule
Direct & 0.356 & 2.560 &
Earliest robust cluster of sustained predictive deterioration &
7/7 stable; guardrail indices 53--56 \\
Surrogate & 0.202 & 2.223 &
Stable post-valley $\Delta_{\mathrm{NL}}$ rebound (nonlinear-structure caution) &
7/7 rebound; guardrail indices 53--55 \\
\bottomrule
\end{tabularx}
\vspace{1mm}
\begin{minipage}{0.98\textwidth}
\scriptsize
\emph{Notes.}
The activity onset is a soft marker for the beginning of broad equity/mechanism response; smaller positive penalties are not classified as harmful or invalid. The upper guardrail is a first-stage screening boundary, not a selected, safe, optimal, or deployment penalty. Direct and Surrogate $\rho$ values are family-specific and are not numerically comparable. For Surrogate, the earlier predictive bend is nonbinding; a cross-metric predictive-harm cluster occurs later than the nonlinear caution. The nonlinear caution records a stable increase in $\Delta_{\mathrm{NL}}$ after its CV valley and does not claim that a visibly pronounced S-shape starts exactly at the guardrail.
\end{minipage}
\end{table}
}

\subsection{How the Valuation-Ratio Shape Changes Along the Path}
\label{subsec:ratio_shape_results}

\oldtext{Scalar diagnostics do not show where along the value distribution the ratio pattern changes. We therefore complement the sweep table with a selection-free ratio-shape view using the same 30 equal-count-bin construction as the baseline motivation figure and the prespecified $\rho$ display anchors within each penalty family.}
\latesttext{Scalar diagnostics do not show where along the value distribution the ratio pattern changes. We therefore complement the sweep table with a selection-free ratio-shape view using the same 30 equal-count-bin construction as the baseline motivation figure, the zero-penalty path origin, and the nominal decade-spaced positive display anchors $\{0.01,0.1,1,10,100\}$ mapped to their nearest tested grid points.}

\begin{figure}[!htbp]
\centering
\safeincludegraphics[width=0.8\textwidth]{img/generated_v12_994/ratio_shape_evolution.pdf}
\caption{\oldtext{Valuation-ratio profiles against sale price at the prespecified display anchors. The horizontal line at 1 is the principal neutrality reference; lines at 0.9 and 1.1 are aggregate appraisal-level reference guides and are not binwise acceptance criteria.}\latesttext{Valuation-ratio profiles against sale price at the zero-penalty path origin and nominal decade-spaced positive display anchors $\rho\in\{0.01,0.1,1,10,100\}$, each represented by the nearest tested grid point. The horizontal line at 1 is the principal neutrality reference; lines at 0.9 and 1.1 are aggregate appraisal-level reference guides and are not binwise acceptance criteria. The display anchors are presentation-only and do not select a penalty strength.}}
\label{fig:ratio_shape_path_placeholder}
\end{figure}

Both formulations attenuate the broad decline in valuation ratios as regularization
increases, but they change the profile differently. 
% \oldtext{The Direct path modifies the shape comparatively smoothly.} 
\oldtext{At the prespecified anchors currently displayed, the Direct profile changes more smoothly than the Surrogate profile.}\latesttext{Across the decade-spaced display snapshots, the Direct profile changes more smoothly than the Surrogate profile.} Under stronger Surrogate regularization, the profile
develops a pronounced non-monotone trough through the lower and middle sale-price
range before recovering at higher prices. \newtext{This visual contrast is descriptive rather than a structural guarantee: Direct covariance controls the affine residual--price component but does not protect against nonlinear conditional-mean departures at stronger regularization. Nonlinear shape is therefore evaluated for both families with the common $\Delta_{\mathrm{NL}}$ diagnostic rather than inferred from visual smoothness alone.} 
% \oldtext{Consequently, a near-neutral global slope or grouped high-versus-low statistic need not imply a locally flat valuation-ratio profile.}
\newtext{Consequently, a near-neutral global slope or high-versus-low grouped statistic need not imply a flat valuation-ratio profile across the value distribution.}
\oldtext{The fixed cross-fitted $\Delta_{\mathrm{NL}}$ nonlinearity gap separates these families. Along the Direct anchors, held-out $\Delta_{\mathrm{NL}}$ remains near $0.116$--$0.132$ and is $0.115$ at $\rho=100$; the 2025 Direct path is similarly stable, ending at $0.119$. The Surrogate path instead falls through moderate regularization, from $0.116$ at $\rho=0$ to $0.099$ near $\rho=0.954$ held-out and from $0.121$ to $0.092$ near $\rho=10.481$ in 2025, then rises at $\rho=100$ to $0.124$ and $0.111$, respectively. That rebound coincides with the non-monotone Surrogate ratio trough above. We do not infer a specific functional shape from the scalar gap alone; the shape label comes from the ratio-profile figure.}
\latesttext{The fixed cross-fitted $\Delta_{\mathrm{NL}}$ nonlinearity gap separates these families. Along the Direct path, held-out $\Delta_{\mathrm{NL}}$ remains near $0.116$--$0.132$ over the previously reported representative configurations and is $0.115$ at $\rho=100$; the 2025 Direct path is similarly stable, ending at $0.119$. The Surrogate path instead falls through moderate regularization, from $0.116$ at $\rho=0$ to $0.099$ near the nominal $\rho=1$ display anchor (tested $\rho=0.954095$) on the held-out sample and from $0.121$ to $0.092$ near the nominal $\rho=10$ display anchor (tested $\rho=10.4811$) in 2025, then rises at $\rho=100$ to $0.124$ and $0.111$, respectively. That rebound coincides with the non-monotone Surrogate ratio trough above. We do not infer a specific functional shape from the scalar gap alone; the shape label comes from the ratio-profile figure.}
\latesttext{The chronological-CV screening analysis adds an earlier warning to this out-of-time picture. For Surrogate, $\Delta_{\mathrm{NL}}$ reaches its CV valley before entering a leave-one-fold-out-stable rebound at the upper candidate guardrail. Event-relative ratio-profile QA does not show a visibly pronounced trough at that exact location in most folds: only 1/7 folds shows the deformation near the caution event, while 6/7 show it farther to the right. We therefore interpret the guardrail as a nonlinear-structure \emph{caution} that precedes the visibly severe high-$\rho$ shape distortion, not as the exact onset of an S-shape.}

\subsection{
\texorpdfstring{
% \oldtext{The Correction Acts on the Intended First-Order Mechanism}
\newtext{First-Order Correction and Residual Structure Along the Path}}{First-Order Correction and Residual Structure Along the Path}
}
\label{subsec:mechanism_results}

% \oldtext{The direct training target is covariance, and $\beta_{\log}$ is its signed first-order slope equivalent up to the fixed variance of log sale price within an evaluation sample. Distance correlation provides a complementary check for remaining nonlinear dependence. Figure~\ref{fig:mechanism_path_placeholder} therefore reports the complete paths rather than a single positive-$\rho$ model.}

{\color{impGreen}The Direct training target is covariance, and $\beta_{\log}$ is its signed first-order slope equivalent up to the fixed variance of log sale price within an evaluation sample. The correlation-ratio nonlinearity gap $\Delta_{\mathrm{NL}}$ then measures conditional-mean explanatory power beyond the best affine residual--price relationship, while distance correlation provides a broader omnibus check for remaining distributional dependence. These diagnostics answer nested but distinct questions---first-order trend, non-affine conditional mean, and broader dependence---and are therefore traced along the complete paths rather than evaluated only at a single positive-$\rho$ model.}

\begin{figure}[!htbp]
\centering
\safeincludegraphics[width=0.8\textwidth]{img/generated_v12_994/mechanism_vs_rho_candidate_region.pdf}
\caption{\oldtext{Mechanism and residual-structure paths versus $\rho$ for Direct and Surrogate on held-out (solid) and 2025 (dashed) evaluations. The $\beta_{\log}=0$ line is a first-order neutrality reference; $\Delta_{\mathrm{NL}}$ and dCor are not given assessor-style zero targets. Gray fill and dashed vertical boundaries mark the frozen family-specific CV-derived descriptive transition span; they are not a selected or recommended penalty interval.}\latesttext{Mechanism and residual-structure paths versus $\rho$ for Direct and Surrogate on held-out (solid) and 2025 (dashed) evaluations, with the CV-derived screening region overlaid unchanged. The light fill marks the candidate region, the dashed left boundary is the soft activity onset, and the solid right boundary is the family-specific upper guardrail. The older prediction/COD transition span is retained only as subordinate context. For Surrogate, the binding upper boundary is the nonlinear-structure caution from the stable CV $\Delta_{\mathrm{NL}}$ rebound; dCor is still improving but flattening at that event and does not set the guardrail.}}
\label{fig:mechanism_path_placeholder}
\end{figure}

The first-order mechanism moves in the intended direction in both families. Along
the held-out Direct anchors, $\beta_{\log}$ changes from $-0.147$ at
$\rho=0$ to $-0.079$ at $\rho=100$; the corresponding 2025 values move from
$-0.161$ to $-0.093$. The Surrogate path moves closer to first-order neutrality,
reaching $-0.001$ on the held-out sample and $-0.015$ in 2025 at $\rho=100$.
% \oldtext{The nonlinear dependence diagnostic is not monotone, however.} 
\newtext{The broader distance-correlation diagnostic is not monotone along the Surrogate path, however.} For the Surrogate,
held-out dCor falls from $0.382$ at $\rho=0$ to $0.250$ near
$\rho=10.481$ and then rises to $0.267$ at $\rho=100$; the 2025 path shows
the same rebound from $0.258$ to $0.266$. 
% \oldtext{This behavior is consistent with the non-monotone ratio shapes above and shows that reducing the targeted first-order association does not eliminate broader nonlinear sale-price-related dependence.} 
\newtext{This rebound is consistent with the non-monotone ratio shapes above and shows that broader sale-price-related dependence can remain even as the targeted first-order association approaches zero. Because distance correlation is an omnibus dependence measure, the rebound does not by itself establish that nonlinear conditional-mean structure has increased; that distinction is assessed with $\Delta_{\mathrm{NL}}$.}
\newtext{Figure~\ref{fig:mechanism_path_placeholder} traces that distinction on the complete grids. On Direct, $\lvert\beta_{\log}\rvert$ declines from $\rho=0$ to $\rho=100$ while $\Delta_{\mathrm{NL}}$ stays in a narrow band and is not increasing at the strongest penalty. On Surrogate, $\lvert\beta_{\log}\rvert$ continues to fall as $\Delta_{\mathrm{NL}}$ first declines and then rises at high $\rho$, including the held-out movement from $0.099$ near $\rho=0.954$ to $0.124$ at $\rho=100$. First-order slope reduction and the correlation-ratio nonlinearity gap are therefore not the same object, and the Surrogate path is the clearer CCAO example of that separation.}
\latesttext{At the Surrogate CV caution event, distance correlation does not independently reverse: its local path is classified as \emph{flattening}, meaning that it is still decreasing but at a slightly smaller rate. This distinction is important for interpretation. The upper screen is triggered by the stable $\Delta_{\mathrm{NL}}$ rebound, not by a claimed simultaneous rebound in every dependence diagnostic.}

\subsection{Accuracy--Equity Trajectories}
\label{subsec:tradeoff_results}

The central assessor-facing result is the path traced jointly by predictive
performance and vertical-equity diagnostics.
\oldtext{PRD and VEI provide complementary ratio-based and grouped views and are therefore the two main-text tradeoff panels. PRB and MKI are retained as companion diagnostics in Appendix~\ref{app:ccao_path_details}.}
\latesttext{The main tradeoff figure reports all four assessor-facing vertical-equity diagnostics---PRD, PRB, MKI, and VEI---against $R^2_P$. Appendix~\ref{app:ccao_path_details} broadens the accuracy dimension to MAE, MAPE, and $\operatorname{RMSE}_{\log P}$ and separately reports mechanism-versus-accuracy trajectories.}
% \oldtext{The previous headline accuracy--equity figure overlaid a centered-recalibration path with the Direct and Surrogate paths.}
\begin{figure}[!htbp]
\centering
\safeincludegraphics[width=0.98\textwidth]{img/generated_v12_994/accuracy_equity_trajectories_inprocessing_only.pdf}
\caption{\oldtext{Accuracy--equity trajectories with $R^2_P$ on the vertical axis and PRD, PRB, MKI, and VEI on the horizontal axis, for held-out and 2025 evaluations. Linear and ordinary LightGBM are context anchors. Arrows indicate increasing $\rho$ and do not mark a selected point.}\latesttext{Accuracy--equity trajectories with $R^2_P$ on the vertical axis and PRD, PRB, MKI, and VEI on the horizontal axis, for held-out and 2025 evaluations. Linear and ordinary LightGBM are context anchors. Dotted vertical lines mark metric neutrality where applicable, and shaded regions are assessor-facing guidance ranges rather than compliance claims. Arrows indicate increasing $\rho$ and do not mark a selected point.}}
\label{fig:accuracy_equity_placeholder}
\end{figure}
\newtext{The trajectories are not simple monotone accuracy--equity frontiers. At moderate regularization, parts of both paths move toward more neutral PRD and VEI while preserving or improving price-scale fit relative to the corresponding custom-objective $\rho=0$ control.} 
% \oldtext{At stronger regularization, the vertical-equity diagnostics continue toward neutrality while $R^2_P$ bends downward.}
\newtext{At stronger regularization, $R^2_P$ bends downward while several vertical-equity diagnostics approach or cross their neutral values; the resulting trajectories can themselves become non-monotone.} This is why the complete path, rather than a single endpoint or a presumed accuracy cost, is 
% \oldtext{the relevant pre-selection empirical object}
\newtext{the relevant empirical object before selecting an operating point}. \oldtext{Accordingly, the main accuracy--equity figure does not highlight the Direct CV transition span as a preferred segment; temporal portability of that span is assessed separately below.}\latesttext{Accordingly, the main accuracy--equity figure does not highlight either family's CV transition span as a preferred segment; temporal portability of those spans is assessed separately below.}

\subsection{Chronological Stability and Other Model-Quality Dimensions}
\label{subsec:path_stability_results}

A regularization path is more informative if its direction is stable across market
periods and if improvements in vertical-equity diagnostics do not conceal large
changes in valuation level, horizontal uniformity, or alternative prediction
metrics. \oldtext{We therefore reserve the full fold-level stability traces and the
prediction/level/dispersion paths for Appendix~\ref{app:ccao_path_details}, while
reporting the main path findings in the figures above.}
\latesttext{We therefore reserve the complete fold-level paths for prediction, valuation level and horizontal uniformity, vertical equity, and mechanism diagnostics for Appendix~\ref{app:ccao_path_details}, while reporting the main path findings in the figures above.}

\begin{oldrevisionblock}
\oldtext{The chronological folds support the direction of the targeted first-order mechanism but do not imply a stable multi-metric operating point. On Direct, $\beta_{\log}$ becomes less negative from $\rho=0$ to $\rho=100$ in all seven folds. Applying the frozen five-metric transition rule to the equal-weight chronological-CV path yields a descriptive Direct span $\rho\in[0.1,1.099]$. All seven leave-one-fold-out aggregations also yield a five-event positive-interior span, so this within-CV pattern is not driven by any single fold. The lower endpoint requires an important qualification: $\rho=0.1$ is the smallest positive penalty on the tested grid and therefore is not an estimated lower transition threshold.}
\latesttext{The chronological folds support reproducible within-development path structure but do not identify a stable multi-metric operating point. After the lower-tail extension, the five Direct CV events all occur at positive interior grid points and define the descriptive span $\rho\in[0.0494,1.099]$. The event that had previously been censored at the first positive point---$\operatorname{RMSE}_{\log P}$---moves inward to $\rho=0.0494$, while the other four Direct event locations are unchanged at the reported grid resolution. All seven Direct leave-one-fold-out aggregations again yield a valid positive-interior five-event span.}

\oldtext{Temporal concordance of the exact event locations is substantially weaker. After freezing the Direct CV span, only one of the five held-out turning events and none of the five 2025 turning events fall inside it. Value-based span regret gives a more nuanced result than event locations alone: held-out $\operatorname{RMSE}_{\log P}$ has zero span regret, whereas the other four held-out primary metrics have strictly positive regret; all five 2025 primary metrics have strictly positive regret. Because no substantive ``small-regret'' threshold is imposed, these signs are not used to label the positive regrets as material or negligible. The evidence is therefore mixed by metric: the Direct span is a reproducible description of the development-period CV path, but it does not define a temporally invariant penalty region.}
\latesttext{Exact temporal concordance is nevertheless weak. None of the five Direct held-out turning events and none of the five Direct 2025 turning events fall inside the frozen v2 CV span. Thus, the lower-tail extension strengthens the interpretation of the Direct span as an interior description of the development-period CV path while providing no evidence that the same event envelope is a temporally invariant operating region.}

\oldtext{The Surrogate does not support the same five-metric construction in the full chronological-CV aggregate because $\operatorname{RMSE}_{\log P}$ is minimized at the $\rho=0$ boundary; only two of seven leave-one-fold-out aggregations produce a common positive-interior span. Thus, moderate Surrogate regularization can still improve several predictive and assessor-facing criteria, but those improvements do not organize into the same common five-metric CV transition structure as Direct. Appendix~\ref{app:ccao_path_details} reports the event locations, leave-one-fold-out diagnostics, and Direct regret values. These are descriptive path-stability results and do not select a model family or operating point.}
\latesttext{The lower-tail extension also changes the Surrogate conclusion. Its CV $\operatorname{RMSE}_{\log P}$ event moves from the $\rho=0$ boundary to the positive interior point $\rho=0.00222$, and its MAE event moves below $0.1$. The five Surrogate CV events therefore define their own descriptive span, $\rho\in[0.00222,0.954]$, and all seven Surrogate leave-one-fold-out aggregations now support a positive-interior five-event span. This does not make the Surrogate span a selected or portable operating range: only one of five held-out events and one of five 2025 events (MAE in each case) fall inside it. Taken together, the two families show that reproducible within-CV transition structure is distinct from temporal portability of the exact operating location. Appendix~\ref{app:ccao_path_details} reports the event locations and leave-one-fold-out diagnostics. These are descriptive path-stability results and do not select a model family or operating point.}
\latesttext{The family-specific CV spans remain Direct $\rho\in[0.0494,1.099]$ and Surrogate $\rho\in[0.00222,0.954]$. They should be interpreted as envelopes of the five prediction/COD turning events, not as equity-neutrality intervals. Appendix Figures~\ref{fig:vertical_equity_event_locations} and~\ref{fig:mechanism_event_locations} show that the closest-to-neutral locations for PRD, PRB, MKI, and VEI generally occur at much larger $\rho$ for both families; predictive/COD turning and equity neutrality therefore occur at distinct scales. Exact out-of-time event concordance remains weak, but the family-symmetric value-based span-regret result in Table~\ref{tab:transition_regret} is more nuanced. Direct normalized regret is roughly $0.005$--$0.056$ on the held-out sample and $0.004$--$0.037$ in 2025. Surrogate MAE regret is $0$ on both held-out and 2025 evaluations; Surrogate $\operatorname{RMSE}_{\log P}$ normalized regret is at most $0.003$; the largest reported Surrogate normalized regret is 2025 $R^2_P\approx 0.083$.}
\oldtext{The corresponding metric-value losses are generally limited in magnitude over the observed paths, although no materiality threshold is imposed.}
\latesttext{These regret values show that exact event-location disagreement and value-based loss are distinct aspects of temporal portability; no materiality threshold is imposed. Direct $\operatorname{RMSE}_{\log P}$ and Surrogate $R^2_P$ CV events are particularly value-flat, while MAE events are comparatively sharper; exact event-$\rho$ movement therefore should not be interpreted as equal substantive movement for all criteria.}
\end{oldrevisionblock}

{\color{latestPurple}
The generic candidate-region screen is the primary path-stability summary for downstream use. Direct has activity onset $\rho=0.355648$ and upper guardrail $2.559548$. The activity onset is nearly identical across leave-one-fold-out aggregates (six of seven at the same grid index and one at the adjacent index), while the Direct guardrail remains within indices 53--56 in all seven recalculations. The upper boundary is therefore read as the conservative left edge of a reproducible predictive-deterioration transition, not as a claim that every larger $\rho$ is dominated.

Surrogate has activity onset $\rho=0.202359$ and upper guardrail $2.222996$. Its first predictive bend, at the earlier MAE/$R^2_P$ cluster, remains descriptive because those metrics have not yet jointly exhausted their $\rho=0$ predictive advantage. The first qualifying cross-metric predictive-harm cluster occurs later, at $\rho=6.866488$. The binding upper guardrail is instead the post-valley $\Delta_{\mathrm{NL}}$ rebound: all seven leave-one-fold-out aggregates detect the same qualitative rebound within indices 53--55, so the caution event is not driven by a single fold. The ratio-profile QA shows that the visually pronounced non-monotone deformation usually appears farther right, which is why the boundary is described as an early nonlinear-structure caution rather than the onset of visible pathology.

The older five-metric prediction/COD transition spans remain a separate descriptive diagnostic: Direct $\rho\in[0.0494,1.099]$ and Surrogate $\rho\in[0.00222,0.954]$. Their exact event locations transfer weakly to the later samples and they are not used as candidate-region endpoints. Held-out and 2025 paths are overlaid only after the CV screening endpoints are frozen; they provide retrospective portability evidence and do not revise the screen. The candidate regions therefore reduce the positive-$\rho$ search space while intentionally leaving the final within-region operating-point choice unresolved.
}

\subsection{Exploratory Six-County Evidence Shows Heterogeneity and Limits}
\label{subsec:attom_results}

% \oldtext{Table~\ref{tab:attom_baselines} shows the unpenalized ATTOM benchmark results. Every county has at least one vertical-equity point estimate outside the reference bands used in the benchmark, but the patterns are not identical. Cook, Allegheny, Maricopa, King, and Middlesex have negative PRB point estimates outside the benchmark band, whereas Miami-Dade's PRB point estimate lies within the band even though its PRD remains high. Miami-Dade is especially useful as a boundary case: its positive PRB and curved ratio pattern show why a global covariance target and a single ratio statistic cannot be treated as complete descriptions of vertical equity.}
\newtext{Table~\ref{tab:attom_baselines} shows the unpenalized ATTOM benchmark results. Every county has at least one vertical-equity point estimate outside the benchmark bands, but the patterns differ materially across jurisdictions. Cook, Allegheny, Maricopa, King, and Middlesex have negative PRB point estimates outside the benchmark band, whereas Miami-Dade's PRB lies within the band even though its PRD remains high. Miami-Dade is therefore a useful boundary case: its positive PRB and curved ratio profile illustrate why neither the global covariance target nor any single ratio statistic provides a complete description of vertical equity.}

{\color{impGreen}
\begin{table}[t]
\centering
\caption{Unpenalized LightGBM results in the six-county ATTOM benchmark. These are held-out model-benchmark point estimates, not official assessor results.}
\label{tab:attom_baselines}
\resizebox{\textwidth}{!}{%
\begin{tabular}{lrrrrrrr}
\toprule
County & $R^2_{\log P}$ & $\operatorname{RMSE}_{\log P}$ & MAPE (\%) & COD & PRD & PRB & Median $r$ \\
\midrule
Cook, IL       & 0.833 & 0.282 & 20.9 & 21.266 & 1.073 & -0.079 & 0.976 \\
Allegheny, PA  & 0.781 & 0.324 & 25.3 & 26.061 & 1.103 & -0.103 & 0.956 \\
Maricopa, AZ   & 0.848 & 0.219 & 14.0 & 14.171 & 1.076 & -0.060 & 0.977 \\
King, WA       & 0.814 & 0.250 & 16.7 & 16.790 & 1.102 & -0.074 & 0.948 \\
Miami-Dade, FL & 0.815 & 0.310 & 25.4 & 24.181 & 1.086 &  0.027 & 1.039 \\
Middlesex, MA  & 0.790 & 0.258 & 18.6 & 18.255 & 1.083 & -0.088 & 1.014 \\
\bottomrule
\end{tabular}}
\end{table}
}
\todo{Table~\ref{tab:attom_baselines} has been consolidated to the intended revised values. Before submission, verify these values against the same frozen county-report artifact used to generate Tables~\ref{tab:attom_selected_penalty} and~\ref{tab:attom_design_appendix}; regenerate all three together if the benchmark changes.}

\newtext{The exploratory benchmark reports $R^2_{\log P}$ because its historical comparison is defined on the log-price scale; the primary CCAO application instead uses price-scale $R^2_P$ as its headline coefficient of determination. MAPE is reported in percentage points, consistently with Section~\ref{subsec:metrics}.}

% \oldtext{Among the four counties whose baseline PRB point estimate is outside the benchmark band, at least one tested penalty places PRB inside the band in all four, with no more than a 1.84\% increase in log-scale RMSE. In contrast, only Maricopa has a tested penalty that places both PRD and PRB point estimates inside their bands.}
\newtext{For an auditable held-out comparison, Table~\ref{tab:attom_selected_penalty} reports the single baseline-penalty configuration chosen for each county on its earlier chronological validation block. The held-out block is not used in that choice. The selected penalty moves PRB inside $[-0.05,0.05]$ in Cook, Allegheny, Maricopa, King, and Middlesex, while Miami-Dade moves from $0.027$ to $0.067$. PRD remains above $1.03$ for all six selected configurations. The corresponding held-out log-RMSE changes range from $-0.20\%$ in King County to $+5.33\%$ in Middlesex County. These descriptive results show that the first-order correction can materially alter value-related bias diagnostics across heterogeneous datasets, but neither the assessor-facing response nor the predictive cost is uniform across counties.}

{\color{impGreen}
\begin{table}[!htbp]
\centering
\scriptsize
\setlength{\tabcolsep}{4.0pt}
\renewcommand{\arraystretch}{1.06}
\caption{Validation-selected covariance-penalty comparison in the exploratory six-county ATTOM benchmark. Selection is performed on the earlier chronological validation block; the reported values are from the held-out block.}
\label{tab:attom_selected_penalty}
\begin{tabular}{lrrrrrrr}
\toprule
County & Selected $\rho$ & PRD$_0$ & PRD$_{\mathrm{pen}}$ & PRB$_0$ & PRB$_{\mathrm{pen}}$ & $\Delta$RMSE$_{\log P}$ (\%) \\
\midrule
Cook, IL & 10.39 & 1.073 & 1.036 & -0.079 & -0.023 & +2.11 \\
Allegheny, PA & 22.21 & 1.103 & 1.047 & -0.103 & -0.004 & +4.16 \\
Maricopa, AZ & 14.47 & 1.076 & 1.055 & -0.060 & -0.020 & +1.04 \\
King, WA & 13.64 & 1.102 & 1.079 & -0.074 & -0.023 & -0.20 \\
Miami-Dade, FL & 4.03 & 1.086 & 1.046 & 0.027 & 0.067 & +3.77 \\
Middlesex, MA & 28.50 & 1.083 & 1.051 & -0.088 & -0.009 & +5.33 \\
\bottomrule
\end{tabular}
\vspace{1mm}
\begin{minipage}{0.98\textwidth}
\scriptsize
\emph{Notes.} Subscript $0$ denotes ordinary LightGBM and ``pen'' denotes the county-specific covariance-penalty configuration selected on chronological validation by the benchmark's Pareto--Nash maximum-volume-hyperrectangle rule over MAPE, COD, $|\mathrm{PRD}-1|$, and $|\mathrm{PRB}|$. The held-out block is evaluated only after the validation-stage choice. The benchmark remains exploratory because the broader county configurations were developed iteratively rather than under one preregistered confirmatory protocol.
\end{minipage}
\end{table}
}

% \oldtext{These results do not establish a universal or confirmatory compliance result and do not imply that the selected penalty dominates simpler post-hoc recalibration; the latter comparison is discussed next.}
\newtext{These results do not establish a universal or confirmatory compliance result, nor do they establish that training-time regularization dominates other possible correction strategies.}

% \oldtext{The previous draft included a separate primary CCAO centered-recalibration results subsection, endpoint table, and overlays in the accuracy--equity figures. Those empirical comparisons are removed from the present paper so that the contribution is scoped to training-time covariance-guided correction.}
\todo{Follow-up paper: evaluate the complete post-hoc calibration frontier---beginning with the one-parameter centered-spread map suggested by the Direct fixed-space result and extending to structured geographic/temporal corrections---against Direct and Surrogate under matched temporal validation, accuracy, and vertical-equity targets.}

\section{Discussion and Conclusions}
\label{sec:discussion}

% \oldtext{The CCAO application answers a limited question. The regenerated baseline comparison establishes a clear application-specific tension: relative to linear regression, ordinary LightGBM improves predictive accuracy and aggregate ratio dispersion while exhibiting a stronger regressive vertical-equity pattern. The primary pre-selection Results therefore report the complete direct and surrogate regularization paths---their ratio-shape changes, mechanism diagnostics, accuracy--equity trajectories, and chronological stability---rather than a single chosen penalty. Claims about a preferred penalty strength or final operating point are intentionally deferred to a separate model-selection analysis.}
\newtext{The CCAO application provides a clear answer to the paper's central model-design question. The covariance-guided objectives can materially attenuate the regressive first-order pattern of the LightGBM baseline, and moderate parts of the regularization paths can do so without an immediate loss of predictive performance. At the same time, the full paths show that stronger global correction can alter valuation level, ratio dispersion, and nonlinear ratio structure. The method therefore provides a practical training-time control for one important component of regressivity, not a complete solution to vertical equity. The present analysis reports the complete paths; selection of a deployment-oriented operating point is deferred.}

The complete paths sharpen that conclusion in two ways. First, moderate
regularization need not impose an immediate accuracy cost: parts of both
custom-objective paths improve price-scale predictive performance while moving
multiple vertical-equity diagnostics toward neutrality. 
% \oldtext{Second, stronger regularization exposes the limits of a global first-order target. The Surrogate can drive $\beta_{\log}$ and grouped diagnostics close to neutral while the ratio-shape profile remains non-monotone and distance correlation rebounds. Global first-order neutrality is therefore informative but not sufficient evidence of locally uniform vertical equity.}
\newtext{Second, stronger regularization exposes the limits of a global first-order target. The Surrogate provides the clearest current CCAO example: $\beta_{\log}$ and grouped diagnostics can approach neutrality while the ratio-shape profile remains non-monotone and distance correlation rebounds. The methodological limitation is broader than that example, however: neither Direct nor Surrogate directly controls the non-affine conditional-mean structure measured by $\Delta_{\mathrm{NL}}$. Global first-order neutrality is therefore informative but not sufficient evidence of a flat value-related profile, broader residual--price independence, or locally uniform vertical equity.}

\oldtext{The transition analysis adds a separate temporal qualification. A targeted mechanism can be directionally stable without yielding a stable multi-metric operating point. Direct exhibits an internally reproducible five-metric CV transition span, including seven-of-seven leave-one-fold-out support, but the exact held-out and forward turning locations show limited concordance with that interval and the value-based regret result is mixed across metrics. The Surrogate does not support the analogous common positive span in the full CV aggregate. Chronological validation is therefore valuable for characterizing the regularization path, but the evidence does not support treating a particular raw-$\rho$ interval as an intrinsic or temporally invariant property of either objective. Any later deployment calibration should specify explicitly which predictive, valuation, and equity criteria define the operating decision and re-evaluate that decision across time.}
\oldtext{The transition analysis adds a separate temporal qualification. After resolving the lower tail of the grid, both Direct and Surrogate exhibit positive-interior five-metric CV transition spans with seven-of-seven leave-one-fold-out support. Exact out-of-time event portability remains weak, however: Direct has zero of five held-out and zero of five 2025 events inside its frozen span, while Surrogate has one of five in each period. Chronological validation is therefore useful for characterizing repeatable path structure, but the evidence does not support treating either raw-$\rho$ interval as an intrinsic or temporally invariant operating region. Any later deployment calibration should specify explicitly which predictive, valuation, and equity criteria define the operating decision and re-evaluate that decision across time.}
\oldtext{The transition analysis adds a separate temporal qualification. After resolving the lower tail of the grid, both Direct and Surrogate exhibit positive-interior five-metric CV transition spans with seven-of-seven leave-one-fold-out support. Exact out-of-time event portability remains weak, however: Direct has zero of five held-out and zero of five 2025 events inside its frozen span, while Surrogate has one of five in each period. The augmented evidence distinguishes a reproducible prediction/COD transition structure from the much larger penalties at which assessor-facing measures or mechanism diagnostics become closest to neutrality. Therefore there is no single empirically supported ``safe'' $\rho$ interval, and a future deployment choice requires an explicit operating criterion. Chronological validation is therefore useful for characterizing repeatable path structure, but the evidence does not support treating either raw-$\rho$ interval as an intrinsic or temporally invariant operating region. Any later deployment calibration should specify explicitly which predictive, valuation, and equity criteria define the operating decision and re-evaluate that decision across time.}
\latesttext{The CV screening analysis adds a more operational, but still deliberately limited, path summary. In CCAO it identifies a Direct positive-$\rho$ candidate region of approximately $0.36$--$2.56$ and a Surrogate region of approximately $0.20$--$2.22$. The lower endpoints are soft activity-onset markers. The Direct upper endpoint is the conservative onset of a stable predictive-deterioration transition; the Surrogate upper endpoint is a stable post-valley rebound in $\Delta_{\mathrm{NL}}$, interpreted as a nonlinear-structure caution that precedes the visibly severe high-$\rho$ ratio-shape deformation. These ranges are screening devices rather than operating-point choices: they are derived from chronological CV, leave the within-region tradeoff unresolved, and are not asserted to be intrinsic or temporally invariant properties of either objective. The older prediction/COD transition spans remain useful descriptive diagnostics but no longer serve as the primary candidate-range interpretation.}

% \oldtext{The value of the approach is operational. It converts a ratio-study concern into a model-training option, works through a standard boosted-tree interface, and can be tuned with chronological validation.}
\newtext{The value of the approach is operational. It converts a ratio-study concern into a model-training option that can be incorporated into a standard boosted-tree workflow: the Surrogate is exactly observation-additive, while the Direct implementation uses the exact covariance gradient with diagonal curvature through the standard LightGBM interface. Both can be evaluated through the same chronological-validation and ratio-study workflow used for the baseline model.} \advisorchange{Those features follow directly from the CCAO problem that motivated the collaboration.}

% \oldtext{The previous version mixed the scope limitations, the native/custom implementation caveat, ATTOM external-validity caveats, and the centered-recalibration comparison in one repeated paragraph.}

% \oldtext{The evidence also bounds the paper's claim. Better data can improve predictive accuracy and assessment equity together. The proposed penalties target a countywide first-order value-related pattern; they do not determine valuation level, guarantee assessor-facing reference ranges, remove nonlinear or local inequity, or establish broader residual--price independence. The ATTOM benchmark provides exploratory evidence about transferability and failure modes rather than confirmatory external validation. Centered recalibration provides a meaningful low-cost comparator: it can reproduce a large first-order correction, while the primary CCAO results show that moderate in-processing configurations can move the same broad family of diagnostics with substantially smaller predictive losses than the validation-neutral recalibration endpoint. The covariance-guided correction should therefore be treated as one auditable modeling option within the broader valuation-review process, not as an equity certificate.}
\newtext{The evidence also bounds the paper's claim. Better data can improve predictive accuracy and assessment equity together. The proposed penalties target a countywide first-order value-related pattern; they do not determine valuation level, guarantee assessor-facing reference ranges, remove nonlinear or local inequity, or establish broader residual--price independence. The ATTOM benchmark provides exploratory evidence about transferability and failure modes rather than confirmatory external validation. The covariance-guided correction should therefore be treated as one auditable modeling option within the broader valuation-review process, not as evidence that training-time regularization is necessary or sufficient for vertical equity.}

\newtext{The fixed-prediction-space analysis sharpens this interpretation rather than creating a separate theoretical claim. The Direct objective has rank-one geometry and, with an intercept in the fixed space, reduces exactly to a centered spread correction; this makes transparent both what the covariance penalty controls and how limited that control can be. The Surrogate, by contrast, is a log-price-distance-weighted projection whose correction can combine multiple learnable modes and therefore need not trace the same path as Direct. This comparison explains why the Surrogate should not be interpreted as a computational approximation to the Direct objective alone. At the same time, neither fixed-space result predicts the exact path of a retrained tree ensemble, and the spectral representation does not by itself explain the empirical high-$\rho$ ratio shape.}

\subsection{\texorpdfstring{
% \oldtext{Assessor Implementation}
\newtext{Practical Implementation Considerations}}{Practical Implementation Considerations}}
\label{sec:implementation_workflow}

% \oldtext{The method can be inserted into an existing AVM with the following sequence.}
\newtext{For a later operational implementation, the correction can be incorporated into an existing AVM development workflow through the following sequence.}

\begin{enumerate}[leftmargin=7mm]
    \item Build a defensible baseline model using verified sales, appropriate time adjustment, relevant property features, and local market information.
    \item Reserve future observations for held-out and forward evaluation and create rolling or blocked validation folds that respect sale chronology.
    \item Report valuation level, COD, PRD, PRB, MKI, VEI, value-group ratio profiles, and the mechanism/residual-structure diagnostics $\beta_{\log}$, $\Delta_{\mathrm{NL}}$, and dCor for the baseline model.
    \item Fit the Direct and Surrogate covariance-guided objectives over prespecified logarithmic grids of $\rho$, including $\rho=0$ as an implementation control.
    \item Summarize the complete regularization paths across chronological validation folds, including fold-level variability.
    \item \oldtext{If a later operating point is required, define its objective and guardrails explicitly and assess temporal portability; do not treat clustering of CV turning points as a universal penalty-selection rule.}
    \item \latesttext{Apply the generic CV screening rule to identify a soft activity onset and a family-specific upper guardrail. This step narrows the positive-$\rho$ configurations passed forward but does not rank the surviving models or select a deployment penalty.}
    \item \latesttext{If a later operating point is required, choose it from the surviving candidate set using an explicit predictive/equity decision rule and assess that decision's temporal and subgroup portability.}
    \item Refit the same fixed paths on the full development samples and evaluate them descriptively on the held-out and forward periods without changing the grid or base learner.
    \item Repeat the diagnostics by geography, property class, value range, and time before any later deployment-oriented model-selection decision.
\end{enumerate}

\begin{table}[t]
\centering
\caption{Different pipeline components address different problems.}
\label{tab:roles}
\begin{tabular}{p{4.2cm}p{10.0cm}}
\toprule
Component & Primary role \\
\midrule
Better property, spatial, and temporal data & Reduce omitted structure and improve the conditional valuation model; may improve accuracy and equity together \\
Constant level calibration & Move the overall ratio level toward its required target; does not change log-residual covariance \\
% \oldtext{One-parameter slope/spread rescaling} & \oldtext{Low-cost post-hoc correction of a global value trend; benchmark for the penalty} \\
Covariance-guided regularization & Discourage a global first-order log-ratio/log-price trend during model training \\
% \oldtext{Stratified and local diagnostics} & \oldtext{Detect offsetting geographic, class-specific, or nonlinear patterns hidden by aggregate measures} \\
% \oldtext{$\Delta_{\mathrm{NL}}$, ratio-shape, and stratified/local diagnostics} & \oldtext{Separate nonlinear conditional-mean departures from the first-order trend and detect offsetting geographic or class-specific patterns hidden by aggregate measures} \\
\newtext{Residual-structure and local diagnostics} & \newtext{Use $\beta_{\log}$, $\Delta_{\mathrm{NL}}$, distance correlation, and ratio-shape diagnostics to distinguish first-order, non-affine conditional-mean, and broader dependence patterns; use stratified diagnostics to detect geographic or property-class patterns hidden by countywide summaries.} \\
\latesttext{Temporal validation and uncertainty analysis} & \latesttext{Characterize how path behavior changes across market periods; by itself, chronological validation does not establish a temporally invariant penalty range} \\
\bottomrule
\end{tabular}
\end{table}

The office should save the penalty definition, full grid, validation predictions, held-out and forward predictions, and all ratio-study outputs. This documentation makes the change auditable and allows the correction to be re-evaluated in later assessment cycles.

\subsection{Limitations}

% \oldtext{First, the CCAO application uses sold properties. Results on sales do not establish performance for the full unsold inventory, particularly when the sale sample is not representative. Second, the primary application is a Python translation of the CCAO workflow rather than an official production run. Third, the Direct penalty targets log-scale covariance while the Surrogate penalizes its sample-additive upper bound; PRD, PRB, and VEI are price-scale statistics with different functional forms. Relatedly, neither Direct nor Surrogate directly minimizes the nonlinear lack-of-fit $\Delta_{\mathrm{NL}}$; strong regularization can therefore reduce the targeted first-order slope without guaranteeing a flatter conditional-mean residual--price relationship. Fourth, the current pre-selection analysis does not yet identify a deployment-oriented penalty strength; that decision is intentionally deferred. Fifth, fold-level variation and the reported group-level uncertainty should accompany any later operational decision. Sixth, a countywide correction can average away local regressivity and progressivity. Seventh, the ATTOM benchmark is exploratory, uses commercially harmonized data, and does not reproduce the six offices' institutional data or a single confirmatory model-selection protocol. Finally, post-hoc centered recalibration provides an important benchmark.}

\newtext{Three groups of limitations delimit the interpretation of the results. First, the evidence is sales-based and application-specific. The primary CCAO analysis uses sold properties and a Python translation of the published residential workflow rather than an official production run; performance on observed sales does not establish performance on the full unsold inventory, particularly under sale-selection differences. The six-county ATTOM benchmark is likewise exploratory and uses commercially harmonized data rather than each assessor office's institutional data and production protocol.}

% \oldtext{Second, the correction is deliberately narrow. Direct targets a global log-scale covariance, while the Surrogate penalizes a sample-additive upper bound with a distinct weighted-loss mechanism. Neither objective directly optimizes PRD, PRB, VEI, valuation level, or the correlation-ratio nonlinearity gap $\Delta_{\mathrm{NL}}$. A countywide first-order correction can therefore coexist with nonlinear, geographic, property-class, or temporal inequity.}
\newtext{Second, the correction is deliberately narrow. Direct targets a global log-scale covariance, while the Surrogate penalizes a sample-additive upper bound with a distinct weighted-loss mechanism. Neither objective directly optimizes PRD, PRB, VEI, valuation level, or the correlation-ratio nonlinearity gap $\Delta_{\mathrm{NL}}$. Moreover, because observed sale price is both the prediction target and the value benchmark, the fixed-space theory shows that part of the in-sample residual--price covariance can arise mechanically under squared-error fitting. A countywide first-order correction can therefore coexist with nonlinear, geographic, property-class, or temporal inequity and should not be interpreted as a direct estimate of inequity relative to latent market value.}

% \oldtext{Third, the present analysis characterizes regularization paths rather than a deployment decision. Fold-level heterogeneity and group-level uncertainty should accompany any later operating-point choice, and simple centered recalibration remains an important benchmark because it can reproduce a substantial first-order correction. The present evidence supports comparing the complete in-processing and recalibration paths rather than claiming that either approach uniformly dominates the other.}
\newtext{Third, the present analysis characterizes regularization paths rather than a deployment decision. Fold-level heterogeneity and group-level uncertainty should accompany any later operating-point choice. The paper also does not establish comparative superiority over post-processing alternatives, which are outside its empirical scope.} \oldtext{The transition analysis sharpens this limitation empirically: even the Direct family, for which the five-metric CV construction survives every leave-one-fold-out recalculation, does not reproduce the same exact turning locations on the later samples, and its value-based portability is mixed across metrics.}\latesttext{The transition analysis sharpens this limitation empirically: both families' five-metric CV constructions survive every leave-one-fold-out recalculation on the augmented grid, yet their exact turning locations transfer weakly to the held-out and 2025 paths.}
\latesttext{The new candidate-region screen should be interpreted with the same restraint. Its leave-one-fold-out stability supports a repeatable development-period screening rule, not a confidence interval or a guarantee that the same raw-$\rho$ endpoints will remain appropriate in a later assessment cycle. Held-out and 2025 overlays are retrospective checks and are not used to move the CV-derived boundaries.}

\subsection{Future Work}

\advisorchange{The next evaluation step is a clean ablation} that compares the baseline model, better spatial and temporal information, the covariance correction, and their combination under the same temporal splits. Comparable-property and neighborhood features must use only information available before the target sale. 
\latesttext{For any future deployment-oriented calibration, the selection objective and guardrails should be frozen before evaluating later assessment cycles, and temporal stability should be reported in both event-location and value-regret terms rather than inferred from a single CV turning region.}

% \oldtext{Preliminary tests of second- and third-degree global recalibration have shown little consistent improvement over a first-degree map, so higher-order corrections should remain secondary unless they add stable out-of-time gains.}

\oldtext{The main methodological extension is structured covariance control by geography, property class, and time, with safeguards for small or unstable groups. That extension should be judged against stratified post-hoc recalibration and against a stronger base AVM.}
\newtext{One natural methodological extension is structured covariance control by geography, property class, and time, with safeguards for small or unstable groups and separate regularization strengths for the dimensions being targeted. A distinct follow-up question is when comparable first-order improvements can be achieved by post-hoc calibration of an already fitted AVM. The fixed-space Direct result makes that comparison natural, but the present paper does not characterize the in-processing/post-processing Pareto frontier. A follow-up study should compare global and structured post-hoc corrections with corresponding training-time penalties under the same temporal splits and the same predictive, assessor-facing, nonlinear-shape, and local-equity diagnostics.} Further work should also add block-bootstrap uncertainty, worst-period validation, and evaluation on the full assessment inventory. Desk review by assessors is needed to determine whether the changed predictions remain understandable and operationally defensible.




\printbibliography


\appendix

\section{Metric Details and Guidance Status}
\label{app:metrics}

\subsection{Assessor-facing guidance status}

The 2013 IAAO Standard on Ratio Studies is the established standard cited for PRD
and COD guidance in this paper \cite{IAAO2013Ratio}. The May 2026 revision of the
ratio-studies document is an exposure draft. Measures introduced or revised there,
including VEI, are reported as descriptive assessor-facing guidance rather than as
final adopted compliance standards \cite{IAAO2026ExposureRatio}. Reference bands
are used as descriptive guidance in the present
% \oldtext{pre-selection}
\newtext{regularization-path} analysis; the paper does not
infer legal compliance from a point estimate.

The weighted mean ratio satisfies
\[
\bar r_W
=
\frac{\sum_i r_iP_i}{\sum_iP_i}
=
\frac{\sum_i\widehat P_i}{\sum_iP_i},
\]
and PRD is $\bar r/\bar r_W$. PRB follows the IAAO median-normalized regression
definition. The full implementation \advisorchange{used in the application should
be released} with the code so that all transformations, trimming rules, and proxy
definitions are reproducible.

\subsection{Mechanism and residual-structure diagnostics}

The three diagnostics introduced after the assessor-facing measures play distinct
roles. The sample slope $\beta_{\log}(S)$ in Eq.~\eqref{eq:log_ratio_slope} is the
direct out-of-sample analogue of the covariance mechanism: on a fixed evaluation
sample it differs from empirical covariance only by the fixed factor
$\widehat{\Var}_S(y)^{-1}$.

For the nonlinear diagnostic, let
\[
Z=\frac{y-\E[y]}{\sqrt{\Var(y)}},
\qquad
m(z)=\E[e\mid Z=z],
\]
and let $\Pi_{\mathrm{aff}}m$ denote the $L^2(P_Z)$ projection of $m$ onto the
affine span of $\{1,Z\}$. Pearson's correlation ratio, also called the
nonparametric coefficient of determination, is
\[
\eta^2(e\mid Z)
=
\frac{\Var\!\left(m(Z)\right)}{\Var(e)}
=
\frac{\Var\!\left(\E[e\mid Z]\right)}{\Var(e)}
\]
\cite{DoksumSamarov1995}. Because $Z$ is standardized, the affine component explains
the fraction $\Corr(e,Z)^2$ of the variance of $e$. Orthogonal projection therefore
gives
\begin{equation}
\label{eq:nonlinearity_gap_projection}
\Delta_{\mathrm{NL}}
=
\eta^2(e\mid Z)-\Corr(e,Z)^2
=
\frac{
\E\!\left[
\left\{
m(Z)-(\Pi_{\mathrm{aff}}m)(Z)
\right\}^2
\right]
}{
\Var(e)
}
\geq 0.
\end{equation}
Thus $\Delta_{\mathrm{NL}}$ isolates the conditional-mean explanatory power that is
not captured by the best affine residual--price relationship. We use
\emph{correlation-ratio nonlinearity gap} as a descriptive name and
$\Delta_{\mathrm{NL}}$ as paper-specific shorthand; neither the name nor the symbol
is claimed as canonical notation from \textcite{DoksumSamarov1995}. That paper
provides the correlation-ratio/nonparametric-$R^2$ foundation and explicitly notes
its use in constructing measures of regression nonlinearity. The integrated
squared-difference goodness-of-fit framework of \textcite{HardleMammen1993} is
closely related to the projection representation above, but it is not the source of
the exact normalized statistic in Eq.~\eqref{eq:nonlinearity_gap_projection}.

For the empirical CCAO analysis, the finite-sample implementation is fixed before
the penalized paths are compared. Each held-out or forward evaluation sample $S$ is
assigned deterministically to five folds using sale identifiers. For each fold, we
fit on the complementary observations (i) an affine regression of $e$ on $Z$ and
(ii) a cubic $B$-spline regression of $e$ on $Z$ with eight quantile knots, and we
score both fits on the omitted fold. Let
$\operatorname{MSE}_{\mathrm{aff}}(S)$ and
$\operatorname{MSE}_{\mathrm{spl}}(S)$ denote the resulting pooled out-of-fold mean
squared errors. The reported sample statistic is
\begin{equation}
\label{eq:nonlinearity_gap_estimator}
\widehat{\Delta}_{\mathrm{NL}}(S)
=
\max\left\{
0,\,
\frac{
\operatorname{MSE}_{\mathrm{aff}}(S)
-
\operatorname{MSE}_{\mathrm{spl}}(S)
}{
\widehat{\Var}_S(e)
}
\right\}.
\end{equation}
The nonnegative truncation respects the population inequality in
Eq.~\eqref{eq:nonlinearity_gap_projection}; finite-sample cross-fitting can otherwise
produce a small negative difference when the flexible fit does not improve
out-of-fold squared error. For readability, the empirical tables use
$\Delta_{\mathrm{NL}}$ as shorthand for
$\widehat{\Delta}_{\mathrm{NL}}(S)$. \oldtext{The estimator is reported on the primary
held-out and 2025 forward paths and is not computed retrospectively for the already
completed chronological CV.}
\latesttext{The estimator is reported on the primary held-out and 2025 forward paths. The same frozen cross-fitted estimator was later applied, fold by fold, to the already-stored chronological validation predictions; those CV values are a retrospective diagnostic and were not used to define the frozen prediction/COD transition span, the $\rho$ grid, or a model-selection rule.}

Finally, population distance correlation is zero if and only if independence under
the usual moment conditions \cite{SzekelyRizzoBakirov2007}. The paper uses its
sample analogue only as an omnibus diagnostic of residual--price dependence; unlike
$\beta_{\log}$ it is unsigned, and unlike $\Delta_{\mathrm{NL}}$ it is not restricted
to conditional-mean structure. It has no assessor-standard reference range.

\todo{Confirm and archive the exact finite-sample distance-correlation implementation
used in the executed metric code (including whether the reported estimator is the
standard biased/V-statistic or an unbiased/U-centered variant) so the replication
materials match the manuscript exactly.}

\section{Squared-Error Prediction and Residual--Value Dependence}
\label{app:population_identity}

Let $Y=\log P$, let $X$ be the observed feature vector, and let
\[
f^\star(X)=\E[Y\mid X]
\]
be the squared-error Bayes predictor. Define $e^\star=f^\star(X)-Y$. If $\E[Y^2]<\infty$, then
\begin{proposition}
\label{prop:bayes_covariance}
\[
\Cov(e^\star,Y)
=
-\E[\Var(Y\mid X)]
\le 0.
\]
\end{proposition}
\begin{proof}
Using the law of total variance,
\[
\Cov(e^\star,Y)
=
\Cov(\E[Y\mid X]-Y,Y)
=
\Var(\E[Y\mid X])-\Var(Y)
=
-\E[\Var(Y\mid X)].
\]
\end{proof}

This identity explains a source of negative log-residual dependence when important price variation remains unresolved after conditioning on $X$. It does not imply that every finite-sample model is regressive under every assessor metric. It also does not imply that adding features cannot improve both accuracy and equity.

\section{Covariance and PRD on the Price Scale}
\label{app:cov_prd}

Let $P>0$ and $r=\widehat P/P>0$. The population PRD can be written as
\[
\mathrm{PRD}
=
\frac{\E[r]\E[P]}{\E[rP]}.
\]
Then
\begin{proposition}
\label{prop:cov_prd}
\[
\operatorname{sign}\!\left(\Cov(r,P)\right)
=
\operatorname{sign}(1-\mathrm{PRD}).
\]
\end{proposition}
\begin{proof}
Since $\E[rP]>0$,
\[
1-\mathrm{PRD}
=
\frac{\E[rP]-\E[r]\E[P]}{\E[rP]}
=
\frac{\Cov(r,P)}{\E[rP]}.
\]
\end{proof}

The penalty in the paper acts on $\Cov(\log r,\log P)$, not on $\Cov(r,P)$. Proposition~\ref{prop:cov_prd} therefore gives directional motivation rather than an equality between the training target and PRD.

\section{\texorpdfstring{\oldtext{Exact Fixed-Space Regularization Path}\newtext{Fixed-Space Regularization Paths: Direct and Surrogate}}{Fixed-Space Regularization Paths: Direct and Surrogate}}
\label{app:theory}

\oldtext{This appendix isolates the geometry of the two penalties when the prediction space is held fixed. The Direct result provides an exact rank-one benchmark for the targeted covariance correction. The Surrogate result below provides the corresponding weighted-projection benchmark and makes precise why its path need not remain one-dimensional. These finite-sample projection results are used for mechanism interpretation only; they do not assert that a fully retrained boosting path remains in a fixed prediction space.}
\newtext{This appendix isolates the data-fit geometry of the two penalties when the prediction space is held fixed, omitting model-specific complexity penalties in order to expose the correction mechanisms cleanly. The Direct result provides an exact rank-one benchmark for the targeted covariance correction. The Surrogate result provides the corresponding log-price-distance-weighted projection benchmark and makes precise why its path need not remain one-dimensional. These finite-sample projection results are used for mechanism interpretation only; they do not assert that a fully retrained boosting path remains in a fixed prediction space or that the fixed-space formulas describe LightGBM's finite boosting trajectory exactly.}

\subsection{\newtext{Direct penalty: rank-one covariance path}}

Let $\Vcal\subseteq\R^n$ be a closed linear prediction space, define
\[
\langle u,v\rangle_n=\frac{1}{n}u^\top v,
\qquad
\|u\|_n^2=\langle u,u\rangle_n,
\qquad
c=y-\bar y\1,
\]
and consider
\begin{equation}
\label{eq:path_problem}
f_\rho\in\argminop_{f\in\Vcal}
\left\{
\|f-y\|_n^2
+\frac{\rho}{2}\langle f-y,c\rangle_n^2
\right\}.
\end{equation}
Let $P_{\Vcal}$ be the orthogonal projection under $\langle\cdot,\cdot\rangle_n$ and set
\[
f_0=P_{\Vcal}y,\qquad
C_0=\langle f_0-y,c\rangle_n,\qquad
g=P_{\Vcal}c,\qquad
A=\|g\|_n^2.
\]

\begin{proposition}[Covariance shrinkage path]
\label{prop:path}
If $A>0$, then
\[
f_\rho=f_0-\frac{\rho}{2}C_\rho g,
\qquad
C_\rho=\frac{C_0}{1+\rho A/2}.
\]
Consequently, the slope of residuals on $y$ is multiplied by
\[
q(\rho)=\frac{1}{1+\rho A/2}.
\]
\end{proposition}

\begin{proof}
For every $h\in\Vcal$, first-order optimality gives
\[
2\langle f_\rho-y,h\rangle_n
+\rho C_\rho\langle c,h\rangle_n=0.
\]
Since $f_0=P_{\Vcal}y$, $\langle f_0-y,h\rangle_n=0$. Writing
$d_\rho=f_\rho-f_0\in\Vcal$ and $g=P_{\Vcal}c$ gives
\[
2d_\rho+\rho C_\rho g=0.
\]
Taking the inner product with $c$ yields
\[
C_\rho=C_0-\frac{\rho}{2}AC_\rho,
\]
which proves the result.
\end{proof}

{\color{impGreen}
\begin{corollary}[Centered-spread form of the fixed-space Direct path]
\label{cor:path_scaling}
If $\1\in\Vcal$, then $g=f_0-\bar y\1$ and
\[
f_\rho
=
\bar y\1
+
b_\rho\bigl(f_0-\bar y\1\bigr),
\qquad
b_\rho
=
1-\frac{\rho}{2}C_\rho.
\]
Moreover, $C_0=-\|f_0-y\|_n^2$. Hence, whenever the unpenalized fit is non-perfect, $C_0<0$, $b_\rho>1$ for $\rho>0$, and $C_\rho$ approaches zero from below without changing sign at any finite $\rho$.
\end{corollary}

\begin{proof}
Because $\1\in\Vcal$,
\[
g=P_{\Vcal}(y-\bar y\1)=f_0-\bar y\1.
\]
Substituting this identity into Proposition~\ref{prop:path} gives the centered-spread representation. Also,
\[
c=(f_0-\bar y\1)+(y-f_0),
\]
and the projection residual $y-f_0$ is orthogonal to $\Vcal$, so
\[
C_0
=
\langle f_0-y,c\rangle_n
=
-\|f_0-y\|_n^2.
\]
The remaining statements follow from $C_\rho=C_0/(1+\rho A/2)$.
\end{proof}

\begin{remark}[Mechanical component of the in-sample residual--outcome covariance]
\label{rem:mechanical_covariance}
\newtext{Corollary~\ref{cor:path_scaling} implies that, in this fixed-space squared-error benchmark with an intercept, $C_0$ is nonpositive even without imposing any equity interpretation: $C_0=-\|f_0-y\|_n^2$. The negative training covariance therefore contains a mechanical residual--outcome component induced by fitting the same observed $y$ that is later used to define the log valuation ratio. This identity does not hold mechanically on a new evaluation sample, and it does not imply that observed out-of-sample regressivity is spurious. Rather, it limits the interpretation of the training penalty: $C(f)$ is a controllable first-order mechanism, not a standalone estimator of inequity relative to latent market value. This is why the empirical analysis evaluates the fitted paths chronologically with proxy-aware assessor measures, ratio profiles, and residual-structure diagnostics.}
\end{remark}

\oldtext{This equivalence is specific to the exact fixed-prediction-space benchmark and motivates the centered recalibration comparator in the previous draft.}
\newtext{This equivalence is specific to the exact fixed-prediction-space benchmark. It reveals that, within that benchmark, Direct covariance regularization is no richer than a one-parameter centered spread correction. This is both a mechanism result and a limitation: whether an explicit post-hoc correction can reproduce the accuracy--equity path of a fully retrained nonlinear learner is a separate empirical question left for follow-up work.} It does not imply that the empirical Direct-LightGBM path is one-dimensional: retraining can change tree splits and leaf structure, and the standard-library Direct implementation uses only the diagonal of the covariance Hessian.
\todo{Future research: use the centered-spread equivalence as the starting point for a separate matched-frontier study of post-hoc versus in-processing correction; this comparison is intentionally not part of the present empirical results.}
}

\begin{proposition}[Accuracy cost]
\label{prop:path_cost}
Under Proposition~\ref{prop:path},
\[
\|f_\rho-y\|_n^2-\|f_0-y\|_n^2
=
\frac{C_0^2}{A}\bigl(1-q(\rho)\bigr)^2.
\]
\end{proposition}

Within a fixed prediction space, the targeted covariance falls to first order near $\rho=0$, while the in-sample squared-error cost is second order. For boosted trees this is a local benchmark only. Retraining can change splits, leaves, and stopping iteration.

\subsection{\newtext{Surrogate penalty: weighted-projection path}}

{\color{impGreen}
Let
\[
D=\diag(c_1^2,\ldots,c_n^2),
\qquad
W_\rho=I+\rho D,
\qquad
e_0=f_0-y,
\]
where $f_0=P_{\Vcal}y$ is the ordinary least-squares projection defined above. The exact fixed-space Surrogate problem corresponding to Eq.~\eqref{eq:surrogate_weighted_loss} is
\begin{equation}
\label{eq:surrogate_fixed_space_problem}
f_\rho^{\mathrm{surr}}
=
\argminop_{f\in\Vcal}
\left\{
\|f-y\|_n^2
+
\rho\langle f-y,D(f-y)\rangle_n
\right\}.
\end{equation}
Because $W_\rho$ is positive definite for every $\rho\geq0$, this problem has a unique solution. Define the weighted inner product
\[
\langle u,v\rangle_{n,W_\rho}
=
\langle u,W_\rho v\rangle_n,
\]
let $P_{\Vcal}^{W_\rho}$ denote projection onto $\Vcal$ under this inner product, and define the self-adjoint positive-semidefinite operator
\[
T=P_{\Vcal}D\big|_{\Vcal}
\]
together with the baseline tail-weighted residual direction
\[
g_{\mathrm{surr}}
=
P_{\Vcal}De_0.
\]}

{\color{impGreen}
\begin{proposition}[Surrogate weighted-projection path]
\label{prop:surrogate_fixed_space}
For every $\rho\geq0$,
\begin{equation}
\label{eq:surrogate_weighted_projection}
f_\rho^{\mathrm{surr}}
=
P_{\Vcal}^{W_\rho}y,
\end{equation}
and its departure from the unpenalized fit is
\begin{equation}
\label{eq:surrogate_correction_operator}
f_\rho^{\mathrm{surr}}-f_0
=
-\rho
\left(I_{\Vcal}+\rho T\right)^{-1}
g_{\mathrm{surr}}.
\end{equation}
\end{proposition}

\begin{proof}
For every $h\in\Vcal$, first-order optimality for~\eqref{eq:surrogate_fixed_space_problem} gives
\[
\langle f_\rho^{\mathrm{surr}}-y,h\rangle_n
+
\rho\langle D(f_\rho^{\mathrm{surr}}-y),h\rangle_n
=
0,
\]
or equivalently
\[
\langle f_\rho^{\mathrm{surr}}-y,h\rangle_{n,W_\rho}=0.
\]
This is exactly the weighted-projection characterization in~\eqref{eq:surrogate_weighted_projection}. Now write
\[
\delta_\rho=f_\rho^{\mathrm{surr}}-f_0\in\Vcal.
\]
Since $f_0=P_{\Vcal}y$, the ordinary projection residual $e_0=f_0-y$ is orthogonal to $\Vcal$. Substituting
$f_\rho^{\mathrm{surr}}-y=e_0+\delta_\rho$ into the first-order condition and projecting the $D$ terms onto $\Vcal$ yields
\[
\delta_\rho
+
\rho T\delta_\rho
+
\rho g_{\mathrm{surr}}
=
0.
\]
Because $T$ is positive semidefinite, $I_{\Vcal}+\rho T$ is invertible for $\rho\geq0$, which proves~\eqref{eq:surrogate_correction_operator}.
\end{proof}
}

{\color{impGreen}
\begin{corollary}[Local Surrogate direction and baseline-loss cost]
\label{cor:surrogate_local}
As $\rho\downarrow0$,
\begin{equation}
\label{eq:surrogate_local_direction}
f_\rho^{\mathrm{surr}}
=
f_0
-
\rho g_{\mathrm{surr}}
+
O(\rho^2)
=
f_0
-
\rho P_{\Vcal}
\left[
c^{\odot2}\odot e_0
\right]
+
O(\rho^2).
\end{equation}
where $\odot$ denotes elementwise multiplication and $c^{\odot2}=(c_1^2,\ldots,c_n^2)^\top$. Moreover,
\begin{equation}
\label{eq:surrogate_baseline_loss_cost}
\|f_\rho^{\mathrm{surr}}-y\|_n^2
-
\|f_0-y\|_n^2
=
\rho^2\|g_{\mathrm{surr}}\|_n^2
+
O(\rho^3).
\end{equation}
\end{corollary}

\begin{proof}
The expansion
\[
(I_{\Vcal}+\rho T)^{-1}
=
I_{\Vcal}-\rho T+O(\rho^2)
\]
in~\eqref{eq:surrogate_correction_operator} gives~\eqref{eq:surrogate_local_direction}. For the second claim, $\delta_\rho=f_\rho^{\mathrm{surr}}-f_0$ belongs to $\Vcal$ and $e_0=f_0-y$ is orthogonal to $\Vcal$, so
\[
\|e_0+\delta_\rho\|_n^2-\|e_0\|_n^2
=
\|\delta_\rho\|_n^2.
\]
Using $\delta_\rho=-\rho g_{\mathrm{surr}}+O(\rho^2)$ gives~\eqref{eq:surrogate_baseline_loss_cost}.
\end{proof}
}

\subsection{\newtext{Direct--Surrogate comparison in a fixed prediction space}}

{\color{impGreen}
The preceding results expose a structural difference that is not visible from the upper-bound relation alone. For the Direct penalty, Proposition~\ref{prop:path} shows that the entire fixed-space path moves along the single direction $P_{\Vcal}c$; when $\1\in\Vcal$, Corollary~\ref{cor:path_scaling} identifies this direction with the centered fitted-value spread $f_0-\bar y\1$. The Direct correction is therefore rank one in this benchmark. For the Surrogate, the initial direction is instead
\[
-g_{\mathrm{surr}}
=
-P_{\Vcal}
\left[c^{\odot2}\odot e_0\right],
\]
the learnable component of the baseline residual pattern after weighting observations by squared log-price distance from the sample center. The complete Surrogate path generally changes direction as $\rho$ varies.}

{\color{impGreen}
To see this explicitly, let $\{u_j\}$ be an orthonormal eigenbasis of $T$ on $\Vcal$, with
\[
Tu_j=\lambda_j u_j,
\qquad
\lambda_j\geq0,
\]
and write
\[
g_{\mathrm{surr}}=\sum_j a_j u_j.
\]
Then Proposition~\ref{prop:surrogate_fixed_space} gives
\begin{equation}
\label{eq:surrogate_spectral_path}
f_\rho^{\mathrm{surr}}-f_0
=
-\sum_j
\frac{\rho}{1+\rho\lambda_j}
a_j u_j.
\end{equation}
Thus, unless the nonzero components of $g_{\mathrm{surr}}$ lie in a single eigenspace of $T$, the relative contribution of the learnable correction modes changes with $\rho$. This contrasts with the one-direction Direct path. It also makes clear that the Surrogate does not inherit the Direct formula
$C_\rho=C_0/(1+\rho A/2)$: the Surrogate minimizes a log-price-distance-weighted prediction loss rather than the signed covariance itself, so monotone covariance shrinkage is not guaranteed by the fixed-space theory. Likewise, the spectral representation does not imply that stronger Surrogate regularization must generate any particular nonlinear ratio shape; the empirical shape and $\Delta_{\mathrm{NL}}$ results remain separate diagnostics.}

\todo{Treat Eq.~\eqref{eq:surrogate_spectral_path} as a mechanism interpretation only. Do not attribute the empirical high-$\rho$ S-shape to particular eigenmodes unless a separate diagnostic explicitly projects the fitted-path changes onto these fixed-space modes.}

{\color{impGreen}
\oldtext{Taken together, the fixed-space results give a parallel interpretation of the two methods. Direct regularization applies a rank-one correction aligned with the centered price direction and, with an intercept in the prediction space, becomes an exact centered spread rescaling. Surrogate regularization changes the projection geometry through the diagonal weights $1+\rho c_i^2$: locally it repairs the learnable tail-weighted residual pattern, and globally it can combine multiple correction modes. In both cases the baseline squared-error cost is second order near $\rho=0$, but only the Direct objective directly guarantees shrinkage of the targeted covariance within this benchmark.}
\newtext{Taken together, the fixed-space results clarify the distinct mechanisms of the two methods. Direct regularization applies a rank-one correction aligned with the projected centered-price direction and, with an intercept in the prediction space, becomes an exact centered spread rescaling. Surrogate regularization changes the projection geometry through the diagonal weights $1+\rho c_i^2$: locally it repairs the learnable tail-weighted residual pattern, and globally it can combine multiple correction modes. In both cases the baseline squared-error cost is second order near $\rho=0$, but only the Direct objective guarantees monotone shrinkage of the targeted covariance in this benchmark. The comparison therefore supports the paper's interpretation of the Surrogate as a related but distinct objective rather than as a purely computational approximation to Direct.}
}

\section{Linear and Boosting Implementations}
\label{app:implementation}


\subsection{Linear Model}

Let $X\in\R^{n\times p}$ and $e=X\theta-y$. With
\[
c=y-\bar y\1,\qquad a=X^\top c,\qquad b=y^\top c,
\]
the direct penalized estimator under loss $\|X\theta-y\|_2^2/n$ satisfies
\[
\left(
\frac{1}{n}X^\top X+\frac{\rho}{2n^2}aa^\top
\right)\theta
=
\frac{1}{n}X^\top y+\frac{\rho}{2n^2}ba.
\]
This is a rank-one modification of the normal equations.

For the surrogate, let $D=\diag(c_1^2,\ldots,c_n^2)$. Then
\[
\left(
\frac{1}{n}X^\top X+\frac{\rho}{n}X^\top DX
\right)\theta
=
\frac{1}{n}X^\top y+\frac{\rho}{n}X^\top Dy.
\]

\subsection{LightGBM Custom-Objective Initialization and Scaling}

For the canonical custom objectives, let $b=\bar y$ denote the mean log sale price
in the current training sample. We fit the custom-objective model to centered labels
$\widetilde y_i=y_i-b$ with zero initial score and add $b$ back to predictions.
Equivalently, the fitted model predicts on the original log-price scale as
$f(x)=b+\widetilde f(x)$. 
% \oldtext{This makes the custom-objective starting prediction match the native squared-error LightGBM initialization while leaving $e_i=f(x_i)-y_i$ and $c_i=y_i-\bar y$ unchanged.}
\newtext{This construction is designed to match the native squared-error LightGBM starting prediction while leaving $e_i=f(x_i)-y_i$ and $c_i=y_i-\bar y$ unchanged.}
\oldtext{The empirical $\rho=0$ audit shows that this initialization does not produce prediction parity with native LightGBM beyond the initial score. Direct and Surrogate $\rho=0$ predictions coincide with each other, and the remaining native/custom gap is the custom-objective API path rather than a positive-$\rho$ objective error.}
\newtext{The corresponding native/custom zero-penalty audit is reported in the next subsection.}

The mathematical objectives in the main text are written as sample averages. In
the LightGBM custom-objective implementation we multiply the complete objective by
$n/2$, a common positive factor that does not change its minimizer. This yields a
native squared-error baseline gradient $g_i=e_i$ and Hessian $h_i=1$ at $\rho=0$.
For the direct objective, the supplied gradient and diagonal Hessian are
\[
g_i^{\mathrm{cov}}
=
e_i+\frac{\rho}{2}z c_i,
\qquad
h_i^{\mathrm{cov,diag}}
=
1+\frac{\rho}{2n}c_i^2,
\qquad
z=\frac{1}{n}\sum_{i=1}^n e_i c_i.
\]
For the surrogate, the supplied derivatives are
\[
g_i^{\mathrm{surr}}=e_i(1+\rho c_i^2),
\qquad
h_i^{\mathrm{surr}}=1+\rho c_i^2.
\]
Thus the direct implementation uses the exact covariance gradient and only the
diagonal of its dense covariance curvature, whereas the surrogate derivatives are
exactly observation-additive. The $\rho=0$ custom-objective path is retained as an
implementation control.

\subsection{\newtext{Zero-Penalty Implementation Audit}}
\label{subsec:zero_control_results}

\newtext{Before attributing within-family changes to positive regularization, we compare ordinary LightGBM with the Direct and Surrogate implementations at $\rho=0$. The two custom-objective controls coincide with each other in the reported metrics, but the currently frozen artifacts are not numerically identical to ordinary LightGBM. Because the supplied zero-penalty gradient and Hessian match squared-error derivatives algebraically, this audit identifies a native-versus-custom training-path difference rather than a positive-penalty objective effect; it does not by itself isolate a unique software cause.}

{\color{impGreen}
\begin{table}[!ht]
\centering
\caption{Implementation-control comparison at $\rho=0$. Prediction differences are log-price deviations from ordinary LightGBM, aligned on sale identifiers and computed independently on the held-out and 2025 samples.}
\label{tab:rho_zero_control}
\begin{tabular}{llrrrrr}
\toprule
Sample & Model & $R^2_P$ & $\operatorname{RMSE}_{\log P}$ & $\beta_{\log}$ & Mean $|\Delta\widehat y|$ & Max $|\Delta\widehat y|$ \\
\midrule
Held-out & Ordinary LightGBM & 0.894 & 0.289 & -0.150 & -- & -- \\
Held-out & Direct, $\rho=0$ & 0.893 & 0.289 & -0.147 & 3.24e-02 & 3.09e-01 \\
Held-out & Surrogate, $\rho=0$ & 0.893 & 0.289 & -0.147 & 3.24e-02 & 3.09e-01 \\
\addlinespace
2025 forward & Ordinary LightGBM & 0.904 & 0.278 & -0.164 & -- & -- \\
2025 forward & Direct, $\rho=0$ & 0.904 & 0.279 & -0.161 & 3.06e-02 & 2.83e-01 \\
2025 forward & Surrogate, $\rho=0$ & 0.904 & 0.279 & -0.161 & 3.06e-02 & 2.83e-01 \\
\bottomrule
\end{tabular}
\end{table}
}

\todo{A later initialization-aligned parity experiment in the project achieved near-numerical native/custom agreement on the tested path. Before submission, rerun that canonical parity audit for both held-out and 2025 refits. If parity holds on both splits, regenerate Table~\ref{tab:rho_zero_control} and any affected path-origin artifacts from that implementation. Keep the audit in the appendix rather than reintroducing it as a main-text Results subsection.}

\todo[inline]{\textcolor{impGreen}{Before freezing the final empirical artifacts, audit the exact executed code against the normalization in Eqs.~\eqref{eq:baseline_learning}, \eqref{eq:direct_penalty}, and \eqref{eq:surrogate_learning} and against the $n/2$ LightGBM scaling stated here. Record the exact executed code/configuration hash and archive the commit or complete working-tree diff that generated the reported artifacts. Because the empirical paths are indexed by numerical values of $\rho$, any factor-of-two or sample-size convention must be reconciled explicitly rather than absorbed silently into a relabeled $\rho$.}}

\subsection{Direct Penalty in Boosting}

For current predictions $f$, define $e=f-y$ and
\[
z=\frac{1}{n}e^\top c.
\]
Under squared loss, the direct objective has
\[
g_i=\frac{2e_i}{n}+\frac{\rho}{n}zc_i,
\]
and exact Hessian
\[
\nabla_f^2
=
\frac{2}{n}I+\frac{\rho}{n^2}cc^\top.
\]
The diagonal approximation uses
\[
h_i^{\mathrm{diag}}
=
\frac{2}{n}+\frac{\rho}{n^2}c_i^2.
\]
For a fixed tree with leaves $I_1,\ldots,I_J$, let
\[
b_j=\sum_{i\in I_j}e_i,
\qquad
s_j=\sum_{i\in I_j}c_i,
\qquad
n_j=|I_j|.
\]
With leaf-weight vector $w$, ridge parameter $\lambda$, and $s=(s_1,\ldots,s_J)^\top$, the exact local quadratic system is
\[
A_Tw^\star=-g_T,
\qquad
g_{T,j}=\frac{2b_j}{n}+\frac{\rho z s_j}{n},
\]
\[
A_T
=
\diag\!\left(\frac{2n_1}{n}+\lambda,\ldots,\frac{2n_J}{n}+\lambda\right)
+\frac{\rho}{n^2}ss^\top.
\]
Thus $w^\star=-A_T^{-1}g_T$ when $A_T$ is nonsingular. This solver retains the cross-observation curvature but requires a custom tree-construction procedure; it is not available through the usual observation-wise API.

\subsection{Surrogate in Boosting}

For
\[
\Psi^{\mathrm{surr}}(f)
=
\frac{1}{n}\sum_i e_i^2c_i^2,
\]
the derivatives are observation-wise:
\[
g_i=\frac{2e_i}{n}+\frac{2\rho}{n}e_ic_i^2,
\qquad
h_i=\frac{2}{n}+\frac{2\rho}{n}c_i^2.
\]
This is the canonical surrogate implementation evaluated in the CCAO regularization path.


\section{\texorpdfstring{
% \oldtext{Primary CCAO Pre-Selection Path Details} 
\newtext{Primary CCAO Regularization-Path Details}}{Primary CCAO Regularization-Path Details}}
\label{app:ccao_path_details}

\oldtext{This appendix contains the complete diagnostics underlying the compact main-text
regularization-path presentation. It is intentionally organized around the full
prespecified paths and not around a selected configuration.}
\oldtext{This appendix contains the complete diagnostics underlying the compact main-text regularization-path presentation. It is organized around the final augmented paths and not around a selected configuration; only the compact display anchors remain prespecified presentation choices.}\latesttext{This appendix contains the complete diagnostics underlying the compact main-text regularization-path presentation. It is organized around the final augmented paths and not around a selected configuration; the compact nominal decade-spaced anchors are presentation-only snapshots of those frozen paths.}

\subsection{Complete Regularization-Path Results}

% \oldtext{The complete pre-selection path is provided}
\newtext{The complete regularization path is provided} in the machine-readable replication
artifacts \texttt{combined\_path\_table.csv} and
\texttt{combined\_path\_table.parquet}. These files contain every family--$\rho$
configuration together with the seven fold values, equal-weight CV mean and
standard deviation, held-out results, and 2025 forward results for the
Section~\ref{subsec:metrics} evaluation metrics. \oldtext{The compact main-text table uses only the prespecified positive display anchors together with the ordinary-LightGBM baseline; zero-penalty custom-objective diagnostics are kept in the implementation appendix. No path row is omitted from the machine-readable results because of its observed performance.}
\latesttext{The compact main-text table uses only the nominal decade-spaced positive display anchors $\{0.01,0.1,1,10,100\}$, mapped to the nearest frozen-grid points, together with the ordinary-LightGBM baseline; zero-penalty custom-objective diagnostics are kept in the implementation appendix. No path row is omitted from the machine-readable results because of its observed performance.}
\latesttext{For completeness, the complementary baseline diagnostics and the complementary representative-$\rho$ diagnostics are reported here rather than in the main empirical narrative.}

{\color{latestPurple}
\begin{table}[!htbp]
\centering
\scriptsize
\setlength{\tabcolsep}{2.4pt}
\renewcommand{\arraystretch}{1.06}
\newcommand{\baselinecompmetric}[1]{\hspace*{0.6em}#1}
\caption{Complementary baseline comparison: valuation level, horizontal uniformity, and mechanism/residual structure.}
\label{tab:ccao_baseline_complementary}
\begin{tabularx}{\textwidth}{@{} >{\raggedright\arraybackslash}p{2.80cm} >{\centering\arraybackslash}X >{\centering\arraybackslash}X >{\centering\arraybackslash}X >{\centering\arraybackslash}X @{}}
\toprule
\textbf{Measure} & \multicolumn{2}{c}{\textbf{Held-out evaluation}} & \multicolumn{2}{c}{\textbf{2025 forward evaluation}} \\
\cmidrule(lr){2-3}\cmidrule(l){4-5}
& \textbf{Linear} & \textbf{\shortstack{Light\\GBM}} & \textbf{Linear} & \textbf{\shortstack{Light\\GBM}} \\
\midrule
\multicolumn{5}{@{}l}{\textbf{Valuation level}} \\
\cmidrule(r){1-1}
\baselinecompmetric{Median ratio $m_S$} & \textbf{0.969} & 0.929 & \textbf{1.020} & 0.950 \\
\baselinecompmetric{Mean ratio $\bar r_S$} & 1.020 & \textbf{0.989} & 1.081 & \textbf{1.015} \\
\baselinecompmetric{Weighted mean $\bar r_{W,S}$} & \textbf{0.979} & 0.924 & \textbf{1.027} & 0.941 \\
\addlinespace[5pt]
\multicolumn{5}{@{}l}{\textbf{Horizontal uniformity}} \\
\cmidrule(r){1-1}
\baselinecompmetric{COD} & 24.7\% & \textbf{21.6\%} & 24.3\% & \textbf{21.3\%} \\
\baselinecompmetric{COV} & 45.2\% & \textbf{39.7\%} & 42.1\% & \textbf{37.2\%} \\
\addlinespace[5pt]
\multicolumn{5}{@{}l}{\textbf{Mechanism and residual structure}} \\
\cmidrule(r){1-1}
\baselinecompmetric{$\beta_{\log}$} & \textbf{-0.092} & -0.150 & \textbf{-0.109} & -0.164 \\
\baselinecompmetric{$\Delta_{\mathrm{NL}}$} & 0.131 & \textbf{0.119} & 0.122 & \textbf{0.121} \\
\baselinecompmetric{$\operatorname{dCor}(e,y)$} & \textbf{0.250} & 0.387 & \textbf{0.269} & 0.422 \\
\bottomrule
\end{tabularx}
\vspace{1mm}
\begin{minipage}{\textwidth}
\scriptsize
\emph{Notes.}
Median, mean, weighted-mean ratio, COD, and COV are assessor-facing valuation-level or horizontal-uniformity diagnostics defined in Section~\ref{subsec:metrics}. The first-order mechanism slope $\beta_{\log}$, correlation-ratio nonlinearity gap $\Delta_{\mathrm{NL}}$, and distance correlation are defined in Eqs.~\eqref{eq:log_ratio_slope}, \eqref{eq:nonlinearity_gap}, and~\eqref{eq:dcor_diagnostic}. For the three mechanism/residual-structure diagnostics, boldface follows the direction toward less first-order, non-affine conditional-mean, and broader residual--price structure, respectively; it should not be read as an assessor-standard equity ranking. Boldface does not denote statistical significance or formal compliance.
\end{minipage}
\end{table}
}

% \begin{table}[!htbp]
% \centering
% \scriptsize
% \renewcommand{\arraystretch}{1.08}
% \caption{\oldtext{Complementary regularization-path summary at the same prespecified display anchors: valuation level, horizontal uniformity, and mechanism/residual structure.}\latesttext{Complementary regularization-path summary at the same nominal decade-spaced display anchors: valuation level, horizontal uniformity, and mechanism/residual structure.}}
% \label{tab:path_anchor_complementary}
% \resizebox{\textwidth}{!}{%
% \begin{tabular}{llrrrrrrrr}
% \toprule
% Method & $\rho$
% & Median $r$
% & Mean $r$
% & W. mean $r$
% & COD
% & COV
% & $\beta_{\log}$
% & $\Delta_{\mathrm{NL}}$
% & dCor \\
% \midrule
% \multicolumn{10}{l}{\textit{Panel A: Held-out evaluation}} \\
% \addlinespace[2pt]
% \multicolumn{10}{l}{\textbf{Baseline models}} \\
% \cmidrule(lr){1-10}
% Linear regression & -- & \textbf{0.969} & 1.020 & \textbf{0.979} & 24.7\% & 45.2\% & \textbf{-0.092} & 0.131 & \textbf{0.250} \\
% Ordinary LightGBM & -- & 0.929 & \textbf{0.989} & 0.924 & \textbf{21.6\%} & \textbf{39.7\%} & -0.150 & \textbf{0.119} & 0.387 \\
% \addlinespace[3pt]
% \multicolumn{10}{l}{\textbf{Direct penalty}} \\
% \cmidrule(lr){1-10}
% \oldtext{Direct} & \oldtext{$\approx0.01$} & \oldtext{\textit{pending canonical v2 extraction}} \\
% \latesttext{Direct} & \latesttext{$\approx0.01$} & \latesttext{0.926} & \latesttext{0.986} & \latesttext{0.922} & \latesttext{21.7\%} & \latesttext{39.9\%} & \latesttext{\textbf{-0.148}} & \latesttext{0.119} & \latesttext{\textbf{0.385}} \\
% Direct & \oldtext{$\approx0.1$}\latesttext{$0.1$} & \oldtext{\textit{source}}\latesttext{0.924} & \oldtext{\textit{source}}\latesttext{0.983} & \oldtext{\textit{source}}\latesttext{0.920} & \oldtext{\textit{source}}\latesttext{21.6\%} & \oldtext{\textit{source}}\latesttext{39.8\%} & \textbf{-0.147} & 0.119 & \textbf{0.382} \\
% Direct & \oldtext{$\approx0.954$}\latesttext{$\approx1$} & \oldtext{\textit{source}}\latesttext{\textbf{0.936}} & \oldtext{\textit{source}}\latesttext{\textbf{0.993}\textsuperscript{*}} & \oldtext{\textit{source}}\latesttext{\textbf{0.936}} & \oldtext{\textit{source}}\latesttext{21.6\%} & \oldtext{\textit{source}}\latesttext{39.9\%} & \textbf{-0.134} & 0.122 & \textbf{0.359} \\
% Direct & \oldtext{$\approx10.481$}\latesttext{$\approx10$} & \oldtext{\textit{source}}\latesttext{\textbf{0.953}} & \oldtext{\textit{source}}\latesttext{\textbf{1.005}\textsuperscript{*}} & \oldtext{\textit{source}}\latesttext{\textbf{0.970}} & \oldtext{\textit{source}}\latesttext{21.9\%} & \oldtext{\textit{source}}\latesttext{\textbf{39.4\%}\textsuperscript{*}} & \textbf{-0.096} & 0.132 & \textbf{0.302} \\
% Direct & 100 & \oldtext{\textit{source}}\latesttext{\textbf{0.968}} & \oldtext{\textit{source}}\latesttext{1.022} & \oldtext{\textit{source}}\latesttext{\textbf{0.993}\textsuperscript{*}} & \oldtext{\textit{source}}\latesttext{24.5\%} & \oldtext{\textit{source}}\latesttext{41.5\%} & \textbf{-0.079}\textsuperscript{*} & \textbf{0.115}\textsuperscript{*} & \textbf{0.265} \\
% \addlinespace[3pt]
% \multicolumn{10}{l}{\textbf{Surrogate penalty}} \\
% \cmidrule(lr){1-10}
% \oldtext{Surrogate} & \oldtext{$\approx0.01$} & \oldtext{\textit{pending canonical v2 extraction}} \\
% \latesttext{Surrogate} & \latesttext{$\approx0.01$} & \latesttext{0.929} & \latesttext{0.987} & \latesttext{0.924} & \latesttext{\textbf{21.6\%}\textsuperscript{*}} & \latesttext{39.9\%} & \latesttext{\textbf{-0.147}} & \latesttext{\textbf{0.116}\textsuperscript{*}} & \latesttext{\textbf{0.383}} \\
% Surrogate & \oldtext{$\approx0.1$}\latesttext{$0.1$} & \oldtext{\textit{source}}\latesttext{0.927} & \oldtext{\textit{source}}\latesttext{0.982} & \oldtext{\textit{source}}\latesttext{0.923} & \oldtext{\textit{source}}\latesttext{\textbf{21.5\%}\textsuperscript{*}} & \oldtext{\textit{source}}\latesttext{39.9\%} & \textbf{-0.138} & \textbf{0.113}\textsuperscript{*} & \textbf{0.367} \\
% Surrogate & \oldtext{$\approx0.954$}\latesttext{$\approx1$} & \oldtext{\textit{source}}\latesttext{0.929} & \oldtext{\textit{source}}\latesttext{0.975} & \oldtext{\textit{source}}\latesttext{\textbf{0.930}} & \oldtext{\textit{source}}\latesttext{\textbf{21.4\%}\textsuperscript{*}} & \oldtext{\textit{source}}\latesttext{\textbf{38.2\%}\textsuperscript{*}} & \textbf{-0.102} & \textbf{0.099}\textsuperscript{*} & \textbf{0.310} \\
% Surrogate & \oldtext{$\approx10.481$}\latesttext{$\approx10$} & \oldtext{\textit{source}}\latesttext{0.927} & \oldtext{\textit{source}}\latesttext{0.953} & \oldtext{\textit{source}}\latesttext{\textbf{0.931}} & \oldtext{\textit{source}}\latesttext{22.2\%} & \oldtext{\textit{source}}\latesttext{\textbf{38.5\%}\textsuperscript{*}} & \textbf{-0.036}\textsuperscript{*} & \textbf{0.103}\textsuperscript{*} & \textbf{0.250} \\
% Surrogate & 100 & \oldtext{\textit{source}}\latesttext{0.923} & \oldtext{\textit{source}}\latesttext{0.937} & \oldtext{\textit{source}}\latesttext{\textbf{0.928}} & \oldtext{\textit{source}}\latesttext{23.5\%} & \oldtext{\textit{source}}\latesttext{40.3\%} & \textbf{-0.001}\textsuperscript{*} & 0.124 & \textbf{0.267} \\
% \midrule
% \addlinespace[2pt]
% \multicolumn{10}{l}{\textit{Panel B: 2025 forward evaluation}} \\
% \addlinespace[2pt]
% \multicolumn{10}{l}{\textbf{Baseline models}} \\
% \cmidrule(lr){1-10}
% Linear regression & -- & \textbf{1.020} & 1.081 & \textbf{1.027} & 24.3\% & 42.1\% & \textbf{-0.109} & 0.122 & \textbf{0.269} \\
% Ordinary LightGBM & -- & 0.950 & \textbf{1.015} & 0.941 & \textbf{21.3\%} & \textbf{37.2\%} & -0.164 & \textbf{0.121} & 0.422 \\
% \addlinespace[3pt]
% \multicolumn{10}{l}{\textbf{Direct penalty}} \\
% \cmidrule(lr){1-10}
% \oldtext{Direct} & \oldtext{$\approx0.01$} & \oldtext{\textit{pending canonical v2 extraction}} \\
% \latesttext{Direct} & \latesttext{$\approx0.01$} & \latesttext{0.948} & \latesttext{\textbf{1.013}\textsuperscript{*}} & \latesttext{0.939} & \latesttext{\textbf{21.2\%}\textsuperscript{*}} & \latesttext{\textbf{37.1\%}\textsuperscript{*}} & \latesttext{\textbf{-0.163}} & \latesttext{0.123} & \latesttext{\textbf{0.420}} \\
% Direct & \oldtext{$\approx0.1$}\latesttext{$0.1$} & \oldtext{\textit{source}}\latesttext{\textbf{0.952}} & \oldtext{\textit{source}}\latesttext{\textbf{1.015}\textsuperscript{*}} & \oldtext{\textit{source}}\latesttext{\textbf{0.942}} & \oldtext{\textit{source}}\latesttext{\textbf{21.2\%}\textsuperscript{*}} & \oldtext{\textit{source}}\latesttext{\textbf{37.2\%}\textsuperscript{*}} & \textbf{-0.160} & 0.121 & \textbf{0.416} \\
% Direct & \oldtext{$\approx0.954$}\latesttext{$\approx1$} & \oldtext{\textit{source}}\latesttext{\textbf{0.957}} & \oldtext{\textit{source}}\latesttext{1.019} & \oldtext{\textit{source}}\latesttext{\textbf{0.954}} & \oldtext{\textit{source}}\latesttext{\textbf{21.2\%}\textsuperscript{*}} & \oldtext{\textit{source}}\latesttext{37.2\%} & \textbf{-0.148} & 0.127 & \textbf{0.392} \\
% Direct & \oldtext{$\approx10.481$}\latesttext{$\approx10$} & \oldtext{\textit{source}}\latesttext{\textbf{0.978}} & \oldtext{\textit{source}}\latesttext{1.032} & \oldtext{\textit{source}}\latesttext{\textbf{0.991}\textsuperscript{*}} & \oldtext{\textit{source}}\latesttext{21.3\%} & \oldtext{\textit{source}}\latesttext{37.4\%} & \textbf{-0.105}\textsuperscript{*} & 0.134 & \textbf{0.318} \\
% Direct & 100 & \oldtext{\textit{source}}\latesttext{\textbf{0.992}\textsuperscript{*}} & \oldtext{\textit{source}}\latesttext{1.049} & \oldtext{\textit{source}}\latesttext{\textbf{1.012}\textsuperscript{*}} & \oldtext{\textit{source}}\latesttext{23.3\%} & \oldtext{\textit{source}}\latesttext{40.3\%} & \textbf{-0.093}\textsuperscript{*} & \textbf{0.119}\textsuperscript{*} & \textbf{0.281} \\
% \addlinespace[3pt]
% \multicolumn{10}{l}{\textbf{Surrogate penalty}} \\
% \cmidrule(lr){1-10}
% \oldtext{Surrogate} & \oldtext{$\approx0.01$} & \oldtext{\textit{pending canonical v2 extraction}} \\
% \latesttext{Surrogate} & \latesttext{$\approx0.01$} & \latesttext{0.948} & \latesttext{\textbf{1.011}\textsuperscript{*}} & \latesttext{0.938} & \latesttext{\textbf{21.2\%}\textsuperscript{*}} & \latesttext{37.2\%} & \latesttext{\textbf{-0.159}} & \latesttext{\textbf{0.119}\textsuperscript{*}} & \latesttext{\textbf{0.413}} \\
% Surrogate & \oldtext{$\approx0.1$}\latesttext{$0.1$} & \oldtext{\textit{source}}\latesttext{0.948} & \oldtext{\textit{source}}\latesttext{\textbf{1.009}\textsuperscript{*}} & \oldtext{\textit{source}}\latesttext{0.939} & \oldtext{\textit{source}}\latesttext{\textbf{21.1\%}\textsuperscript{*}} & \oldtext{\textit{source}}\latesttext{\textbf{36.8\%}\textsuperscript{*}} & \textbf{-0.153} & \textbf{0.116}\textsuperscript{*} & \textbf{0.402} \\
% Surrogate & \oldtext{$\approx0.954$}\latesttext{$\approx1$} & \oldtext{\textit{source}}\latesttext{0.950} & \oldtext{\textit{source}}\latesttext{\textbf{1.003}\textsuperscript{*}} & \oldtext{\textit{source}}\latesttext{\textbf{0.946}} & \oldtext{\textit{source}}\latesttext{\textbf{20.9\%}\textsuperscript{*}} & \oldtext{\textit{source}}\latesttext{\textbf{36.1\%}\textsuperscript{*}} & \textbf{-0.120} & \textbf{0.096}\textsuperscript{*} & \textbf{0.343} \\
% Surrogate & \oldtext{$\approx10.481$}\latesttext{$\approx10$} & \oldtext{\textit{source}}\latesttext{0.949} & \oldtext{\textit{source}}\latesttext{0.979} & \oldtext{\textit{source}}\latesttext{\textbf{0.948}} & \oldtext{\textit{source}}\latesttext{21.5\%} & \oldtext{\textit{source}}\latesttext{\textbf{35.9\%}\textsuperscript{*}} & \textbf{-0.050}\textsuperscript{*} & \textbf{0.092}\textsuperscript{*} & \textbf{0.258}\textsuperscript{*} \\
% Surrogate & 100 & \oldtext{\textit{source}}\latesttext{0.948} & \oldtext{\textit{source}}\latesttext{0.964} & \oldtext{\textit{source}}\latesttext{\textbf{0.945}} & \oldtext{\textit{source}}\latesttext{22.7\%} & \oldtext{\textit{source}}\latesttext{\textbf{36.8\%}\textsuperscript{*}} & \textbf{-0.015}\textsuperscript{*} & \textbf{0.111}\textsuperscript{*} & \textbf{0.266}\textsuperscript{*} \\
% \bottomrule
% \end{tabular}}
% \vspace{1mm}
% \begin{minipage}{\textwidth}
% \scriptsize
% \emph{Notes.}
% \latesttext{The nominal display anchors $\{0.01,0.1,1,10,100\}$ correspond to tested values $\{0.0104811,0.1,0.954095,10.4811,100\}$ on the frozen augmented grid.}
% \latesttext{Formatting comparisons are evaluated using the unrounded canonical values; entries that appear equal after display rounding can therefore receive different boldface formatting.} The three valuation-level ratios are compared by distance from one; COD and COV summarize horizontal dispersion. For the mechanism/residual-structure diagnostics, smaller $|\beta_{\log}|$, $\Delta_{\mathrm{NL}}$, and dCor indicate less first-order, non-affine conditional-mean, and broader residual--price structure, respectively; these directions are diagnostic rather than assessor-standard equity targets. Within the two baseline rows, boldface marks the baseline closer to the relevant ideal or preferred direction. For Direct and Surrogate rows, boldface marks improvement over Ordinary LightGBM, while an asterisk marks improvement over both baselines. Boldface and asterisks do not denote statistical significance, formal compliance, or model selection.
% \end{minipage}
% \end{table}
% }


\begin{table}[!htbp]
\centering
\scriptsize
\renewcommand{\arraystretch}{1.08}
\caption{\latesttext{Complementary regularization-path summary at the same nominal decade-spaced display anchors: valuation level, horizontal uniformity, and mechanism/residual structure.}}
\label{tab:path_anchor_complementary}
\resizebox{\textwidth}{!}{%
\begin{tabular}{llrrrrrrrr}
\toprule
Method & $\rho$
& Median $r$
& Mean $r$
& W. mean $r$
& COD
& COV
& $\beta_{\log}$
& $\Delta_{\mathrm{NL}}$
& dCor \\
\midrule
\multicolumn{10}{l}{\textit{Panel A: Held-out evaluation}} \\
\addlinespace[2pt]
\multicolumn{10}{l}{\textbf{Baseline models}} \\
\cmidrule(lr){1-10}
Linear regression & -- & \textbf{0.969} & 1.020 & \textbf{0.979} & 24.7\% & 45.2\% & \textbf{-0.092} & 0.131 & \textbf{0.250} \\
Ordinary LightGBM & -- & 0.929 & \textbf{0.989} & 0.924 & \textbf{21.6\%} & \textbf{39.7\%} & -0.150 & \textbf{0.119} & 0.387 \\
\addlinespace[3pt]

\multicolumn{10}{l}{\textbf{Direct penalty}} \\
\cmidrule(lr){1-10}
\latesttext{Direct} & \latesttext{$\approx0.01$} & \latesttext{0.926} & \latesttext{0.986} & \latesttext{0.922} & \latesttext{21.7\%} & \latesttext{39.9\%} & \latesttext{\textbf{-0.148}} & \latesttext{0.119} & \latesttext{\textbf{0.385}} \\
Direct & \latesttext{$0.1$} & \latesttext{0.924} & \latesttext{0.983} & \latesttext{0.920} & \latesttext{21.6\%} & \latesttext{39.8\%} & \textbf{-0.147} & 0.119 & \textbf{0.382} \\
Direct & \latesttext{$\approx1$} & \latesttext{\textbf{0.936}} & \latesttext{\textbf{0.993}\textsuperscript{*}} & \latesttext{\textbf{0.936}} & \latesttext{21.6\%} & \latesttext{39.9\%} & \textbf{-0.134} & 0.122 & \textbf{0.359} \\
Direct & \latesttext{$\approx10$} & \latesttext{\textbf{0.953}} & \latesttext{\textbf{1.005}\textsuperscript{*}} & \latesttext{\textbf{0.970}} & \latesttext{21.9\%} & \latesttext{\textbf{39.4\%}\textsuperscript{*}} & \textbf{-0.096} & 0.132 & \textbf{0.302} \\
Direct & 100 & \latesttext{\textbf{0.968}} & \latesttext{1.022} & \latesttext{\textbf{0.993}\textsuperscript{*}} & \latesttext{24.5\%} & \latesttext{41.5\%} & \textbf{-0.079}\textsuperscript{*} & \textbf{0.115}\textsuperscript{*} & \textbf{0.265} \\
\addlinespace[3pt]

\multicolumn{10}{l}{\textbf{Surrogate penalty}} \\
\cmidrule(lr){1-10}
\latesttext{Surrogate} & \latesttext{$\approx0.01$} & \latesttext{0.929} & \latesttext{0.987} & \latesttext{0.924} & \latesttext{\textbf{21.6\%}\textsuperscript{*}} & \latesttext{39.9\%} & \latesttext{\textbf{-0.147}} & \latesttext{\textbf{0.116}\textsuperscript{*}} & \latesttext{\textbf{0.383}} \\
Surrogate & \latesttext{$0.1$} & \latesttext{0.927} & \latesttext{0.982} & \latesttext{0.923} & \latesttext{\textbf{21.5\%}\textsuperscript{*}} & \latesttext{39.9\%} & \textbf{-0.138} & \textbf{0.113}\textsuperscript{*} & \textbf{0.367} \\
Surrogate & \latesttext{$\approx1$} & \latesttext{0.929} & \latesttext{0.975} & \latesttext{\textbf{0.930}} & \latesttext{\textbf{21.4\%}\textsuperscript{*}} & \latesttext{\textbf{38.2\%}\textsuperscript{*}} & \textbf{-0.102} & \textbf{0.099}\textsuperscript{*} & \textbf{0.310} \\
Surrogate & \latesttext{$\approx10$} & \latesttext{0.927} & \latesttext{0.953} & \latesttext{\textbf{0.931}} & \latesttext{22.2\%} & \latesttext{\textbf{38.5\%}\textsuperscript{*}} & \textbf{-0.036}\textsuperscript{*} & \textbf{0.103}\textsuperscript{*} & \textbf{0.250} \\
Surrogate & 100 & \latesttext{0.923} & \latesttext{0.937} & \latesttext{\textbf{0.928}} & \latesttext{23.5\%} & \latesttext{40.3\%} & \textbf{-0.001}\textsuperscript{*} & 0.124 & \textbf{0.267} \\

\midrule
\addlinespace[2pt]
\multicolumn{10}{l}{\textit{Panel B: 2025 forward evaluation}} \\
\addlinespace[2pt]

\multicolumn{10}{l}{\textbf{Baseline models}} \\
\cmidrule(lr){1-10}
Linear regression & -- & \textbf{1.020} & 1.081 & \textbf{1.027} & 24.3\% & 42.1\% & \textbf{-0.109} & 0.122 & \textbf{0.269} \\
Ordinary LightGBM & -- & 0.950 & \textbf{1.015} & 0.941 & \textbf{21.3\%} & \textbf{37.2\%} & -0.164 & \textbf{0.121} & 0.422 \\
\addlinespace[3pt]

\multicolumn{10}{l}{\textbf{Direct penalty}} \\
\cmidrule(lr){1-10}
\latesttext{Direct} & \latesttext{$\approx0.01$} & \latesttext{0.948} & \latesttext{\textbf{1.013}\textsuperscript{*}} & \latesttext{0.939} & \latesttext{\textbf{21.2\%}\textsuperscript{*}} & \latesttext{\textbf{37.1\%}\textsuperscript{*}} & \latesttext{\textbf{-0.163}} & \latesttext{0.123} & \latesttext{\textbf{0.420}} \\
Direct & \latesttext{$0.1$} & \latesttext{\textbf{0.952}} & \latesttext{\textbf{1.015}\textsuperscript{*}} & \latesttext{\textbf{0.942}} & \latesttext{\textbf{21.2\%}\textsuperscript{*}} & \latesttext{\textbf{37.2\%}\textsuperscript{*}} & \textbf{-0.160} & 0.121 & \textbf{0.416} \\
Direct & \latesttext{$\approx1$} & \latesttext{\textbf{0.957}} & \latesttext{1.019} & \latesttext{\textbf{0.954}} & \latesttext{\textbf{21.2\%}\textsuperscript{*}} & \latesttext{37.2\%} & \textbf{-0.148} & 0.127 & \textbf{0.392} \\
Direct & \latesttext{$\approx10$} & \latesttext{\textbf{0.978}} & \latesttext{1.032} & \latesttext{\textbf{0.991}\textsuperscript{*}} & \latesttext{21.3\%} & \latesttext{37.4\%} & \textbf{-0.105}\textsuperscript{*} & 0.134 & \textbf{0.318} \\
Direct & 100 & \latesttext{\textbf{0.992}\textsuperscript{*}} & \latesttext{1.049} & \latesttext{\textbf{1.012}\textsuperscript{*}} & \latesttext{23.3\%} & \latesttext{40.3\%} & \textbf{-0.093}\textsuperscript{*} & \textbf{0.119}\textsuperscript{*} & \textbf{0.281} \\
\addlinespace[3pt]

\multicolumn{10}{l}{\textbf{Surrogate penalty}} \\
\cmidrule(lr){1-10}
\latesttext{Surrogate} & \latesttext{$\approx0.01$} & \latesttext{0.948} & \latesttext{\textbf{1.011}\textsuperscript{*}} & \latesttext{0.938} & \latesttext{\textbf{21.2\%}\textsuperscript{*}} & \latesttext{37.2\%} & \latesttext{\textbf{-0.159}} & \latesttext{\textbf{0.119}\textsuperscript{*}} & \latesttext{\textbf{0.413}} \\
Surrogate & \latesttext{$0.1$} & \latesttext{0.948} & \latesttext{\textbf{1.009}\textsuperscript{*}} & \latesttext{0.939} & \latesttext{\textbf{21.1\%}\textsuperscript{*}} & \latesttext{\textbf{36.8\%}\textsuperscript{*}} & \textbf{-0.153} & \textbf{0.116}\textsuperscript{*} & \textbf{0.402} \\
Surrogate & \latesttext{$\approx1$} & \latesttext{0.950} & \latesttext{\textbf{1.003}\textsuperscript{*}} & \latesttext{\textbf{0.946}} & \latesttext{\textbf{20.9\%}\textsuperscript{*}} & \latesttext{\textbf{36.1\%}\textsuperscript{*}} & \textbf{-0.120} & \textbf{0.096}\textsuperscript{*} & \textbf{0.343} \\
Surrogate & \latesttext{$\approx10$} & \latesttext{0.949} & \latesttext{0.979} & \latesttext{\textbf{0.948}} & \latesttext{21.5\%} & \latesttext{\textbf{35.9\%}\textsuperscript{*}} & \textbf{-0.050}\textsuperscript{*} & \textbf{0.092}\textsuperscript{*} & \textbf{0.258}\textsuperscript{*} \\
Surrogate & 100 & \latesttext{0.948} & \latesttext{0.964} & \latesttext{\textbf{0.945}} & \latesttext{22.7\%} & \latesttext{\textbf{36.8\%}\textsuperscript{*}} & \textbf{-0.015}\textsuperscript{*} & \textbf{0.111}\textsuperscript{*} & \textbf{0.266}\textsuperscript{*} \\

\bottomrule
\end{tabular}}

\vspace{1mm}
\begin{minipage}{\textwidth}
\scriptsize
\emph{Notes.}
\latesttext{The nominal display anchors $\{0.01,0.1,1,10,100\}$ correspond to tested values $\{0.0104811,0.1,0.954095,10.4811,100\}$ on the frozen augmented grid.}
\latesttext{Formatting comparisons are evaluated using the unrounded canonical values; entries that appear equal after display rounding can therefore receive different boldface formatting.} The three valuation-level ratios are compared by distance from one; COD and COV summarize horizontal dispersion. For the mechanism/residual-structure diagnostics, smaller $|\beta_{\log}|$, $\Delta_{\mathrm{NL}}$, and dCor indicate less first-order, non-affine conditional-mean, and broader residual--price structure, respectively; these directions are diagnostic rather than assessor-standard equity targets. Within the two baseline rows, boldface marks the baseline closer to the relevant ideal or preferred direction. For Direct and Surrogate rows, boldface marks improvement over Ordinary LightGBM, while an asterisk marks improvement over both baselines. Boldface and asterisks do not denote statistical significance, formal compliance, or model selection.
\end{minipage}
\end{table}



\oldtext{The machine-readable combined path table now includes held-out and 2025 $\Delta_{\mathrm{NL}}$ for Linear, ordinary LightGBM, and every Direct/Surrogate grid point. The fixed cross-fitted estimator specification in Appendix~\ref{app:metrics} is archived with those values. Fold-level $\Delta_{\mathrm{NL}}$ is not added to the already completed chronological CV. Split-specific $\rho=0$ prediction-parity audits and the centered-recalibration validation record, including $b^\star$ and row-level recalibrated predictions, are stored in the same result root.}
\oldtext{The machine-readable combined path table includes held-out and 2025 $\Delta_{\mathrm{NL}}$ for Linear, ordinary LightGBM, and every Direct/Surrogate grid point. The fixed cross-fitted estimator specification in Appendix~\ref{app:metrics} is archived with those values. Fold-level $\Delta_{\mathrm{NL}}$ is not added retrospectively to the already completed chronological CV. Split-specific $\rho=0$ prediction-parity audits are stored in the same result root.}
\latesttext{The machine-readable combined path table includes held-out and 2025 $\Delta_{\mathrm{NL}}$ for Linear, ordinary LightGBM, and every Direct/Surrogate grid point. Chronological-CV $\Delta_{\mathrm{NL}}$ was later reconstructed, fold by fold, from the already-frozen validation predictions using the same fixed cross-fitted estimator; those CV values are archived separately and were not used to define the frozen prediction/COD transition span, the $\rho$ grid, or a model-selection rule. Split-specific $\rho=0$ prediction-parity audits are stored in the same result root.}


\oldtext{The supplementary transition analysis is computed read-only from the same frozen 994-tree regularization paths. Its frozen protocol, event tables, fold and leave-one-fold-out summaries, retrospective temporal-concordance calculations, span-regret diagnostics, event-sharpness diagnostics, and paper-asset manifest are archived in the transition-analysis replication artifacts. No model is refit and the original path grid and display anchors are unchanged.}
\oldtext{The v2 transition analysis is computed read-only from the final augmented 994-tree path table. The original $0.1$--$100$ fits and v1 transition artifacts are preserved unchanged; v2 appends 32 lower-tail positive values per family under the same frozen data, split, seed, predictor set, and 994-tree configuration. The 448 new chronological-CV fits were completed first, the augmented CV transition state was frozen, and only then were the same 32 new values evaluated on the held-out and 2025 samples (128 additional out-of-time fits). The five-metric transition rule and the prespecified display anchors remain unchanged. The lower-grid protocol, event tables, fold and leave-one-fold-out summaries, retrospective temporal-concordance calculations, event-sharpness diagnostics, paper-asset manifest, and hashes linking v1 and v2 are archived in the transition-analysis replication artifacts.}
\oldtext{The v2 transition analysis is computed read-only from the final augmented 994-tree path table. The original $0.1$--$100$ fits and v1 transition artifacts are preserved unchanged; v2 appends 32 lower-tail positive values per family under the same frozen data, split, seed, predictor set, and 994-tree configuration. The 448 new chronological-CV fits were completed first, the augmented CV transition state was frozen, and only then were the same 32 new values evaluated on the held-out and 2025 samples (128 additional out-of-time fits). The five-metric transition rule and the prespecified display anchors remain unchanged. The lower-grid protocol, event tables, fold and leave-one-fold-out summaries, retrospective temporal-concordance calculations, the completed v3 artifacts \texttt{transition\_oos\_span\_regret\_v3}, \texttt{transition\_lofo\_endpoint\_summary\_v3}, \texttt{vertical\_equity\_event\_locations}, \texttt{mechanism\_event\_locations}, and \texttt{event\_sharpness\_summary\_v3}, the v3 figure QA manifest, paper-asset manifest, and hashes linking v1, v2, and v3 are archived in the transition-analysis replication artifacts.}
\latesttext{The v2 transition analysis is computed read-only from the final augmented 994-tree path table. The original $0.1$--$100$ fits and v1 transition artifacts are preserved unchanged; v2 appends 32 lower-tail positive values per family under the same frozen data, split, seed, predictor set, and 994-tree configuration. The 448 new chronological-CV fits were completed first, the augmented CV transition state was frozen, and only then were the same 32 new values evaluated on the held-out and 2025 samples (128 additional out-of-time fits). The five-metric transition rule remains unchanged and is independent of the later presentation-only display-anchor convention. The lower-grid protocol, event tables, fold and leave-one-fold-out summaries, retrospective temporal-concordance calculations, the completed v3 artifacts \texttt{transition\_oos\_span\_regret\_v3}, \texttt{transition\_lofo\_endpoint\_summary\_v3}, \texttt{vertical\_equity\_event\_locations}, \texttt{mechanism\_event\_locations}, and \texttt{event\_sharpness\_summary\_v3}, the v3 figure QA manifest, paper-asset manifest, and hashes linking v1, v2, and v3 are archived in the transition-analysis replication artifacts.}

\subsection{\texorpdfstring{\textcolor{latestPurple}{Generic CV Candidate-Region Screen}}{Generic CV Candidate-Region Screen}}
\label{subsec:candidate_screen_appendix}

{\color{latestPurple}
The candidate-region analysis is computed read-only from the same frozen augmented path and uses no CCAO-specific $\rho$ constants. All event detection is performed on $x=\log_{10}\rho$ at observed positive grid points. The minimum segment size and breakpoint-clustering tolerance are functions of the grid length and median log-grid spacing rather than fixed raw-$\rho$ distances. A scale-equivariance check multiplies every positive $\rho$ by a constant while leaving all metric paths unchanged; the selected event indices are unchanged and the reported $\rho$ values rescale by the same factor. Synthetic Direct-like, Surrogate-like, and no-pathology paths verify that the screen detects the intended qualitative regimes and abstains when a guardrail is unsupported. Two independent Phase-A runs reproduce the same machine-readable endpoint hashes.

For the benefit-side activity onset, the screen uses neutral-distance versions of PRD, PRB, MKI, VEI, and $\beta_{\log}$ and requires at least three of five supported active-improvement regimes. For Direct, the upper guardrail is the earliest qualifying cluster of sustained predictive-deterioration events among $R^2_P$, MAE, MAPE, and $\operatorname{RMSE}_{\log P}$. For Surrogate, that first predictive bend is recorded but does not bind by itself. Predictive harm additionally requires exhaustion of the corresponding within-family $\rho=0$ advantage, and the final Surrogate guardrail is the earlier of a stable predictive-harm cluster and a stable post-valley $\Delta_{\mathrm{NL}}$ rebound. In the CCAO path the nonlinear caution binds before the predictive-harm cluster. Distance correlation is classified as flattening rather than rebounding at the caution event, and the ratio-profile QA shows visible deformation mostly farther along the path. Thus the Surrogate upper endpoint is an early structural warning, not a claim that visible pathology starts exactly there.

The final CCAO positive-$\rho$ candidate regions are Direct $[0.355648,\,2.559548]$ and Surrogate $[0.202359,\,2.222996]$. The lower endpoints are soft activity markers. All candidate-region endpoints are frozen from chronological CV before held-out and 2025 overlays are inspected; no out-of-time result changes an endpoint.
}

\todo{Before submission, make sure the released replication package and manuscript provenance reference the final v2.1 screening specification rather than the historical v2 Surrogate rule. The utility-level method note should be explicitly versioned so that the earlier ``predictive bend'' guardrail cannot be mistaken for the final predictive-harm/nonlinear-caution rule.}

\subsection{\texorpdfstring{\textcolor{latestPurple}{Cross-Metric Turning Points and Temporal Concordance}}{Cross-Metric Turning Points and Temporal Concordance}}
\label{subsec:transition_appendix}

\oldtext{The main text reports only the interpretation of the five-metric transition diagnostic. The objects below provide the auditable event locations and value-based checks. The Direct interval is always described as a \emph{CV-derived descriptive transition span}, never as a recommended, safe, optimal, or deployment range.}
\oldtext{The main text reports only the interpretation of the five-metric transition diagnostic. The objects below provide the auditable event locations and fold/leave-one-fold-out stability checks. Any family-specific interval is described only as a \emph{CV-derived descriptive transition span}, never as a recommended, safe, optimal, preferred, or deployment range.}
\latesttext{The main text reports only the interpretation of the five-metric transition diagnostic. The objects below provide the auditable event locations, fold/leave-one-fold-out stability checks, and family-symmetric value-based span regret. Any family-specific interval is described only as a \emph{CV-derived descriptive transition span}, never as a recommended, safe, optimal, preferred, or deployment range.}

\begin{table}[!htbp]
\centering
\scriptsize
\setlength{\tabcolsep}{3.2pt}
\renewcommand{\arraystretch}{1.08}
\caption{\oldtext{Cross-metric turning-event summary for the frozen chronological-CV diagnostic and its retrospective temporal comparison. The Direct interval is a CV-derived descriptive transition span, not a selected or recommended penalty range.}\latesttext{Cross-metric turning-event summary for the frozen augmented-grid chronological-CV diagnostic and its retrospective temporal comparison. The family-specific intervals are CV-derived descriptive transition spans, not selected or recommended penalty ranges.}}
\label{tab:transition_summary}
\begin{tabularx}{\textwidth}{@{} l >{\raggedright\arraybackslash}X >{\raggedright\arraybackslash}X @{}}
\toprule
& \textbf{Direct} & \textbf{Surrogate} \\
\midrule
CV $R^2_P$ max & $\rho=1.099$ (interior $+$)
& $\rho=0.829$ (interior $+$) \\
CV $\operatorname{MAE}_P$ min & $\rho=0.543$ (interior $+$)
& \oldtext{$\rho=0.1$ (interior $+$)}\latesttext{$\rho=0.0569$ (interior $+$)} \\
CV $\operatorname{MAPE}_P$ min & $\rho=0.471$ (interior $+$)
& $\rho=0.829$ (interior $+$) \\
CV $\operatorname{RMSE}_{\log P}$ min & \oldtext{$\rho=0.1$ (interior $+$)}\latesttext{$\rho=0.0494$ (interior $+$)}
& \oldtext{$\rho=0$ (boundary ($\rho=0$))}\latesttext{$\rho=0.00222$ (interior $+$)} \\
CV COD min & $\rho=0.543$ (interior $+$)
& $\rho=0.954$ (interior $+$) \\
\addlinespace
Common five-metric CV span
& \oldtext{valid positive-interior span $[0.1,\,1.099]$; $\log_{10}$-width $1.04$}\latesttext{valid positive-interior span $[0.0494,\,1.099]$; $\log_{10}$-width $1.35$}
& \oldtext{not supported}\latesttext{valid positive-interior span $[0.00222,\,0.954]$; $\log_{10}$-width $2.63$} \\
First-positive-grid note
& \oldtext{$\rho=0.1$ is the smallest positive tested value, not an estimated lower threshold}\latesttext{lower span endpoint $0.0494$ is interior; smallest positive tested value is $0.00110$}
& \oldtext{---}\latesttext{lower span endpoint $0.00222$ is interior; smallest positive tested value is $0.00110$} \\
LOFO five-event support
& \oldtext{7/7; valid endpoints $[0.1,\,0.153]$ to $[1.099,\,1.265]$}\oldtext{7/7}\latesttext{7/7; lower $[0.004498,\,0.065513]$, upper $[1.098541,\,1.264855]$}
& \oldtext{2/7}\oldtext{7/7}\latesttext{7/7; lower $[0.001099,\,0.010481]$, upper $[0.828643,\,1.676833]$} \\
Held-out exact concordance
& \oldtext{1/5}\latesttext{0/5}
& \oldtext{0/5}\latesttext{1/5 (MAE)} \\
2025 exact concordance
& 0/5
& \oldtext{0/5}\latesttext{1/5 (MAE)} \\
\oldtext{Value-based reading}\latesttext{Temporal reading}
& \oldtext{mixed by metric}\latesttext{valid full-CV span and 7/7 LOFO support, but 0/5 exact concordance in both later periods}
& \oldtext{no common five-metric positive CV span}\latesttext{valid full-CV span and 7/7 LOFO support, but only 1/5 exact concordance in each later period} \\
\bottomrule
\end{tabularx}
\vspace{1mm}
\begin{minipage}{\textwidth}
\scriptsize
\emph{Notes.}
Event locations are discrete observed grid points on equal-weight chronological CV.
A common span is reported only when all five events are positive and interior.
\latesttext{The five-metric rule is unchanged from v1; only the near-zero grid resolution was extended before the new lower-tail out-of-time fits were run.}
Held-out and 2025 comparisons are retrospective temporal concordance, not prospective confirmation.
No $\rho$ is selected.
\latesttext{Leave-one-fold-out 7/7 records that a positive-interior five-event span exists in every omitted-fold aggregation; it does not imply precise endpoint stability. Direct $\log_{10}$-width min/median/max is $1.286/1.347/2.449$ with $0/7$ first-positive lower-bound hits; Surrogate $\log_{10}$-width min/median/max is $2.204/2.633/2.939$ with $1/7$ first-positive lower-bound hits. Lower-end location is materially less stable than the upper end.}
\end{minipage}
\end{table}

\begin{figure}[!htbp]
\centering
\safeincludegraphics[width=0.95\textwidth]{img/generated_v12_994/paper_transition_event_locations.pdf}
\caption{\oldtext{Turning-event locations for the five fixed criteria. Gray shading, where present, is the frozen CV-derived descriptive transition span and is not a selected or recommended penalty interval.}\latesttext{Turning-event locations for the five fixed prediction/COD criteria. Gray fill and dashed vertical boundaries mark the frozen family-specific CV-derived descriptive transition span; they are not a selected or recommended penalty interval.}}
\label{fig:transition_event_locations}
\end{figure}

\begin{figure}[!htbp]
\centering
\safeincludegraphics[width=0.95\textwidth]{img/generated_v12_994/vertical_equity_event_locations.pdf}
\caption{\oldtext{Closest-to-neutral observed-grid locations for PRD, PRB, MKI, and VEI. Gray fill and dashed vertical boundaries mark the already-frozen five-primary-metric CV-derived descriptive transition span; these events do not redefine that span and do not select $\rho$.}\latesttext{Closest-to-neutral observed-grid locations for PRD, PRB, MKI, and VEI. Gray fill and dashed vertical boundaries mark the already-frozen CV-derived descriptive transition span constructed from the five prediction/COD criteria; these events do not redefine that span and do not select $\rho$.}}
\label{fig:vertical_equity_event_locations}
\end{figure}

\begin{figure}[!htbp]
\centering
\safeincludegraphics[width=0.95\textwidth]{img/generated_v12_994/mechanism_event_locations.pdf}
\caption{\oldtext{Mechanism turning-event locations corresponding to minimum $|\beta_{\log}|$ and minimum dCor. Gray fill and dashed vertical boundaries mark the frozen five-primary-metric CV-derived descriptive transition span as context only. No CV $\Delta_{\mathrm{NL}}$ event is shown because $\Delta_{\mathrm{NL}}$ was not computed retrospectively for chronological CV. These events do not define a selected or recommended range.}\oldtext{Mechanism turning-event locations corresponding to minimum $|\beta_{\log}|$ and minimum dCor. Gray fill and dashed vertical boundaries mark, for context only, the frozen CV-derived descriptive transition span constructed from the five prediction/COD criteria. No CV $\Delta_{\mathrm{NL}}$ event is shown because $\Delta_{\mathrm{NL}}$ was not computed retrospectively for chronological CV. These events do not define a selected or recommended range.}\latesttext{Mechanism turning-event locations corresponding to minimum $|\beta_{\log}|$, minimum $\Delta_{\mathrm{NL}}$, and minimum dCor. Gray fill and dashed vertical boundaries mark, for context only, the frozen CV-derived prediction/COD transition span. Chronological-CV $\Delta_{\mathrm{NL}}$ events use the same frozen estimator applied retrospectively to stored validation predictions; they do not redefine that span or select $\rho$.}}
\label{fig:mechanism_event_locations}
\end{figure}

\latesttext{A post-hoc six-metric CV equity/mechanism bend diagnostic did not yield a common leave-one-fold-out-stable span, so no second $\rho$ band is reported.}

\begin{oldrevisionblock}
\noindent\textbf{Previous Direct-only span-regret table.}\par\smallskip
\oldtext{Value-based cost of restricting the Direct out-of-time paths to the frozen CV-derived descriptive span. No small/large regret threshold is imposed.}\par\smallskip
\centering
\scriptsize
\setlength{\tabcolsep}{2.6pt}
\renewcommand{\arraystretch}{1.08}
\resizebox{\textwidth}{!}{%
\begin{tabular}{lrrrrrrr}
\toprule
\oldtext{Metric} & \oldtext{$\rho^{*}$} & \oldtext{value$^{*}$} & \oldtext{$\rho_{\mathrm{span}}$} & \oldtext{value$_{\mathrm{span}}$} & \oldtext{raw regret} & \oldtext{norm. regret} & \oldtext{$\log_{10}$ dist.} \\
\midrule
\multicolumn{8}{l}{\oldtext{\textit{Panel A: held-out evaluation}}} \\
\addlinespace[2pt]
\oldtext{$R^2_P$} & \oldtext{3.393} & \oldtext{0.901} & \oldtext{1.099} & \oldtext{0.899} & \oldtext{0.0018} & \oldtext{0.056} & \oldtext{0.490} \\
\oldtext{$\operatorname{MAE}_P$} & \oldtext{1.931} & \oldtext{\$73,924} & \oldtext{0.829} & \oldtext{\$74,409} & \oldtext{\$485} & \oldtext{0.044} & \oldtext{0.245} \\
\oldtext{$\operatorname{MAPE}_P$} & \oldtext{1.265} & \oldtext{21.0\%} & \oldtext{0.829} & \oldtext{21.1\%} & \oldtext{0.033 pp} & \oldtext{0.011} & \oldtext{0.061} \\
\oldtext{$\operatorname{RMSE}_{\log P}$} & \oldtext{0.720} & \oldtext{0.289} & \oldtext{0.720} & \oldtext{0.289} & \oldtext{0} & \oldtext{0} & \oldtext{0} \\
\oldtext{COD} & \oldtext{1.265} & \oldtext{21.52} & \oldtext{1.099} & \oldtext{21.53} & \oldtext{0.015} & \oldtext{0.005} & \oldtext{0.061} \\
\midrule
\multicolumn{8}{l}{\oldtext{\textit{Panel B: 2025 forward evaluation}}} \\
\addlinespace[2pt]
\oldtext{$R^2_P$} & \oldtext{1.931} & \oldtext{0.912} & \oldtext{1.099} & \oldtext{0.910} & \oldtext{0.0014} & \oldtext{0.037} & \oldtext{0.245} \\
\oldtext{$\operatorname{MAE}_P$} & \oldtext{2.560} & \oldtext{\$76,558} & \oldtext{1.099} & \oldtext{\$76,848} & \oldtext{\$289} & \oldtext{0.023} & \oldtext{0.367} \\
\oldtext{$\operatorname{MAPE}_P$} & \oldtext{2.223} & \oldtext{20.6\%} & \oldtext{0.309} & \oldtext{20.6\%} & \oldtext{0.052 pp} & \oldtext{0.020} & \oldtext{0.306} \\
\oldtext{$\operatorname{RMSE}_{\log P}$} & \oldtext{1.456} & \oldtext{0.278} & \oldtext{0.625} & \oldtext{0.278} & \oldtext{0.0001} & \oldtext{0.003} & \oldtext{0.122} \\
\oldtext{COD} & \oldtext{2.223} & \oldtext{21.05} & \oldtext{0.309} & \oldtext{21.11} & \oldtext{0.062} & \oldtext{0.027} & \oldtext{0.306} \\
\bottomrule
\end{tabular}}
\vspace{1mm}
\begin{minipage}{\textwidth}
\scriptsize
\oldtext{\emph{Notes.} This Direct-only table belongs to the v1 span analysis. Its values are not carried forward as final v2 evidence because the final augmented experiment gives both families valid CV spans and therefore requires the regret diagnostic to be regenerated symmetrically.}
\end{minipage}
\end{oldrevisionblock}

{\color{latestPurple}
\begin{table}[!htbp]
\centering
\scriptsize
\setlength{\tabcolsep}{2.6pt}
\renewcommand{\arraystretch}{1.08}
\caption{Value-based cost of restricting each out-of-time path to its family-specific frozen CV-derived descriptive transition span. No materiality threshold is imposed.}
\label{tab:transition_regret}
\resizebox{\textwidth}{!}{%
\begin{tabular}{llrrrrr}
\toprule
Family & Metric & $\rho^{*}$ & $\rho_{\mathrm{span}}$ & raw regret & norm.\ regret & $\log_{10}$ dist. \\
\midrule
\multicolumn{7}{l}{\textit{Panel A: Held-out evaluation}} \\
\addlinespace[2pt]
Direct & $R^2_P$ & 3.393 & 1.099 & 0.0018 & 0.056 & 0.490 \\
Direct & $\operatorname{MAE}_P$ & 1.931 & 0.829 & \$485 & 0.044 & 0.245 \\
Direct & $\operatorname{MAPE}_P$ & 1.265 & 0.829 & 0.033 pp & 0.011 & 0.061 \\
Direct & $\operatorname{RMSE}_{\log P}$ & 0.0160 & 0.0569 & 0.000461 & 0.013 & 0.490 \\
Direct & COD & 1.265 & 1.099 & 0.01451 & 0.005 & 0.061 \\
Surrogate & $R^2_P$ & 2.947 & 0.954 & 0.0004 & 0.042 & 0.490 \\
Surrogate & $\operatorname{MAE}_P$ & 0.0494 & 0.0494 & \$0 & 0 & 0 \\
Surrogate & $\operatorname{MAPE}_P$ & 1.265 & 0.954 & 0.001 pp & 0.001 & 0.122 \\
Surrogate & $\operatorname{RMSE}_{\log P}$ & 0.00168 & 0.00391 & 0 & 0.000 & 0.122 \\
Surrogate & COD & 1.265 & 0.954 & 0.01622 & 0.007 & 0.122 \\
\midrule
\multicolumn{7}{l}{\textit{Panel B: 2025 forward evaluation}} \\
\addlinespace[2pt]
Direct & $R^2_P$ & 1.931 & 1.099 & 0.0014 & 0.037 & 0.245 \\
Direct & $\operatorname{MAE}_P$ & 2.560 & 1.099 & \$289 & 0.023 & 0.367 \\
Direct & $\operatorname{MAPE}_P$ & 2.223 & 0.309 & 0.052 pp & 0.020 & 0.306 \\
Direct & $\operatorname{RMSE}_{\log P}$ & 0.0121 & 0.625 & 0.000108 & 0.004 & 0.612 \\
Direct & COD & 2.223 & 0.309 & 0.06171 & 0.027 & 0.306 \\
Surrogate & $R^2_P$ & 2.947 & 0.954 & 0.0008 & 0.083 & 0.490 \\
Surrogate & $\operatorname{MAE}_P$ & 0.115 & 0.115 & \$0 & 0 & 0 \\
Surrogate & $\operatorname{MAPE}_P$ & 1.456 & 0.954 & 0.095 pp & 0.057 & 0.184 \\
Surrogate & $\operatorname{RMSE}_{\log P}$ & 0.00193 & 0.00222 & 0.000157 & 0.003 & 0.061 \\
Surrogate & COD & 1.456 & 0.829 & 0.09823 & 0.053 & 0.184 \\
\bottomrule
\end{tabular}}
\vspace{1mm}
\begin{minipage}{\textwidth}
\scriptsize
\emph{Notes.}
Each family is restricted to its own frozen CV-derived descriptive transition span.
Normalized regret is raw regret divided by the full-path metric range when that range is nonzero.
\oldtext{Zero means the out-of-time global observed-grid optimum is available within the frozen span.}
\latesttext{An exact zero regret means that the split-specific global optimum \emph{value} is attained by at least one observed-grid point inside the frozen span; the event location $\rho^{*}$ may still lie outside the span when the optimum value is tied across grid points. Displayed regret values are rounded, and the machine-readable table retains the unrounded values.}
Exact event concordance and value-based regret answer different questions.
No interpolation, near-optimality tolerance, or materiality threshold is imposed, and no $\rho$ is selected.
\end{minipage}
\end{table}
}

\oldtext{Exact event concordance is deliberately interpreted jointly with Table~\ref{tab:transition_regret}: a turning point can move outside the CV interval even when the best value available inside the interval is close to the global discrete-grid optimum. Event-sharpness diagnostics are retained in the machine-readable replication package to document this issue without defining a post-hoc near-optimality tolerance.}
\oldtext{Exact event concordance is used here only as a strict location diagnostic. Event-sharpness diagnostics remain in the machine-readable replication package rather than being promoted to an additional selection rule.}
\latesttext{Exact event concordance is used here only as a strict location diagnostic. Machine-readable event-sharpness diagnostics document flat versus sharp events without defining a near-optimality tolerance or a new $\rho$ interval.}

\oldtext{The appendix previously included a separate PRB/MKI-versus-$R^2_P$ companion figure.}
\latesttext{That standalone companion is removed from the active presentation because PRB and MKI now appear in the four-metric main accuracy--equity figure and in the complete appendix tradeoff atlas.}

\subsection{Prediction, Valuation-Level, and Horizontal-Uniformity Paths}

\oldtext{Throughout the $\rho$-evolution figures below, gray shading denotes the family-specific CV-derived descriptive transition span when one exists. It is a path-description aid, not a selected, safe, preferred, or recommended penalty interval.}
\latesttext{Throughout the $\rho$-evolution figures below, gray fill and dashed vertical boundaries mark the frozen family-specific CV-derived descriptive transition span when one exists. It is a path-description aid, not a selected, safe, preferred, or recommended penalty interval.}

\begin{figure}[!htbp]
\centering
\safeincludegraphics[width=0.7\textwidth]{img/generated_v12_994/predictive_metric_paths_candidate_region.pdf}\\[2mm]
\safeincludegraphics[width=0.7\textwidth]{img/generated_v12_994/level_uniformity_paths_candidate_region.pdf}
\caption{\oldtext{Predictive-metric paths and valuation-level/uniformity paths along the Direct and Surrogate grids, with held-out and 2025 evaluations overlaid. Gray fill and dashed vertical boundaries mark the frozen family-specific CV-derived descriptive transition span. Valuation-level panels include the ratio $=1$ reference.}\latesttext{Predictive-metric and valuation-level/uniformity paths along the Direct and Surrogate grids, with held-out and 2025 evaluations overlaid. The light fill marks the frozen CV-derived candidate region; the dashed left boundary is activity onset and the solid right boundary is the upper guardrail. The older prediction/COD transition span is retained as subordinate context. Valuation-level panels include the ratio $=1$ reference.}}
\label{fig:other_metric_paths_placeholder}
\end{figure}\begin{figure}[!htbp]
\centering
\safeincludegraphics[width=0.8\textwidth]{img/generated_v12_994/vertical_equity_metric_paths_candidate_region.pdf}
\caption{\oldtext{Vertical-equity metric paths versus $\rho$ for Direct and Surrogate. Gray fill and dashed vertical boundaries mark the frozen family-specific CV-derived descriptive transition span. Reference lines mark PRD $=1$, PRB $=0$, MKI $=1$, and VEI $=0$.}\latesttext{Vertical-equity metric paths versus $\rho$ for Direct and Surrogate with the frozen CV-derived candidate region overlaid. The dashed left boundary is the soft activity onset and the solid right boundary is the upper guardrail; the older prediction/COD span is subordinate context. Reference lines mark PRD $=1$, PRB $=0$, MKI $=1$, and VEI $=0$.}}
\label{fig:vertical_equity_metric_paths}
\end{figure}

\subsection{\newtext{Exact Rolling-Origin Fold Structure}}

\newtext{Table~\ref{tab:fold_structure} records the exact cumulative windows and sample sizes underlying the seven-fold chronological design described in Section~\ref{subsec:path_design}. In every fold, the chronologically newest 10\% of the cumulative window is used for validation; because that rule is identical across folds, it is stated once here rather than repeated as a table column.}

{\color{impGreen}
\begin{table}[!ht]
\centering
\caption{Rolling-origin validation structure for the final CCAO development sample. In every fold, the newest 10\% of observations in the cumulative window form the validation block.}
\label{tab:fold_structure}
\begin{tabular}{clrr}
\toprule
Fold & Cumulative window & $n_{\mathrm{train}}$ & $n_{\mathrm{val}}$ \\
\midrule
1 & 2016-01-01--2017-04-01 & 46,888 & 5,209 \\
2 & 2016-01-01--2018-07-01 & 100,776 & 11,197 \\
3 & 2016-01-01--2019-10-01 & 151,187 & 16,798 \\
4 & 2016-01-01--2021-01-01 & 200,908 & 22,323 \\
5 & 2016-01-01--2022-04-01 & 252,486 & 28,053 \\
6 & 2016-01-01--2023-07-01 & 298,022 & 33,113 \\
7 & 2016-01-01--2023-11-09 & 310,147 & 34,460 \\
\bottomrule
\end{tabular}
\end{table}
}

\subsection{Chronological Fold Stability}

\oldtext{The previous appendix included a separate four-metric fold-stability figure for $R^2_P$, PRD, VEI, and $\beta_{\log}$.}
\latesttext{That summary figure is removed from the active presentation because the comprehensive fold figures immediately below now report prediction, valuation level and horizontal uniformity, vertical equity, and mechanism paths separately for all available CV metrics.}\begin{figure}[!htbp]
\centering
\safeincludegraphics[width=0.8\textwidth]{img/generated_v12_994/cv_predictive_metric_paths_candidate_region.pdf}
\caption{\oldtext{Chronological-fold predictive-metric paths (thin gray) and equal-weight CV means (thick). Gray fill and dashed vertical boundaries mark the frozen family-specific CV-derived descriptive transition span.}\latesttext{Chronological-fold predictive-metric paths (thin gray) and equal-weight CV means (thick), with the CV-derived candidate region overlaid. The dashed activity-onset boundary is soft; the solid upper guardrail marks the family-specific screening cutoff. The older prediction/COD transition span remains secondary context.}}
\label{fig:cv_predictive_metric_paths}
\end{figure}\begin{figure}[!htbp]
\centering
\safeincludegraphics[width=0.8\textwidth]{img/generated_v12_994/cv_level_uniformity_paths_candidate_region.pdf}
\caption{\oldtext{Chronological-fold valuation-level and uniformity paths (thin gray) and equal-weight CV means (thick). Gray fill and dashed vertical boundaries mark the frozen family-specific CV-derived descriptive transition span. Valuation-level panels include the ratio $=1$ reference.}\latesttext{Chronological-fold valuation-level and uniformity paths (thin gray) and equal-weight CV means (thick), with the frozen CV-derived candidate region overlaid. The activity onset is a soft left marker and the upper guardrail is the screening boundary. Valuation-level panels include the ratio $=1$ reference.}}
\label{fig:cv_level_uniformity_paths}
\end{figure}\begin{figure}[!htbp]
\centering
\safeincludegraphics[width=0.8\textwidth]{img/generated_v12_994/cv_vertical_equity_metric_paths_candidate_region.pdf}
\caption{\oldtext{Chronological-fold vertical-equity paths (thin gray) and equal-weight CV means (thick). Gray fill and dashed vertical boundaries mark the frozen family-specific CV-derived descriptive transition span. Reference lines mark PRD $=1$, PRB $=0$, MKI $=1$, and VEI $=0$.}\latesttext{Chronological-fold vertical-equity paths (thin gray) and equal-weight CV means (thick), with the frozen CV-derived candidate region overlaid. The lower boundary marks activity onset, while the upper boundary is the family-specific screening guardrail. Reference lines mark PRD $=1$, PRB $=0$, MKI $=1$, and VEI $=0$.}}
\label{fig:cv_vertical_equity_metric_paths}
\end{figure}\begin{figure}[!htbp]
\centering
\safeincludegraphics[width=0.8\textwidth]{img/generated_v12_994/cv_mechanism_metric_paths_candidate_region.pdf}
\caption{\oldtext{Chronological-fold mechanism paths for $\beta_{\log}$, $\Delta_{\mathrm{NL}}$, and distance correlation (thin gray) and equal-weight CV means (thick). The $\beta_{\log}=0$ line is a first-order neutrality reference; $\Delta_{\mathrm{NL}}$ and dCor are not given forced zero lines. Gray fill and dashed vertical boundaries mark the frozen family-specific CV-derived prediction/COD transition span. CV $\Delta_{\mathrm{NL}}$ was reconstructed from frozen validation predictions with the same fixed estimator used out of time; it does not define that span or select $\rho$.}\latesttext{Chronological-fold mechanism paths for $\beta_{\log}$, $\Delta_{\mathrm{NL}}$, and distance correlation (thin gray) and equal-weight CV means (thick), with the CV-derived candidate region overlaid. For Surrogate, the solid upper guardrail is the stable post-valley $\Delta_{\mathrm{NL}}$ caution event; dCor is classified as flattening rather than rebounding there. The older prediction/COD span is shown only as context.}}
\label{fig:cv_mechanism_metric_paths}
\end{figure}

\latesttext{Descriptively, on the frozen tested grid, the retrospectively reconstructed CV $\Delta_{\mathrm{NL}}$ paths also differ across the two objectives: the Direct equal-weight CV mean reaches its smallest observed value, $0.08685$, only at the upper tested boundary $\rho=100$, whereas the Surrogate CV mean reaches an interior minimum of $0.08677$ at $\rho\approx1.677$. These locations are not selected operating points; they further illustrate that first-order, nonlinear conditional-mean, broader-dependence, and assessor-facing vertical-equity diagnostics need not share a common transition scale.}

\begin{figure}[!htbp]
\centering
\safeincludegraphics[width=0.8\textwidth]{img/generated_v12_994/ratio_shape_cv_transition_span_only.pdf}
\caption{\oldtext{Ratio-shape profiles restricted to the prespecified display anchors that lie inside each family's frozen CV-derived descriptive transition span. Families without a valid common positive span show no penalized curves.}\oldtext{Ratio-shape profiles restricted to the prespecified display anchors that lie inside each family's frozen CV-derived descriptive transition span. The horizontal line at 1 is the principal neutrality reference; lines at 0.9 and 1.1 are aggregate appraisal-level reference guides and are not binwise acceptance criteria.}\latesttext{Ratio-shape profiles restricted to the nominal decade-spaced display anchors whose nearest tested grid points lie inside each family's frozen CV-derived descriptive transition span. The horizontal line at 1 is the principal neutrality reference; lines at 0.9 and 1.1 are aggregate appraisal-level reference guides and are not binwise acceptance criteria. The filtering is presentation-only and does not redefine the span.}}
\label{fig:ratio_shape_cv_transition_span_only}
\end{figure}\begin{figure}[!htbp]
\centering
\safeincludegraphics[width=0.95\textwidth]{img/generated_v12_994/tradeoff_equity_vs_accuracy_heldout.pdf}
\caption{\oldtext{Held-out assessor-facing diagnostics (horizontal) versus predictive metrics (vertical) along the Direct and Surrogate paths.}\latesttext{Held-out assessor-facing diagnostics (horizontal) versus predictive metrics (vertical) along the Direct and Surrogate paths. Dotted vertical lines mark metric neutrality where applicable, and shaded regions are assessor-facing guidance ranges rather than compliance claims.}}
\label{fig:tradeoff_equity_vs_accuracy_heldout}
\end{figure}\begin{figure}[!htbp]
\centering
\safeincludegraphics[width=0.95\textwidth]{img/generated_v12_994/tradeoff_equity_vs_accuracy_2025.pdf}
\caption{\oldtext{2025 assessor-facing diagnostics (horizontal) versus predictive metrics (vertical) along the Direct and Surrogate paths.}\latesttext{2025 assessor-facing diagnostics (horizontal) versus predictive metrics (vertical) along the Direct and Surrogate paths. Dotted vertical lines mark metric neutrality where applicable, and shaded regions are assessor-facing guidance ranges rather than compliance claims.}}
\label{fig:tradeoff_equity_vs_accuracy_2025}
\end{figure}\begin{figure}[!htbp]
\centering
\safeincludegraphics[width=0.95\textwidth]{img/generated_v12_994/tradeoff_mechanism_vs_accuracy_heldout.pdf}
\caption{\oldtext{Held-out mechanism/residual diagnostics (horizontal) versus predictive metrics (vertical).}\latesttext{Held-out mechanism/residual diagnostics (horizontal) versus predictive metrics (vertical). The dotted vertical line at $\beta_{\log}=0$ is a first-order neutrality reference where applicable; $\Delta_{\mathrm{NL}}$ and dCor are not given assessor-style zero targets.}}
\label{fig:tradeoff_mechanism_vs_accuracy_heldout}
\end{figure}\begin{figure}[!htbp]
\centering
\safeincludegraphics[width=0.95\textwidth]{img/generated_v12_994/tradeoff_mechanism_vs_accuracy_2025.pdf}
\caption{\oldtext{2025 mechanism/residual diagnostics (horizontal) versus predictive metrics (vertical).}\latesttext{2025 mechanism/residual diagnostics (horizontal) versus predictive metrics (vertical). The dotted vertical line at $\beta_{\log}=0$ is a first-order neutrality reference where applicable; $\Delta_{\mathrm{NL}}$ and dCor are not given assessor-style zero targets.}}
\label{fig:tradeoff_mechanism_vs_accuracy_2025}
\end{figure}

\subsection{VEI Percentile-Group Profiles}

\begin{figure}[!htbp]
\centering
\safeincludegraphics[width=0.8\textwidth]{img/generated_v12_994/vei_percentile_group_profile.pdf}
\caption{VEI percentile-group median valuation ratios for Linear and ordinary LightGBM on the held-out and 2025 samples, with deterministic 90\% bootstrap intervals.}
\label{fig:vei_group_profile_placeholder}
\end{figure}

\section{CCAO Reproducibility Requirements}
\label{app:ccao_results}

% \oldtext{A replication package for the CCAO application should contain the complete penalty grid; fold-level, held-out, and forward predictions for every valid configuration; the frozen baseline hyperparameters and configuration hash; the ordinary LightGBM and custom-objective $\rho=0$ prediction-parity audit for both out-of-time splits; the fixed cross-fitted $\Delta_{\mathrm{NL}}$ estimator specification and resulting values; and the complete primary centered-recalibration validation record and predictions.}
\oldtext{The planned CCAO replication package should contain the complete penalty grid; fold-level, held-out, and forward predictions for every valid configuration; the fixed baseline hyperparameters and configuration hash; the ordinary-LightGBM/custom-objective $\rho=0$ parity audit; the fixed $\Delta_{\mathrm{NL}}$ estimator specification and resulting values; and the complete centered-recalibration validation record and predictions.}
\newtext{The planned CCAO replication package should contain the complete penalty grid; fold-level, held-out, and forward predictions for every valid configuration; the fixed baseline hyperparameters and configuration hash; the ordinary-LightGBM/custom-objective $\rho=0$ parity audit; and the fixed $\Delta_{\mathrm{NL}}$ estimator specification and resulting values.} \oldtext{It should also archive the frozen transition-analysis protocol and status files, exact event tables, fold and leave-one-fold-out summaries, Direct span-regret and event-sharpness tables, the canonical paper-asset manifest, and hashes linking those outputs to the same 994-tree data/configuration/split identity.}\oldtext{It should also archive the original and augmented grid manifests and hashes, the v2 CV-freeze status, exact family-specific event tables, fold and leave-one-fold-out summaries, retrospective concordance outputs, event-sharpness diagnostics, any regenerated family-symmetric span-regret table, the canonical paper-asset manifest, and hashes linking all v2 outputs to the same 994-tree data/configuration/split identity.}
\latesttext{It should also archive the original and augmented grid manifests and hashes, the v2 CV-freeze status, exact family-specific event tables, fold and leave-one-fold-out summaries, retrospective concordance outputs, the completed v3 artifacts \texttt{transition\_oos\_span\_regret\_v3}, \texttt{transition\_lofo\_endpoint\_summary\_v3}, \texttt{vertical\_equity\_event\_locations}, \texttt{mechanism\_event\_locations}, and \texttt{event\_sharpness\_summary\_v3}, the reconstructed chronological-CV $\Delta_{\mathrm{NL}}$ tables, the v4 decade-anchor mapping, the post-hoc bend-diagnostic status, the v3/v4 figure QA manifests, the canonical paper-asset manifest, and hashes linking all v2/v3/v4 outputs to the same 994-tree data/configuration/split identity.}
\latesttext{The final screening package should additionally archive the frozen generic v2.1 method specification, activity-onset and upper-guardrail tables, leave-one-fold-out endpoints, the Surrogate predictive-bend/predictive-harm/nonlinear-caution event-state table, raw-path QA, scale-equivariance and synthetic-behavior tests, the no-hardcoded-$\rho$ audit, and deterministic output hashes.} Before operational use, the same analysis should report block-bootstrap intervals and ratio curves by township, property class, value range, and time. These outputs separate model-selection uncertainty, implementation effects, and local variation.
\todo{Before submission, verify which listed artifacts are actually included in the released replication package. If all are archived, replace ``should contain'' with an assertive description; otherwise delete items that are not released rather than promising them.}

\oldtext{Exploratory ATTOM Benchmark and Recalibration}
\section{\newtext{Exploratory ATTOM Benchmark}}
\label{app:attom}

% \oldtext{The exploratory ATTOM benchmark uses single-family recorder transfers that pass the benchmark price and data-quality filters. Each transfer is joined to the most recent assessor-history record available before the sale and to lagged contextual variables. The final 20\% of eligible sales form a chronological test block, with a validation block inside the remaining data.}

\newtext{The exploratory ATTOM benchmark uses Recorder \texttt{TRANSFERAMOUNT} as the sale-price target for property-use code 385 (single-family), with a minimum qualified sale price of \$50,000. Each eligible transfer is joined as of the sale date to the most recent Assessor History snapshot that closed before the transaction, to the Tax Assessor location crosswalk, and to tract-level ACS information lagged two years. Prior-sale features are constructed using only earlier transactions. Within each county, eligible modeled sales are ordered chronologically; the final 20\% form the held-out block, and the final 10\% of the remaining observations form the validation block used for all model/configuration choices.}

{\color{impGreen}
\begin{table}[!htbp]
\centering
\scriptsize
\setlength{\tabcolsep}{3.6pt}
\renewcommand{\arraystretch}{1.05}
\caption{County-specific sample and modeling configuration for the exploratory ATTOM benchmark.}
\label{tab:attom_design_appendix}
\resizebox{\textwidth}{!}{%
\begin{tabular}{lrrrrlll}
\toprule
County & Modeled sales & Train & Held-out & Held-out period & Feature set & LightGBM configuration \\
\midrule
Cook, IL & 170,118 & 136,094 & 34,024 & 2024-07-09--2025-12-29 & \texttt{status\_quo\_augmented} & \texttt{cv\_top2\_r2} \\
Allegheny, PA & 72,841 & 58,272 & 14,569 & 2024-09-13--2025-12-30 & \texttt{status\_quo\_augmented} & \texttt{cv\_top1\_r2} \\
Maricopa, AZ & 581,525 & 465,220 & 116,305 & 2023-05-09--2025-12-31 & \texttt{attom\_market\_history} & \texttt{cv\_top2\_r2} \\
King, WA & 186,017 & 148,813 & 37,204 & 2023-08-23--2025-12-31 & \texttt{attom\_market\_history} & \texttt{cv\_top2\_r2} \\
Miami-Dade, FL & 137,343 & 109,874 & 27,469 & 2023-03-17--2025-12-30 & \texttt{status\_quo\_augmented} & \texttt{cv\_top1\_r2} \\
Middlesex, MA & 83,350 & 66,680 & 16,670 & 2024-04-08--2025-12-31 & \texttt{status\_quo\_augmented} & \texttt{cv\_top1\_r2} \\
\midrule
Total & 1,231,194 & 984,953 & 246,241 & -- & -- & -- \\
\bottomrule
\end{tabular}}
\vspace{1mm}
\begin{minipage}{0.98\textwidth}
\scriptsize
\emph{Notes.} ``Modeled sales'' are the property-use-385 observations entering the county model after the benchmark's qualification, history-matching, and data-availability steps. The held-out block is the final 20\% of each county's chronologically ordered modeled sample. Feature-set and LightGBM-configuration labels are benchmark configuration identifiers and are not official assessor specifications.
\end{minipage}
\end{table}
}

% \oldtext{The benchmark compares the ordinary LightGBM fit, the custom-objective zero-penalty control, positive covariance-penalty values, and centered recalibration maps of the baseline log predictions. Recalibration parameters are estimated without using the test block. Constant level calibration is reported separately because it changes the overall ratio level but not the targeted log-residual covariance.}

\oldtext{The benchmark compares ordinary LightGBM, the custom-objective zero-penalty control, positive \texttt{LGBCovPenalty[diff]} values, and the centered one-parameter recalibration family. County-specific feature sets, LightGBM configurations, penalty strengths, and recalibration strengths are chosen using development/chronological-validation information only.}
\oldtext{The benchmark compares ordinary LightGBM, the custom-objective zero-penalty control, and positive \texttt{LGBCovPenalty[diff]} values. County-specific feature sets, LightGBM configurations, and penalty strengths are chosen using development/chronological-validation information only. The benchmark applies the Pareto--Nash maximum-volume-hyperrectangle rule over MAPE, COD, $|\mathrm{PRD}-1|$, and $|\mathrm{PRB}|$ to the validation-stage penalty choice; the held-out block is evaluated only after that choice and final refit.}
\newtext{The substantive benchmark compares ordinary LightGBM with positive \texttt{LGBCovPenalty[diff]} values; the custom-objective zero-penalty fit is retained only as an implementation control. County-specific feature sets, LightGBM configurations, and penalty strengths are chosen using development/chronological-validation information only. The benchmark applies the Pareto--Nash maximum-volume-hyperrectangle rule over MAPE, COD, $|\mathrm{PRD}-1|$, and $|\mathrm{PRB}|$ to the validation-stage penalty choice; the held-out block is evaluated only after that choice and final refit.}

The current county analysis was developed iteratively and is treated as exploratory. \oldtext{A confirmatory version must freeze the feature set, LightGBM settings, penalty grid, recalibration family, and selection rule before opening any county's test block.} \newtext{A confirmatory version must freeze the feature set, LightGBM settings, penalty grid, and selection rule before opening any county's test block.} It must report the full non-dominated frontier for prediction error, PRD, PRB, COD, 
% \oldtext{assessment level}
\newtext{valuation level}, and value-decile ratio curves. Until that protocol is rerun, the county results are evidence about possible behavior and failure modes rather than definitive external validation.

\section*{Data, Code, and Institutional Statement}

\advisorchange{The primary application uses CCAO data} and a translated version of the office's published residential modeling workflow. The multi-county benchmark uses licensed ATTOM-derived records that cannot be redistributed through the paper. Reproducibility therefore depends on archiving the custom objectives, preprocessing specifications, temporal splits, and evaluation code, together with instructions for authorized data users. The results reported here are research findings, not official CCAO assessment decisions. The paper makes no claim that CCAO has adopted the correction or placed it into production.



% FINAL-PASS AUTHOR CHECKLIST (remove from clean submission source)
\todo[inline]{Final pass before submission: (1) resolve the canonical two-split $\rho=0$ parity audit, update the appendix audit and affected path-origin artifacts, and keep the implementation check out of the main Results; (2) verify the executed PRB/VEI code uses Eq.~\eqref{eq:market_value_proxy}; (3) insert exact CCAO extract/version and build/retrieval date; (4) regenerate both accuracy--equity figures without post-hoc recalibration; (5) verify Tables~\ref{tab:attom_baselines}, \ref{tab:attom_selected_penalty}, and \ref{tab:attom_design_appendix} come from the same frozen county artifact; (6) confirm which replication artifacts are actually released; and (7) compile a clean version with all \texttt{\string\oldtext}, \texttt{\string\newtext}, and visible TODO markup removed.}

\end{document}
